\chapter{Linear algebra}
\label{chap:linalg}

In several branches of science, `algebra' means `linear algebra': the study
of vector spaces and linear maps, that is, of the category of modules over a
ring $R$ in the very special case in which $R$ is a \emph{field} (and often
one restricts attention to very special fields, such as $\mathbb{R}$ or
$\mathbb{C}$).

This will be one of the main themes in this chapter. However, I will stress
that much of what can be done over a field can in fact be done over less
special rings. In fact, I will argue that working out the general theory over
\emph{integral domains} produces invaluable tools when we are working over
\emph{fields}: the paramount example being canonical forms for matrices with
entries in a field, which will be obtained as a corollary of the
classification of finitely generated modules over a PID.

Throughout the main body of this chapter (but not necessarily in the
exercises) $R$ will denote an integral domain; most of the theory can be
extended to arbitrary commutative\footnote{The hypothesis of commutativity
is very convenient, as it allows us to identify the notion of \emph{left}-
and \emph{right}-modules, and the fact that integral domains have fields of
fractions simplifies many arguments. The reader should keep in mind that
many results in this chapter can be extended to more general rings.} rings
without major difficulty.

\section{Free modules revisited}
\label{sec:linalg:free-modules-revisited}


\subsection{R-Mod}
\label{subsec:linalg:r-mod}
For generalities on modules, see \S{}III.5. A \emph{module} over $R$ is an
abelian group $M$, endowed with an action of $R$. The action of $r \in R$
on $m \in M$ is denoted $rm$: there is a notational bias towards
\emph{left}-modules (but the distinction between left- and right-modules
will be immaterial here as $R$ is commutative). The defining axioms of a
module tell us that for all $r_1, r_2, r \in R$ and $m, m_1, m_2 \in M$,
\begin{itemize}
\item $(r_1 + r_2)m = r_1 m + r_2 m$,
\item $1m = m$ and $(r_1 r_2)m = r_1(r_2 m)$,
\item $r(m_1 + m_2) = rm_1 + rm_2$.
\end{itemize}
Modules over a ring $R$ form the category $R$-\textsf{Mod}, which we
encountered in Chapter~III. This category reflects subtle and important
properties of $R$, and we are going to attempt to uncover some of these
features. As a first approximation, we look at the `full subcategory'
(cf.\ Exercise~I.3.8) of $R$-\textsf{Mod} whose objects are \emph{free}
modules and in which morphisms are ordinary morphisms in $R$-\textsf{Mod},
that is, $R$-linear group homomorphisms.

The goal of the first few sections of this chapter is to give a very
explicit description of this subcategory: in the case of finitely generated
free modules, this can be done by means of \emph{matrices} with entries in
$R$. Later in the chapter, we will see that matrices may also be used to
describe important classes of nonfree modules.

\subsection{Linear independence and bases}
\label{subsec:linalg:linear-independence-and-bases}
The reader is invited to review the definition of free $R$-modules given in
\S{}III.6.3: $F^R(S)$ denotes an $R$-module containing a given set $S$ and
universal with respect to the existence of a set-map from $S$. We proved
(Claim~III.6.3) that the module $R^{\oplus S}$ with `one component for each
element of $S$' gives an explicit realization of $F^R(S)$.

The main point of this subsection will be that, for reasonable rings $R$,
the set $S$ can be recovered\footnote{This is not necessarily the case if
$R$ is not commutative.} `abstractly' from the free module $F^R(S)$. For
example, it will follow easily that $R^m \cong R^n$ if and only if $m = n$;
this is the first indication that the category of (finitely generated) free
modules does indeed admit a simple description.
Our main tool will be the famous concepts of \emph{linearly independent
subsets} and \emph{bases}. It is easy to be imprecise in defining these
notions. In order to avoid obvious traps, I will give the definitions for
\emph{indexed sets} (cf.\ \S{}I.2.2), that is, for functions $i : I \to M$
from a (nonempty) indexing set $I$ to a given module $M$. The reader should
think of $i$ as a selection of elements of $M$, allowing for the possibility
that the elements $m_\alpha \in M$ corresponding to $\alpha \in I$ may not
all be distinct.

Recall that for all sets $I$ there is a canonical injection $j : I \to
F^R(I)$ and any function $i : I \to M$ determines a unique $R$-module
homomorphism $\varphi : F^R(I) \to M$ making the diagram
\[\begin{tikzcd}
  F^R(I) \arrow[r, "\varphi"] & M \\
  I \arrow[u, "j"] \arrow[ur, "i"']
\end{tikzcd}\]
commute: this is precisely the universal property satisfied by $F^R(I)$.
\begin{defn}{}{linalg:1.1}
We say that the indexed set $i : I \to M$ is \emph{linearly independent} if
$\varphi$ is injective; $i$ is \emph{linearly dependent} otherwise. We say
that $i$ \emph{generates} $M$ if $\varphi$ is surjective.
\end{defn}
Put in a slightly messier, but perhaps more common, way, an indexed set $S
= \{m_\alpha\}_{\alpha \in I}$ of elements of $M$ is linearly independent
if the only vanishing linear combination
\[\sum_{\alpha \in I} r_\alpha m_\alpha = 0\]
is obtained by choosing $r_\alpha = 0$ $\forall \alpha \in I$; $S$ is
linearly dependent otherwise. The indexed set generates $M$ if every
element of $M$ may be written as $\sum_{\alpha \in I} r_\alpha m_\alpha$
for some choice of $r_\alpha$. (As a notational warning/reminder, keep in
mind that only finite sums are defined in a module; therefore, in this
context the notation $\sum$ stands for a finite sum. When writing
$\sum_{\alpha \in I} m_\alpha$ in an ordinary module, one is implicitly
assuming that $m_\alpha = 0$ for all but finitely many $\alpha \in I$.)
Using indexed sets in \cref{linalg:1.1} takes care of obvious
particular cases such as $m$ and $m$ being linearly dependent since $1
\cdot m + (-1) \cdot m = 0$: if $i : I \to M$ is not itself injective,
then $\varphi$ is surely not injective (by the commutativity of the diagram
and the injectivity of $j$), so that $i$ is linearly dependent in this
case. Because of this fact, the datum of a linearly independent $i : I \to
M$ amounts to a (special) choice of \emph{distinct} elements of $M$; the
temptation to identify the elements of $I$ with the corresponding elements
of $M$ is historically irresistible and essentially harmless. Thus, it is
common to speak about linearly dependent/independent \emph{subsets} of $M$.
I will conform to this common practice; the reader should parse the
statements carefully and correct obvious imprecisions that may arise.
A simple application of Zorn's lemma shows that every module has maximal
linearly independent subsets. In fact, it gives the following conveniently
stronger statement:

\begin{lem}{}{linalg:1.2}
Let $M$ be an $R$-module, and let $S \subseteq M$ be a linearly independent
subset. Then there exists a maximal linearly independent subset of $M$
containing $S$.
\end{lem}
\begin{proof}
Consider the family $\mathscr{S}$ of linearly independent subsets of $M$
containing $S$, ordered by inclusion. Since $S$ is linearly independent,
$\mathscr{S} \neq \emptyset$. By Zorn's lemma, it suffices to verify that
every chain in $\mathscr{S}$ has an upper bound. Indeed, the union of a
chain of linearly independent subsets containing $S$ is also linearly
independent: because any relation of linear dependence only involves
finitely many elements and these elements would all belong to one subset in
the chain.
\end{proof}
\begin{rem}{}{linalg:1.3}
This statement is in fact known to be equivalent to the axiom of choice;
therefore, the use of Zorn's lemma in one form or another cannot be
by-passed.
\end{rem}

Note that the singleton $\{2\} \subseteq \Z$ is a `maximal
linearly independent subset' of $\Z$, but it does not generate
$\Z$. In general this is an additional requirement, leading to the
definition of a basis.

\begin{defn}{}{linalg:1.4}
An indexed set $B \to M$ is a \emph{basis} if it generates $M$ and is
linearly independent.
\end{defn}
Again, one often talks of bases as `subsets' of the module $M$; since the
images of all $b \in B$ are necessarily distinct elements, this is rather
harmless. When $B$ is finite (or at any rate countable), the extra
information carried by the indexed set can be encoded by \emph{ordering}
the elements of $B$; to emphasize this, one talks about \emph{ordered
bases}.

Bases are necessarily maximal linearly independent subsets and minimal
generating subsets; this holds over every ring. What will make modules over
a \emph{field}, i.e., vector spaces, so special is that the converse will
also hold.

In any case, only very special modules admit bases:
\begin{lem}{}{linalg:1.5}
An $R$-module $M$ is free if and only if it admits a basis. In fact, $B
\subseteq M$ is a basis if and only if the natural homomorphism $R^{\oplus
B} \to M$ is an isomorphism.
\end{lem}

\begin{proof}
This is immediate from \cref{linalg:1.1}: if $B \subseteq M$ is
linearly independent and generates $M$, then the corresponding homomorphism
$R^{\oplus B} \to M$ is injective and surjective. Conversely, if $\varphi :
R^{\oplus B} \to M$ is an isomorphism, then $B$ is identified with a subset
of $M$ which generates it (because $\varphi$ is surjective) and is linearly
independent (because $\varphi$ is injective).
\end{proof}
By \cref{linalg:1.5}, the choice of a basis $B$ of a free module $M$
amounts to the choice of an isomorphism $R^{\oplus B} \cong M$; this will
be an important observation in due time (e.g., in \S{}2.2).

Once a basis $B$ has been chosen for a free module $M$, then every element
$m \in M$ can be written uniquely as a linear combination
\[m = \sum_{b \in B} r_b b\]
with $r_b \in R$. As always, remember that all but finitely many of the
coefficients $r_b$ are $0$ in any such expression.

\subsection{Vector spaces}
\label{subsec:linalg:vector-spaces}
\cref{linalg:1.5} is all that is needed to prove the fundamental
observation that modules over a field are necessarily free. Recall that
modules over a field $k$ are called \emph{$k$-vector spaces} (Example
III.5.5). Elements of a vector space are called (surprise, surprise)
\emph{vectors}, while elements of the field are called\footnote{This
terminology is also often used for free modules over any ring.}
\emph{scalars}.

By \cref{linalg:1.5}, proving that vector spaces are free modules
amounts to proving that they admit bases; \cref{linalg:1.2} reduces
the matter to the following:

\begin{lem}{}{linalg:1.6}
Let $R = k$ be a field, and let $V$ be a $k$-vector space. Let $B$ be a
maximal linearly independent subset of $V$; then $B$ is a basis of $V$.
\end{lem}

Again, this should be contrasted with the situation over rings: $\{2\}$ is
a maximal linearly independent subset of $\Z$, but it is not a
basis.
\begin{proof}
Let $v \in V$, $v \notin B$. Then $B \cup \{v\}$ is not linearly
independent, by the maximality of $B$; therefore, there exist $c_0, \dots,
c_t \in k$ and (distinct) $b_1, \dots, b_t \in B$ such that
\[c_0 v + c_1 b_1 + \dots + c_t b_t = 0,\]
with not all $c_0, \dots, c_t$ equal to $0$. Now, $c_0 \neq 0$: otherwise
we would get a linear dependence relation among elements of $B$. Since $k$
is a field, $c_0$ is a unit; but then
\[v = (-c_0^{-1} c_1) b_1 + \dots + (-c_0^{-1} c_t) b_t,\]
proving that $v$ is in the span of $B$. It follows that $B$ generates $V$,
as needed.
\end{proof}
Summarizing,

\begin{prop}{}{linalg:1.7}
Let $R = k$ be a field, and let $V$ be a $k$-vector space. Let $S$ be a
linearly independent set of vectors of $V$. Then there exists a basis $B$
of $V$ containing $S$.

In particular, $V$ is free as a $k$-module.
\end{prop}

\begin{proof}
Put \cref{linalg:1.2}, \cref{linalg:1.5}, and \cref{linalg:1.6}
together.
\end{proof}
We could also contemplate this situation from the `mirror' point of view of
generating sets:

\begin{lem}{}{linalg:1.8}
Let $R = k$ be a field, and let $V$ be a $k$-vector space. Let $B$ be a
minimal generating set for $V$; then $B$ is a basis of $V$.

Every set generating $V$ contains a basis of $V$.
\end{lem}

\begin{proof}
Exercise~1.6.
\end{proof}

\cref{linalg:1.8} also fails on more general rings (Exercise~1.5). To
reiterate, over fields (but not over general rings) a subset $B$ of a
vector space is a basis $\iff$ it is a maximal linearly independent subset
$\iff$ it is a minimal generating set.

\subsection{\texorpdfstring{Recovering $B$ from $F^R(B)$}{Recovering B from F\^{}R(B)}}
\label{subsec:linalg:recovering-b-from-f-r-b}
We are ready for the `reconstruction' of a set $B$ (up to a bijection!)
from the corresponding free module $F^R(B)$. This is the result
justifying the notion of \emph{dimension} of a vector space, or, more
generally, the \emph{rank} of a free module. Again, we prove a
somewhat stronger statement.

\begin{prop}{}{linalg:1.9}
Let $R$ be an integral domain, and let $M$ be a free $R$-module. Let
$B$ be a maximal linearly independent subset of $M$, and let $S$ be a
linearly independent subset. Then\footnote{Here, $|A|$ denotes the
\emph{cardinality} of the set $A$, a notion with which the reader is
hopefully familiar. The reader will not lose much by only considering
the case in which $B$, $S$ are finite sets; but the fact is true for
`infinite-dimensional spaces' as well, as the argument shows.}
$|S| \leq |B|$.

In particular, any two maximal linearly independent subsets of a free
module over an integral domain have the same cardinality.
\end{prop}
\begin{proof}
By taking fields of fractions, the general case over an integral
domain is easily reduced to the case of vector spaces over a field;
see Exercise~1.7. We may then assume that $R = k$ is a field and
$M = V$ is a $k$-vector space.

We have to prove that there is an injective map $j : S \hookrightarrow
B$, and this can be done by an inductive process, replacing elements
of $B$ by elements of $S$ `one-by-one'. For this, let $\leq$ be a
well-ordering on $S$, let $v \in S$, and assume we have defined $j$
for all $w \in S$ with $w < v$. Let $B'$ be the set obtained from $B$
by replacing all $j(w)$ by $w$, for $w < v$, and assume (inductively)
that $B'$ is still a maximal linearly independent subset of $V$. Then
I claim that $j(v) \in B'$ may be defined so that
\begin{itemize}
\item $j(v) \neq j(w)$ for all $w < v$;
\item the set $B''$ obtained from $B'$ by replacing $j(v)$ by $v$ is
  still a maximal linearly independent subset.
\end{itemize}
(Transfinite) induction (Claim~V.3.2) then shows that $j$ is defined
and injective on $S$, as needed.

To verify my claim, since $B'$ is a maximal linearly independent set,
$B' \cup \{v\}$ is linearly dependent (as an indexed
set\footnote{This allows for the possibility that $v \in B'$
`already'. In this case, the reader can check that the process I am
about to describe gives $j(v) = v$.}), so that there exists a linear
dependence relation
\[
  \tag{$*$} c_0v + c_1b_1 + \dots + c_tb_t = 0
\]
with not all $c_i$ equal to zero and the $b_i$ distinct in $B'$.
Necessarily $c_0 \neq 0$ (because $B'$ is linearly independent); also,
necessarily not all the $b_i$ with $c_i \neq 0$ are elements of $S$
(because $S$ is linearly independent). Without loss of generality we
may then assume that $c_1 \neq 0$ and $b_1 \in B' \setminus S$. This
guarantees that $b_1 \neq j(w)$ for all $w < v$; I set $j(v) = b_1$.

All that is left now is the verification that the set $B''$ obtained
by replacing $b_1$ by $v$ in $B'$ is a maximal linearly independent
subset. But by using $(*)$ to write
\[
  v = -c_0^{-1}c_1b_1 - \dots - c_0^{-1}c_tb_t,
\]
this is an easy consequence of the fact that $B'$ is a maximal
linearly independent subset. Further details are left to the reader.
\end{proof}

\begin{exam}{}{linalg:1.10}
An \emph{uncountable} subset of $\mathbb{C}[x]$ is necessarily
linearly dependent. Indeed, $\mathbb{C}[x]$ has a countable basis over
$\mathbb{C}$: for example, $\{1, x, x^2, x^3, \dots\}$.
\end{exam}

\begin{coro}{}{linalg:1.11}
Let $R$ be an integral domain, and let $A$, $B$ be sets. Then
\[
  F^R(A) \cong F^R(B) \iff \text{there is a bijection } A \cong B.
\]
\end{coro}
\begin{proof}
Exercise~1.8.
\end{proof}

\begin{rem}{}{linalg:1.12}
We have learned in \cref{linalg:1.2} that we can `complete' every
linearly independent subset $S$ to a maximal one. The argument used
in the proof of \cref{linalg:1.9} shows that we can in fact do this by
borrowing elements of a given maximal linearly independent subset.
\end{rem}

\begin{rem}{}{linalg:1.13}
As a particular case of \cref{linalg:1.11}, we see that if $R$ is an
integral domain, then $R^m \cong R^n$ if and only if $m = n$. This
says that integral domains satisfy the `IBN (Invariant Basis Number)
property'.

Strange as it may seem, this obvious-looking fact does \emph{not} hold
over arbitrary rings: for example, the ring of endomorphisms of an
infinite-dimensional vector space does not satisfy the IBN property.
On the other hand, integral domains are a bit of an overshoot: all
commutative rings satisfy the IBN property (Exercise~1.11).
\end{rem}

One way to think about this is that the category of finitely generated
free modules over (say) an integral domain is `classified' by
$\Z^{\geq 0}$: up to isomorphisms, there is exactly one
finitely generated free module for any nonnegative integer. The task
of describing this category then amounts to describing the
homomorphisms between objects corresponding to two given nonnegative
integers; this will be done in $\S$2.1.

As a byproduct of the result of \cref{linalg:1.9}, we can now give the
following important definition.

\begin{defn}{}{linalg:1.14}
Let $R$ be an integral domain. The \emph{rank} of a free $R$-module
$M$, denoted $\operatorname{rk}_R M$, is the cardinality of a maximal
linearly independent subset of $M$. The rank of a vector space is
called the \emph{dimension}, denoted $\dim_k V$.
\end{defn}

This definition will in fact be adopted for more general finitely
generated modules when the time comes, in $\S$5.3.

Finite-dimensional vector spaces over a fixed field form a category.
Since vector spaces are free modules (\cref{linalg:1.7}),
\cref{linalg:1.11} implies that two finite-dimensional vector spaces
are isomorphic if and only if they have the same dimension.

The subscripts $R$, $k$ are often omitted, if the context permits. But
note that, for example, viewing the complex numbers as a real vector
space, we have $\dim_{\mathbb{R}} \mathbb{C} = 2$, while
$\dim_{\mathbb{C}} \mathbb{C} = 1$. So some care is warranted.

\cref{linalg:1.9} tells us that every linearly independent subset $S$
of a free $R$-module $M$ must have cardinality lower than or equal to
$\operatorname{rk}_R M$. Similarly, every generating set must have
cardinality higher than or equal to the rank. Indeed,

\begin{prop}{}{linalg:1.15}
Let $R$ be an integral domain, and let $M$ be a free $R$-module;
assume that $M$ is generated by $S$: $M = \langle S \rangle$. Then $S$
contains a maximal linearly independent subset of $M$.
\end{prop}
\begin{proof}
By Exercise~1.7 we may assume that $R$ is a field and $M = V$ is a
vector space. Use Zorn's lemma to obtain a linearly independent
subset $B \subseteq S$ which is maximal among subsets of $S$. Arguing
as in the proof of \cref{linalg:1.6} shows that $S$ is in the span of
$B$, and it follows that $B$ generates $V$. Thus $B$ is a basis, and
hence a maximal linearly independent subset of $V$, as needed.
\end{proof}

\begin{rem}{}{linalg:1.16}
I have used again the trick of switching from an integral domain to
its field of fractions. The second part of the argument would not
work over an arbitrary integral domain, since maximal linearly
independent subsets over an integral domain are not generating sets
in general.

Another standard method to reduce questions about modules over
arbitrary commutative rings to vector spaces is to mod out by a
maximal ideal; cf.\ Exercise~1.9 and following.
\end{rem}

\subsection*{Exercises}

\begin{exer}{$\lnot$}{exer:linalg:1.1}
Prove that $\mathbb{R}$ and $\mathbb{C}$ are isomorphic \emph{as
$\Q$-vector spaces}. (In particular, $(\mathbb{R}, +)$ and
$(\mathbb{C}, +)$ are isomorphic as groups.) [II.4.4]
\end{exer}
\begin{exer}{$\lnot$}{exer:linalg:1.2}
Prove that the sets listed in Exercise~III.1.4 are all $\mathbb{R}$-vector
spaces, and compute their dimensions. [1.3]
\end{exer}
\begin{exer}{}{exer:linalg:1.3}
Prove that $\mathfrak{su}(2) \cong \mathfrak{so}_3(\mathbb{R})$ as
$\mathbb{R}$-vector spaces. (This is immediate, and not particularly
interesting, from the dimension computation of Exercise~1.2. However,
these two spaces may be viewed as the tangent spaces to
$\mathrm{SU}(2)$, resp., $\mathrm{SO}_3(\mathbb{R})$, at $I$; the
surjective homomorphism $\mathrm{SU}(2) \to \mathrm{SO}_3(\mathbb{R})$
you constructed in Exercise~II.8.9 induces a more `meaningful'
isomorphism $\mathfrak{su}(2) \to \mathfrak{so}_3(\mathbb{R})$. Can you
find this isomorphism?)
\end{exer}
\begin{exer}{}{exer:linalg:1.4}
Let $V$ be a vector space over a field $k$. A \emph{Lie bracket} on
$V$ is an operation $[\cdot,\cdot] : V \times V \to V$ such that
\begin{itemize}
\item $(\forall u,v,w \in V), (\forall a,b \in k)$,
\[
  [au+bv,w] = a[u,w]+b[v,w], \qquad [w,au+bv] = a[w,u]+b[w,v],
\]
\item $(\forall v \in V)$, $[v,v] = 0$,
\item and $(\forall u,v,w \in V)$,
  $[[u,v],w] + [[v,w],u] + [[w,u],v] = 0$.
\end{itemize}
(This axiom is called the \emph{Jacobi identity}.) A vector space
endowed with a Lie bracket is called a \emph{Lie algebra}. Define a
category of Lie algebras over a given field.

Prove the following:
\begin{itemize}
\item In a Lie algebra $V$, $[u,v] = -[v,u]$ for all $u,v \in V$.
\item If $V$ is a $k$-algebra (Definition~III.5.7), then
  $[v,w] := vw - wv$ defines a Lie bracket on $V$, so that $V$ is a
  Lie algebra in a natural way.
\item This makes $\mathfrak{gl}_n(\mathbb{R})$, $\mathfrak{gl}_n(\mathbb{C})$
  into Lie algebras. The sets listed in Exercise~III.1.4 are all Lie
  algebras, with respect to a Lie bracket induced from $\mathfrak{gl}$.
\item $\mathfrak{su}(2)$ and $\mathfrak{so}_3(\mathbb{R})$ are
  isomorphic as Lie algebras over $\mathbb{R}$.
\end{itemize}
\end{exer}
\begin{exer}{\exerused}{exer:linalg:1.5}
Let $R$ be an integral domain. Prove or disprove the following:
\begin{itemize}
\item Every linearly independent subset of a free $R$-module may be
  completed to a basis.
\item Every generating subset of a free $R$-module contains a basis.
\end{itemize}
[$\S$1.3]
\end{exer}
\begin{exer}{\exerused}{exer:linalg:1.6}
Prove \cref{linalg:1.8}. [$\S$1.3]
\end{exer}
\begin{exer}{\exerused}{exer:linalg:1.7}
Let $R$ be an integral domain, and let $M = R^{\oplus A}$ be a free
$R$-module. Let $K$ be the field of fractions of $R$, and view $M$ as
a subset of $V = K^{\oplus A}$ in the evident way. Prove that a subset
$S \subseteq M$ is linearly independent in $M$ (over $R$) if and only
if it is linearly independent in $V$ (over $K$). Conclude that the
rank of $M$ (as an $R$-module) equals the dimension of $V$ (as a
$K$-vector space). Prove that if $S$ generates $M$ over $R$, then it
generates $V$ over $K$. Is the converse true? [$\S$1.4]
\end{exer}
\begin{exer}{\exerused}{exer:linalg:1.8}
Deduce \cref{linalg:1.11} from \cref{linalg:1.9}. [$\S$1.4]
\end{exer}
\begin{exer}{\exerused}{exer:linalg:1.9}
Let $R$ be a commutative ring, and let $M$ be an $R$-module. Let
$\mathfrak{m}$ be a maximal ideal in $R$, such that $\mathfrak{m}M = 0$
(that is, $rm = 0$ for all $r \in \mathfrak{m}$, $m \in M$). Define in
a natural way a vector space structure over $R/\mathfrak{m}$ on $M$.
[$\S$1.4]
\end{exer}
\begin{exer}{$\lnot$}{exer:linalg:1.10}
Let $R$ be a commutative ring, and let $F = R^{\oplus B}$ be a free
module over $R$. Let $\mathfrak{m}$ be a maximal ideal of $R$, and let
$k = R/\mathfrak{m}$ be the quotient field. Prove that
$F/\mathfrak{m}F \cong k^{\oplus B}$ as $k$-vector spaces. [1.11]
\end{exer}
\begin{exer}{\exerused}{exer:linalg:1.11}
Prove that commutative rings satisfy the IBN property. (Use
Proposition~V.3.5 and Exercise~1.10.) [$\S$1.4]
\end{exer}
\begin{exer}{}{exer:linalg:1.12}
Let $V$ be a vector space over a field $k$, and let
$R = \operatorname{End}_{k\text{-}\mathsf{Vect}}(V)$ be its ring of
endomorphisms (cf.\ Exercise~III.5.9). (Note that $R$ is \emph{not}
commutative in general.)
\begin{itemize}
\item Prove that
  $\operatorname{End}_{k\text{-}\mathsf{Vect}}(V \oplus V) \cong R^4$
  as an $R$-module.
\item Prove that $R$ does not satisfy the IBN property if
  $V = k^{\oplus \mathbb{N}}$. (Note that $V \cong V \oplus V$ if
  $V = k^{\oplus \mathbb{N}}$.)
\end{itemize}
\end{exer}
\begin{exer}{$\lnot$}{exer:linalg:1.13}
Let $A$ be an abelian group such that $\operatorname{End}_{\mathsf{Ab}}(A)$
is a field of characteristic $0$. Prove that $A \cong \Q$.
(Hint: Prove that $A$ carries a $\Q$-vector space structure;
what must its dimension be?) [IX.2.13]
\end{exer}
\begin{exer}{$\lnot$}{exer:linalg:1.14}
Let $V$ be a finite-dimensional vector space, and let
$\varphi : V \to V$ be a homomorphism of vector spaces. Prove that
there is an integer $n$ such that $\ker\varphi^{n+1} = \ker\varphi^n$
and $\operatorname{im}\varphi^{n+1} = \operatorname{im}\varphi^n$.

Show that both claims may fail if $V$ has infinite dimension. [1.15]
\end{exer}
\begin{exer}{}{exer:linalg:1.15}
Consider the question of Exercise~1.14 for free $R$-modules $F$ of
finite rank, where $R$ is an integral domain that is not a field. Let
$\varphi : F \to F$ be an $R$-module homomorphism.

What property of $R$ immediately guarantees that
$\ker\varphi^{n+1} = \ker\varphi^n$ for $n \gg 0$?

Show that there is an $R$-module homomorphism $\varphi : F \to F$
such that $\operatorname{im}\varphi^{n+1} \subsetneq
\operatorname{im}\varphi^n$ for all $n \geq 0$.
\end{exer}
\begin{exer}{$\lnot$}{exer:linalg:1.16}
Let $M$ be a module over a ring $R$. A \emph{finite composition
series} for $M$ (if it exists) is a decreasing sequence of submodules
\[
  M = M_0 \supsetneq M_1 \supsetneq \dots \supsetneq M_m = \langle 0 \rangle
\]
in which all quotients $M_i/M_{i+1}$ are \emph{simple} $R$-modules
(cf.\ Exercise~III.5.4). The \emph{length} of a series is the number
of strict inclusions. The \emph{composition factors} are the
quotients $M_i/M_{i+1}$.

Prove a Jordan-H\"older theorem for modules: any two finite
composition series of a module have the same length and the same
(multiset of) composition factors. (Adapt the proof of
Theorem~IV.3.2.)

We say that $M$ has \emph{length} $m$ if $M$ admits a finite
composition series of length $m$. This notion is well-defined as a
consequence of the result you just proved. [1.17, 1.18, 3.20, 7.15]
\end{exer}
\begin{exer}{}{exer:linalg:1.17}
Prove that a $k$-vector space $V$ has finite length as a module over
$k$ (cf.\ Exercise~1.16) if and only if it is finite-dimensional and
that in this case its length equals its dimension.
\end{exer}
\begin{exer}{}{exer:linalg:1.18}
Let $M$ be an $R$-module of finite length $m$ (cf.\ Exercise~1.16).
\begin{itemize}
\item Prove that every submodule $N$ of $M$ has finite length
  $n \leq m$. (Adapt the proof of Proposition~IV.3.4.)
\item Prove that the `descending chain condition' (d.c.c.) for
  submodules holds in $M$. (Use induction on the length.)
\item Prove that if $R$ is an integral domain that is not a field and
  $F$ is a free $R$-module, then $F$ has finite length if and only if
  it is the $0$-module.
\end{itemize}
\end{exer}
\begin{exer}{}{exer:linalg:1.19}
Let $k$ be a field, and let $f(x) \in k[x]$ be any polynomial. Prove
that there exists a multiple of $f(x)$ in which all exponents of
nonzero monomials are prime integers. (Example: for
$f(x) = 1+x^5+x^6$,
\[
  (1+x^5+x^6)(2x^2-x^3+x^5-x^8+x^9-x^{10}+x^{11})
  = 2x^2-x^3+x^5+2x^7+2x^{11}-x^{13}+x^{17}.
\]
)

(Hint: $k[x]/(f(x))$ is a finite-dimensional $k$-vector space.)
\end{exer}
\begin{exer}{$\lnot$}{exer:linalg:1.20}
Let $A$, $B$ be sets. Prove that the free groups $F(A)$, $F(B)$
($\S$II.5) are isomorphic if and only if there is a bijection
$A \cong B$. (For the interesting direction: remember that
$F(A) \cong F(B) \implies F^{ab}(A) \cong F^{ab}(B)$, by
Exercise~II.7.12). This extends the result of Exercise~II.7.13 to
possibly infinite sets $A$, $B$. [II.5.10]
\end{exer}

\section{Homomorphisms of free modules, I}
\label{sec:linalg:homomorphisms-of-free-modules-i}


\subsection{Matrices}
\label{subsec:linalg:matrices}
% As pointed out in \S{}1, Corollary 1.11 amounts to a classification

% of free modules over an integral domain: if F is a free module, then there is a
% set A (determined up to a bijection) such that F = \sim 
%  R\oplus A . The choice of such an
% isomorphism is precisely the same thing as the choice of a basis of F (Lemma
% 1.5).
% This is simultaneously good and bad news. The good news is that if F1 , F2 are
% free, then it must be possible to `understand'6
% HomR(F1 , F2)
% entirely in terms of the corresponding sets A1 , A2 such that F1 = \sim 
%  R\oplus A1 , F2 = \sim 
%  R\oplus A2 .
% That is, this set of morphisms in R-Mod may be identified with
% HomR (R\oplus A1 , R\oplus A2 ).
% The bad news is that this identification HomR(F1, F2 ) = \sim 
%  HomR (R\oplus A1 , R\oplus A2 )
% is not `canonical', because it depends on the chosen isomorphisms F1 \sim 
%  = R\oplus A1 ,
% 6 For notational simplicity I will denote Hom
% R-Mod (M, N ) by HomR (M, N ); this is very com-
% mon, and no confusion is likely.
% F2 \sim 
%  = R\oplus A2 , that is, on the choice of bases. We will therefore have to do
%  some work
% to deal with this ambiguity (in \S{}2.2).
% The point of this subsection is to deal with the good news, that is, describe
% HomR (R\oplus A1 , R\oplus A2 ). This can be done in a particularly convenient
% way when the
% free modules are finitely generated, the case we are going to analyze more
% carefully.
% Also recall (from \S{}III.5.2) that one of the good features of the category
% R-Mod
% is that the set of morphisms HomR (M, N ) between two R-modules is itself an
% R-module, in a natural way. The task is then to describe as explicitly as
% possible7
% HomR (Rn , Rm )
% as an R-module, for every choice of m, n \in  Z\geq 0.
% This will be done by means of matrices with entries in R. I trust that the
% reader is familiar with the general notion of an m \times  n matrix; I have
% occasionally
% used matrices in examples given in previous chapters, and they have showed up
% in several exercises. An m \times  n matrix with entries in R is simply a choice
% of mn
% elements of R. It is common to arrange these elements as an array consisting of
% m
% rows and n columns:
% (
%  )
% r11 r12 \dots  r1n
% | |
%  r21 r22 \dots  r2n |
%  |
% (rij )i=1,\cdot \cdot \cdot  j=1,\cdot \cdot \cdot ,m ,n
%  = (
%  | .
%  .
%  .
%  ..
%  .
%  ..
%  .
%  .. . |
%  )
% rm1 rm2 \dots  rmn
% with rij \in  R. For any given m, n the set Mm,n (R) of m \times  n matrices
% with entries
% in R is an abelian group under entrywise addition8
% (aij ) + (bij ) := (aij + b ij )
% (cf. Example II.1.5); this is in fact an R-module, under the action
% r(aij ) := (raij )
% for r \in  R. From this point of view, the set of m \times  n matrices is simply
% a copy of
% the R-module Rmn .
% However, there are other interesting operations on these sets. If A = (aik ) is
% an m \times  p matrix and B = (bkj ) is a p \times  n matrix, then one may
% define the product
% of A and B as
% ?
%  p
% A \cdot  B = (aik ) \cdot  (bkj ) := (
%  aik bkj );
% k=1
% this operation is (clearly) distributive with respect to addition and compatible
% with
% the R-module structure. It is just a tiny bit messier to check that it is
% associative,
% in the sense that if A, B, C are matrices of size, respectively, m \times  p, p
% \times  q, and
% q \times  n, then
% (A \cdot  B) \cdot  C = A \cdot  (B \cdot  C).
% 7 I am now writing Rn for what was denoted R\oplus n in previous chapters. This
% is a slight abuse
% of language, but it makes for easier typesetting.
% 8Note that I will be dropping the extra subscripts giving the range of the
% indices and will
% write (rij ) rather than (rij )i=1,\cdot \cdot \cdot  ,m : these subscripts are
% an eyesore, and the context usually
% j=1,\cdot \cdot \cdot  ,n
% makes the information redundant.
% The reader should have no trouble reconstructing the proof of this fact (Exer-
% cise 2.2).
% In particular, we have a binary operation on the abelian group Mn (R) of square
% n \times  n-matrices, and this operation is associative, distributive w.r.t. +,
% and admits
% the identity element
%  (
%  )
% 1 0 \cdot \cdot \cdot  0
% |0 |
%  1 \dots  0|
%  |
% | ( .. . .. .
%  . .
%  . .. . |
%  )
% 0 0 \cdot \cdot \cdot  1
% (the identity matrix, denoted In). That is, Mn (R) is a ring (Example III.1.6)
% and
% in fact an R-algebra (Exercise 2.3). Except in very special cases (such as n =
% 1)
% this ring is not commutative.
% Matrices of type n \times  1 are called column n-vectors; matrices of type 1
% \times  m
% are called row m-vectors, for evident reasons. An element of a free R-module Rn
% is nothing but the choice of n elements of R, and we can arrange these elements
% into row or column vectors if we please; the standard choice is to represent
% them
% as column vectors:
%  ( )
% v1
% |v2 |
% | |
% v = | . | \in  Rn .
% ( .. )
% vn
% I will denote by ei the elements of the `standard basis' of Rn :
% ( )
%  ( )
%  ( )
% |0|
%  |1|
%  |0|
% | |
%  | |
%  | |
% e1 = | . | , e2 = | . | , . . . , en = | . | ,
% ( .. )
%  ( .. )
%  ( .. )
% so that
%  (
%  )
% v1
%  n
% ?
% v = | ( ... | ) =
%  vj ej .
% vn
%  j=1
% The elements vj \in  R are the `components' of v.
% Interpreting elements of Rn as column vectors, we can act on Rn with an m \times
% n
% matrix, by left-multiplication: if A = (aij ) is an m \times  n matrix and v \in
% Rn is a
% column vector, the product is a column vector in Rm :
% (
%  ) ( ) (
%  )
% a11
%  a12
%  \cdot \cdot \cdot 
%  a1n
%  v1
%  a11v1 + a12 v2 + \dots  + a1n vn
% | |
%  a21
%  a22
%  \cdot \cdot \cdot 
%  a2n |
%  | | |v2 | | |
%  | a21v1 + a22 v2 + \dots  + a2n vn |
%  |
%  m
% A \cdot  v = | ( ..
%  .
%  .
%  ..
%  ..
%  .
%  .. .
%  | )
%  \cdot  (
%  | .. .
%  | )
%  = |
%  (
%  .
%  ..
%  | )
%  \in  R .
% am1
%  am2
%  \cdot \cdot \cdot 
%  amn
%  vn
%  am1 v1 + am2 v2 + \dots  + amn vn
% Lemma 2.1. For all m \times  n matrices A with entries in R:
% \bullet  The function \phi  : Rn \to  Rm defined by \phi (v) = A \cdot  v is a homomorphism of
% R-modules.
% \bullet  Every R-module homomorphism Rn \to  Rm is determined in this way by a
% unique m \times  n matrix.
% Proof. The first point follows immediately from the elementary properties of ma-
% trix multiplication recalled above: \forall r, s \in  R, \forall v, w \in  Rn
% \phi (rv + sw) = A \cdot  (rv + sw) = rA \cdot  v + sA \cdot  w = r\phi (v) + s\phi (w)
% as needed.
% For the second point, let \phi  : Rn \to  Rm be a homomorphism of R-modules; let
% aij be the i-th component of \phi (ej ), so that
% (
%  )
% a1j
% \phi (ej ) = |
%  ( ... |
%  ) .
% amj
% Then A = (aij ) is an m \times  n matrix, and \forall v \in  Rn with components
% vj ,
% (
%  )
%  (
%  )
% a11 v1 + a12 v2 + \dots  + a1nvn
%  a1j v j
% | |
%  a21 v1 + a22 v2 + \dots  + a2nvn |
%  | ?
%  n | | a2j v j | |
%  ?
%  n
% A \cdot  v = |
%  (
%  .
%  ..
%  | )
%  =
%  j=1
%  | ( . .. | )
%  =
%  j=1
%  vj \phi (ej )
% am1 v1 + am2 v2 + \dots  + amn vn
%  amj vj
% ?
%  n
% = \phi (
%  vj ej ) = \phi (v),
% j=1
% as needed. The homomorphism induced by a nonzero matrix is manifestly nontriv-
% ial; this implies that the matrix A associated to a homomorphism \phi  is
% uniquely
% determined by \phi .
%  ?
% The reader should attempt to remember the recipe associating a matrix A to
% a homomorphism \phi  as in the proof of Lemma 2.1: the j-th column of A is
% simply
% the column vector \phi (ej ). The collection of these vectors determines \phi
% ---since a
% homomorphism is determined by its action on a set of generators.
% Lemma 2.1 yields the promised explicit description of HomR (Rn, Rm):
% Corollary 2.2. The correspondence introduced in Lemma 2.1 gives an isomorphism
% of R-modules
% Mm,n (R) \sim 
%  = HomR (Rn , Rm ).
% Proof. The reader will check that the correspondence is a bijective homomorphism
% of R-modules; this is enough, by Exercise III.5.12.
%  ?
% This is very good news, and it gets even better. Don't forget that R-Mod is a
% category; that is, we can compose morphisms. Therefore, there is a function9
% HomR (R p , Rm ) \times  HomR (Rn , Rp ) \to  HomR (Rn , Rm ),
% 9 Here I am reversing the order of the Hom sets on the left w.r.t. the
% convention used in
% Definition I.3.1.
% mapping (\phi , \psi ) to \phi  \circ  \psi ; on the other hand, the matrix
% product gives a function
% Mm,p (R) \times  Mp,n (R) \to  Mm,n (R),
% mapping (A, B) to A \cdot  B. That is, we have a diagram
% Mm,p (R) \times  Mp,n (R)
%  / Mm,n (R)
% \sim 
%  \sim 
% \sqcup 
%  \sqcup 
% Hom R (Rp, Rm ) \times  HomR (Rn , Rp )
%  / HomR (Rn , Rm )
% where the vertical maps are (induced by) the isomorphisms obtained in Corol-
% lary 2.2.
% Lemma 2.3. This diagram commutes. That is, the matrix corresponding to a
% composition \phi  \circ  \psi  is the product of the matrices corresponding to
% \phi  and \psi .
% Proof. This follows immediately from the associativity of matrix multiplication:
% for v \in  Rn and A \in  Mm,p (R), B \in  Mp,n (R),
% A \cdot  (B \cdot  v) = (A \cdot  B) \cdot  v;
% that is, the successive action of B and A is equivalent to the action of A \cdot
% B.
%  ?
% Here is a summary of the situation. Given an integral domain R, we can
% construct a `toy category' as follows: we let Z\geq 0 be the set of objects, and
% we
% define the set of morphisms Hom(n, m) to be the set of matrices Mm,n (R), with
% composition given by the matrix product (cf. Exercise I.3.6). Then we have
% verified
% that this toy category is `essentially the same as' the category of finitely
% generated
% free R-modules and R-module homomorphisms10 .
% For example, by virtue of Proposition 1.7, if R = k is a field, then R-Mod =
% k-Vect consists exclusively of free k-modules. The toy category I just described
% is
% a faithful snapshot of the category of finite-dimensional k-vector spaces.

\subsection{Change of basis}
\label{subsec:linalg:change-of-basis}
% It is time to deal with the bad news.

% Free modules are only defined up to isomorphism, as is any structure satisfying
% a universal property. Writing a free module as a direct sum R\oplus A amounts to
% making the choice of one specific realization. As I pointed out at the end of
% \S{}1.2,
% by Lemma 1.5 this choice is equivalent to the choice of a basis of the module.
% The price to pay for representing homomorphisms of free modules by matrices
% is that this correspondence relies on the choice of a basis: we say that it is
% not
% canonical. It is essential to be able to keep track of this choice. For finitely
% generated
% free modules, this also boils down to the action of a matrix, as we proceed to
% see.
% Let F be a free module, and choose two bases A, B for F ; assume that F is
% finitely generated, so that A, B are finite sets, and further |A| = |B| (= rk F
% ) by
% 10 The reader should not take this statement too seriously: we have identified
% together all
% free modules isomorphic to a given one, which is a rather drastic operation. We
% will take a more
% formal look at this situation in Example VIII.1.8.
% Proposition 1.9. The two bases correspond to two isomorphisms
% R\oplus A
%  \phi 
%  / F,
%  R \oplus B
%  \psi 
%  / F .
% Then
% R\oplus A
%  \psi -1
%  \circ \phi 
%  / R\oplus B
% is an isomorphism, which corresponds to a matrix as seen in Lemma 2.1; we may
% call this matrix NA B . Explicitly, the j-th column of NA B is the image of the
% j-th
% element of A, viewed as a column vector in R\oplus B . That is, letting11
% A = (a1 , . . . , ar ),
%  B = (b1, . . . , br ),
% we have NA B = (nij )i=1,\cdot \cdot \cdot  ,r , with
% j=1,\cdot \cdot \cdot  ,r
% ?
%  r
% \psi  -1 \circ  \phi (aj ) =
%  nij bi .
% i=1
% This equality is written in R\oplus B ; in practice one views the basis vectors
% aj , bi as
% vectors of F and would therefore simply write
% ?
%  r
% aj =
%  nij bi ,
% i=1
% that is, the corresponding equality in F .
% Definition 2.4. The matrix NA B is called the matrix of the change of basis.
%  ?
% Since it represents an isomorphism, the matrix of the change of basis is neces-
% sarily an invertible matrix.
% To my knowledge there is no established convention on the notation for this ma-
% trix, and in any case I would find it futile to try to remember any such
% convention.
% Any choice of notation will lead to a pleasant `calculus': for example, the
% definition
% given above implies immediately NB A = (NA B )-1 and NB C NA B = NA C . In my
% experi-
% ence, any such manipulation is immediately clarified by writing isomorphisms out
% explicitly, so no convention is necessary.
% For example, let's work out the action of a change of basis on the matrix
% representation of a homomorphism \alpha  : F \to  G of two free modules. The
% diagram
% taking care of the needed manipulations is
% R\oplus A R\oplus C
% B
%  J
%  J J J \phi 
%  J
%  J J
%  %
%  \alpha 
%  y
% t
%  tt
%  \rho 
%  tt
%  t
%  t
%  D
% \nu A
%  \sqcup  t t t
%  \psi 
%  t t
%  t
%  9 F
%  / G \sigma 
%  JJ
%  J
%  \sqcup 
%  \mu  C
% eJ
%  JJ
% R\oplus B
%  R\oplus D
% Let NA B be the matrix of \nu A
%  B
%  = \psi -1 \circ  \phi , as above, and let MC D be the matrix
% for \mu C D
%  = \sigma -1
%  \circ  \rho .
% 11 Ordering the elements of a basis is convenient, for example because it allows
% us to talk
% about the `j-th element' of the basis. Hence we have the `tuple' notation.
% Choose the basis A for F and C for G. Then the matrix representing \alpha  will
% be
% the matrix P corresponding to
% \rho -1 \circ  \alpha  \circ  \phi  : R \oplus A \to  R\oplus C .
% If on the other hand we choose the basis B for F and D for G, then we represent
% \alpha  by the matrix Q corresponding to
% \sigma  -1 \circ  \alpha  \circ  \psi  : R\oplus B \to  R\oplus D .
% Now,
% \sigma  -1 \circ  \alpha  \circ  \psi  = (\mu D
%  C \circ  \rho 
% -1
%  ) \circ  \alpha  \circ  (\phi  \circ  (\nu A
%  B )-1
% ) = \mu D
%  C \circ  (\rho 
% -1
%  \circ  \alpha  \circ  \phi ) \circ  (\nu A
%  B )-1
%  ,
% and by Lemma 2.3 this shows
% Q = MC D \cdot  P \cdot  (NA B )-1 = MC D \cdot  P \cdot  NB A .
% This is hardly surprising, of course: starting from a vector expressed as a
% combina-
% tion of b's, NB A converts it into a combination of a's; P acts on it by giving
% its image
% under \alpha  as a combination of c's; and MC D converts it into d's. That
% accomplishes
% Q's job, as it should.
% Summarizing this discussion,
% Proposition 2.5. Let \alpha  : F \to  G be a homomorphism of finitely generated
% free
% modules, and let P be a matrix representing it with respect to any choice of
% bases
% for F and G. Then the matrices representing \alpha  with respect to any other
% choice of
% bases are all and only the matrices of the form
% M \cdot  P \cdot  N,
% where M and N are invertible matrices.
% Definition 2.6. Two matrices P, Q \in  Mm,n (R) are equivalent if they represent
% the same homomorphism of free modules Rn \to  Rm up to a choice of basis.
%  ?
% This is manifestly an equivalence relation. Properly speaking, the `abstract'
% homomorphism \alpha  : F \to  G is not represented by one matrix as much as by
% the whole
% equivalence class with respect to this relation. Proposition 2.5 gives a
% computational
% interpretation of equivalence of matrices: P and Q are equivalent if and only if
% there
% are invertible M and N such that Q = M P N .

\subsection{Elementary operations and Gaussian elimination}
\label{subsec:linalg:elementary-operations-and-gaussian-elimination}
% The idea is now to

% capitalize on Proposition 2.5, as follows: given a homomorphism \alpha  : F \to
% G between
% two free modules, find `special' bases in F and G so that the matrix of \alpha
% takes a
% particularly convenient form. That is, look for a particularly convenient matrix
% in
% each equivalence class with respect to the relation introduced in Definition
% 2.6.
% For this, we can start with random bases in F and G, representing \alpha  by a
% matrix P and then (by Proposition 2.5) multiply P on the right and left by
% invertible
% matrices, in order to bring P into whatever form is best suited for our needs.
% There is an even more concrete way to deal with equivalence computationally.
% Consider the following three `elementary (row/column) operations' that can be
% performed on a matrix P :
% \bullet  switch two rows (or two columns) of P ;
% \bullet  add to one row (resp., column) a multiple of another row (resp., column);
% \bullet  multiply all entries in one row (or column) of P by a unit of R.
% Proposition 2.7. Two matrices P, Q \in  Mm,n (R) are equivalent if Q may be
% obtained from P by a sequence of elementary operations.
% Proof. To see that elementary operations produce equivalent matrices, it
% suffices
% (by Proposition 2.5) to express them as multiplications on the left or right12
% by
% invertible matrices. Indeed, these operations may be performed by suitably
% multi-
% plying by the matrices obtained from the identity matrix by performing the same
% operation. For example, multiplying on the left by
% (
%  )
% 1 0 0 0
% |0 0 0 1|
% |
%  |
% (0 0 1 0)
% 0 1 0 0
% interchanges the second and fourth row of a 4 \times  n matrix; multiplying on
% the right
% by
% (
%  )
% 1 0 c
% (0 1 0)
% 0 0 1
% adds to the third column of a m \times  3 matrix the c-multiple of the first
% column. The
% diligent reader will formalize this discussion and prove that all these matrices
% are
% indeed invertible (Exercise 2.5).
%  ?
% The matrices corresponding to the elementary operations are called elementary
% matrices. Linear algebra over arbitrary rings would be a great deal simpler if
% Proposition 2.7 were an `if and only if' statement. This boils down to the
% question
% of whether every invertible matrix may be written as a product of elementary
% matrices, that is, whether elementary matrices generate the `general linear
% group'.
% Definition 2.8. The n-th general linear group over the ring R, denoted GLn (R),
% is the group of units in Mn (R), that is, the group of invertible n \times  n
% matrices with
% entries in R.
%  ?
% Brief mentions of this notion have occurred already; cf. Example II.1.5 and
% several exercises in previous chapters.
% The elementary matrices are elements of GLn (R); in fact, the inverse of an
% elementary matrix is (of course) again an elementary matrix. The following
% obser-
% vation is surely known to the reader, in one form or another; in view of the
% foregoing
% considerations, it says that the relation introduced in Definition 2.6 is under
% good
% control over fields.
% 12 Multiplying on the left acts on the rows of the matrix; multiplying on the
% right acts on its
% columns.
% Proposition 2.9. Let R = k be a field, and let n \geq  0 be an integer. Then GLn
% (k)
% is generated by elementary matrices.
% Thus, two matrices are equivalent over a field if and only if they are linked by
% a sequence of elementary operations.
% Proof. Let A = (aij ) be an n \times  n invertible matrix. In particular, some
% entry in
% the first column of A is nonzero; by performing a row switch if necessary, we
% may
% assume that a11 is nonzero. Multiplying the first row by a-1
%  11 , we may assume that
% a11 = 1:
%  (
%  )
%  a12 . . . a1n
% | a21 a22 . . . a1n |
% |
%  |
% | ( ..
%  .
%  ..
%  .
%  ..
%  .
%  .. . | )
%  .
% an1 an2 . . . ann
% Adding to the second row the (-a21 )-multiple of the first row clears the (2, 1)
% entry.
% After performing the analogous operation on all rows, we may assume that the
% only
% nonzero entry in the first column is the (1, 1) entry:
% (
%  )
% 1 a12 . . . a1n
% |0 a\setminus  22 . . . a\setminus  1n |
% |
%  |
% | ( ..
%  .
%  ..
%  .
%  ..
%  .
%  .. . | )
%  .
% 0 a\setminus  n2
%  ...
%  a\setminus  nn
% Similarly, adding to the second column the (-a12 )-multiple of the first column
% clears the (1, 2) entry. Performing this operation on all columns reduces the
% matrix
% to the form
%  (
%  )
% 1 0 ...
%  ?
%  ?
% |0 a\setminus  22 . . . a\setminus  1n |
%  1 0
% |
%  |
% | ( ..
%  .
%  ..
%  .
%  ..
%  .
%  .. . | )
%  =
%  0 A\setminus 
%  ,
% 0 a\setminus  n2 . . . a\setminus  nn
% where A\setminus  denotes a (clearly invertible) (n - 1) \times  (n - 1)
% matrix13 .
% Repeating the process on A\setminus  and on subsequent smaller matrices reduces
% A to
% the identity matrix In . In other words, In may be obtained from A by a sequence
% of elementary operations:
% In = M \cdot  A \cdot  N,
% where M and N are products of elementary matrices. But then
% A = M -1 \cdot  N -1
% is itself a product of elementary matrices, yielding the statement.
%  ?
% Note that, with notation as in the preceding proof, A-1 = N \cdot  M ; thus, the
% process explained in the proof may be used to compute the inverse of a matrix
% (also
% cf. Exercise 3.5).
% When applied only to the rows of a matrix, the simplification of a matrix by
% means of elementary operations is called Gaussian elimination. This corresponds
% 13 The vertical and horizontal lines alert the reader to the fact that the
% sectors of the matrix
% are themselves matrices; this notation is called a block matrix.
% to multiplying the given matrix on the left by a product of elementary matrices,
% and it suffices in order to reduce any square invertible matrix to the identity
% (Ex-
% ercise 2.15).
% I will sloppily call `Gaussian elimination' the more drastic process including
% column operations as well as row operations; this is in line with our focus on
% equivalence rather than `row-equivalence'. Applied to any rectangular matrix,
% this
% process yields the following:
% Proposition 2.10. Over a field, every m \times  n matrix is equivalent to a
% matrix of
% the form
%  ?
%  ?
% Ir 0
% (where r \leq  min(m, n) and `0' stands for null matrices of appropriate sizes).
% Different matrices of the type displayed in Proposition 2.10 are inequivalent
% (for example by rank considerations; cf. \S{}3.3). Thus, Proposition 2.10
% describes all
% equivalence classes of matrices over a field and shows that for any given m, n
% there
% are in fact only finitely many such classes (over a field!).

\subsection{Gaussian elimination over Euclidean domains}
\label{subsec:linalg:gaussian-elimination-over-euclidean-domains}
% With due care, Gauss-

% ian elimination may be performed over every Euclidean domain. A 2 \times  2
% example
% should suffice to illustrate the general case: let
% a b
% \in  M2 (R),
% c d
% for a Euclidean domain R, with Euclidean valuation N . After switching rows
% and/or
% columns if necessary, we may assume that N (a) is the minimum of the valuations
% of all entries in the matrices. Division with remainder gives
% b = aq + r
% with r = 0 or N (r) < N (a). Adding to the second column the (-q)-multiple of
% the
% first produces the matrix
% a
%  r
% .
% c d - qc
% If r = 0, so that N (r) < N (a), we begin again and shuffle rows and columns so
% that
% the (1, 1) entry has minimum valuation. This process may be repeated, but after
% a
% finite number of steps the (1, 2) entry will have to vanish: because valuations
% are
% nonnegative integers and at each iteration the valuation of the (1, 1) entry
% decreases.
% Trivial variations of the same procedure will clear the (2, 1) entry as well,
% producing a matrix
% e 0
% .
% 0 f
% Now (this is the cleverest part) I claim that we may assume that e divides f in
% R,
% with no remainder. Indeed, otherwise we can add the second row to the first,
% e f
% ,
% 0 f
% and start all over with this new matrix. Again, the effect of all the operations
% will
% be to decrease the valuation of the (1, 1) entry, so after a final number of
% steps we
% must reach the condition e | f .
% The reader will have some fun describing this process in the general case. The
% end result is the following useful remark:
% Proposition 2.11. Let R be a Euclidean domain, and let P \in  Mm,n(R). Then P
% is equivalent to a matrix of the form14
% (
%  )
% d1 \dots  0 0
% | . .
%  |
% | |
%  ..
%  . . ... ... |
%  |
% |
%  |
% ( 0 \dots  dr 0 )
%  \cdot \cdot \cdot 
% with d1 | \dots  | dr .
% This is called the Smith normal form of the matrix.
% Proposition 2.11 suffices to prove a weak form of one of our main goals (a
% classi-
% fication theorem for finitely generated modules over PIDs) over Euclidean
% domains;
% this will hopefully be clear in a little while, but the reader can gain the
% needed in-
% sight right away by contemplating Proposition 2.11 vis-\`{a}-vis the
% classification of
% finite abelian groups (Exercise 2.19).
% Remark 2.12. As a consequence of the preceding considerations and arguing as
% in the proof of Proposition 2.9, we see that GLn (R) is generated by elementary
% matrices if R is a Euclidean domain. The reader may think that some cleverness
% in
% handling Gaussian elimination may extend this to more general rings, but this
% will
% not go too far: there are examples of PIDs R for which GLn(R) is not generated
% by elementary matrices.
% On the other hand, some cleverness does manage to produce a Smith normal
% form for any matrix with entries in a PID: the presence of a good gcd suffices
% to
% adapt the procedure sketched above (but one may need more than elementary row
% and column operations). We will take a direct approach to this question in
% \S{}5. ?

\subsection*{Exercises}

% 2.1. Prove that the subset of M2 (R) consisting of matrices of the form
% 1 0
% r 1
% is a group under matrix multiplication and is isomorphic to (R, +).
% 2.2. ? Prove that matrix multiplication is associative. [\S{}2.1]
% 14 Here r \leq  min(m, n), the bottom-right 0 stands for a null (m - r) \times
% (n - r) matrix, etc.
% 2.3. ? Prove that both Mn (R) and HomR (Rn , Rn ) are R-algebras in a natural
% way
% and the bijection HomR (Rn , Rn ) \sim 
%  = Mn (R) of Corollary 2.2 is an isomorphism of
% R-algebras. (Cf. Exercise III.5.9.) In particular, if the matrix M corresponds
% to
% the homomorphism \phi  : Rn \to  Rn , then M is invertible in Mn (R) if and only
% if \phi
% is an isomorphism. (Note that R is commutative by default.) [\S{}2.1, \S{}3.2,
% \S{}6.1]
% 2.4. Prove Corollary 2.2.
% 2.5. ? Give a formal argument proving Proposition 2.7. [\S{}2.3]
% 2.6. \lnot  A matrix with entries in a field is in row echelon form if
% \bullet  its nonzero rows are all above the zero rows and
% \bullet  the leftmost nonzero entry of each row is 1, and it is strictly to the right of the
% leftmost nonzero entry of the row above it.
% The matrix is further in reduced row echelon form if
% \bullet  the leftmost nonzero entry of each row is the only nonzero entry in its column.
% The leftmost nonzero entries in a matrix in row echelon form are called pivots.
% Prove that any matrix with entries in a field can be brought into reduced
% echelon form by a sequence of elementary operations on rows. (This is what is
% more properly called Gaussian elimination.) [2.7, 2.9]
% 2.7. \lnot  Let M be a matrix with entries in a field and in reduced row echelon
% form
% ?
% (Exercise 2.6). Prove that if a row vector r is a linear combination
%  ai ri of the
% nonzero rows of M , then ai equals the component of r at the position
% corresponding
% to the pivot on the i-th row of M . Deduce that the nonzero rows of M are
% linearly
% independent. [2.9]
% 2.8. \lnot  Two matrices M , N are row-equivalent if M = P N for an invertible
% ma-
% trix P . Prove that this is indeed an equivalence relation, and that two
% matrices
% with entries in a field are row-equivalent if and only if one may be obtained
% from
% the other by a sequence of elementary operations on rows. [2.9, 2.12]
% 2.9. \lnot  Let k be a field, and consider row-equivalence (Exercise 2.8) on the
% set of m\times
% n matrices Mm,n (k). Prove that each equivalence class contains exactly one
% matrix
% in reduced row echelon form (Exercise 2.6). (Hint: To prove uniqueness, argue by
% contradiction. Let M , N be different row-equivalent reduced row echelon
% matrices;
% assume that they have the minimum number of columns with this property. If the
% leftmost column at which M and N differ is the k-th column, use the minimality
% to prove that M , N may be assumed to be of the form
% ?
%  ?
% ?
%  ?
% Ik-1 \ast 
% or
%  Ik-1 \ast  .
%  \ast 
% Use Exercise 2.7 to obtain a contradiction.)
% The unique matrix in reduced row echelon form that is row-equivalent to a
% given matrix M is called the reduced echelon form of M . [2.11]
% 2.10. ? The row space of a matrix M is the span of its rows; the column space
% of M is the span of its column. Prove that row-equivalent matrices have the same
% row space and isomorphic column spaces. [2.12, \S{}3.3]
% 2.11. Let k be a field and M \in  Mm,n (k). Prove that the dimension of the
% space
% spanned by the rows of M equals the number of nonzero rows in the reduced
% echelon
% form of M (cf. Exercise 2.9).
% 2.12. \lnot  Let k be a field, and consider row-equivalence on Mm,n (k)
% (Exercise 2.8).
% By Exercise 2.10, row-equivalent matrices have the same row space. Prove that,
% conversely, there is exactly one row-equivalence class in Mm,n (k) for each
% subspace
% of kn of dimension \leq  m. [2.13, 2.14]
% 2.13. \lnot  The set of subspaces of given dimension in a fixed vector space is
% called
% a Grassmannian. In Exercise 2.12 you have constructed a bijection between the
% Grassmannian of r-dimensional subspaces of k n and the set of reduced row
% echelon
% matrices with n columns and r nonzero rows.
% For r = 1, the Grassmannian is called the projective space. For a vector space V
% ,
% the corresponding projective space PV is the set of `lines' (1-dimensional
% subspaces)
% in V . For V = kn , PV may be denoted Pn-1
%  k
%  , and the field k may be omitted if it is
% clear from the context. Show that Pn-1
%  k
%  may
%  be written as a union k n-1 \cup  k n-2 \cup 
% \dots  \cup  k1
%  \cup  k0
%  , and describe each of these subsets `geometrically'.
% Thus, P n-1 is the union of n `cells'15 , the largest one having dimension n - 1
% (accounting for the choice of notation). Similarly, all Grassmannians may be
% written
% as unions of cells. These are called Schubert cells.
% Prove that the Grassmannian of (n - 1)-dimensional subspaces of kn admits a
% cell decomposition entirely analogous to that of Pn-1
%  k
%  . (This phenomenon will be
% explained in Exercise VIII.5.17.) [VII.2.20, VIII.4.7, VIII.5.17]
% 2.14. ? Show that the Grassmannian Grk (2, 4) of 2-dimensional subspaces of k4
% is
% the union of 6 Schubert cells: k4 \cup  k 3 \cup  k 2 \cup  k 2 \cup  k 1 \cup
% k 0 . (Use Exercise 2.12; list
% all the possible reduced echelon forms.) [VIII.4.8]
% 2.15. ? Prove that a square matrix with entries in a field is invertible if and
% only
% if it is equivalent to the identity, if and only if it is row-equivalent to the
% identity,
% if and only if its reduced echelon form is the identity. [\S{}2.3, 3.5]
% 2.16. Prove Proposition 2.10.
% 2.17. Prove Proposition 2.11.
% 2.18. Suppose \alpha  : Z3 \to  Z2 is represented by the matrix
% -6 12 18
% -15 36 54
% with respect to the standard bases. Find bases of Z3 , Z2 with respect to which
% \alpha
% is given by a matrix of the form obtained in Proposition 2.11.
% 15 Here, a `cell' is simply a subset endowed with a natural bijection with
% k\setminus  for some ?.

\section{Homomorphisms of free modules, II}
\label{sec:linalg:homomorphisms-of-free-modules-ii}


%Prove Corollary IV.6.5 again as a corollary of Proposition 2.11}
% In fact,

% prove the more general fact that every finitely generated abelian group is a
% direct
% sum of cyclic groups. [\S{}II.6.3, \S{}IV.6.1, \S{}IV.6.3, \S{}2.4, \S{}4.1,
% \S{}4.3]
% The work in \S{}2 accomplishes the goal of describing HomR (F, G), where F and G
% are free R-modules of finite rank; as we have seen, this can be done most
% explicitly
% if, for example, R is a field or a Euclidean domain. We are not quite done,
% though:
% even over fields, it is important to `understand the answer'---that is, examine
% the
% classification of homomorphisms into finitely many classes, highlighted at the
% end
% of \S{}2.3. Also, the frequent appearance of invertible matrices makes it
% necessary
% to develop criteria to tell whether a given matrix is or is not in the general
% linear
% group. I begin by pointing out a straightforward application of the
% classification of
% homomorphisms.

\subsection{Solving systems of linear equations}
\label{subsec:linalg:solving-systems-of-linear-equations}
% A moment's thought reveals that

% Gaussian elimination, through the auspices of statements such as Proposition
% 2.9,
% gives us a tool to solve systems of linear equations over a field or a Euclidean
% domain16 . This is another point which does not seem to warrant a careful
% treatment
% or memory gymnastic: writing things out carefully should take care of producing
% tools as need be. To describe a typical situation, suppose
% \{
% |
%  | a11x1 + \dots  + a1n xn = b1
% ...
% |
%  \}a m1x1
%  + \cdot \cdot \cdot  + amn xn
%  = b
% m
% is a system of m equations in n unknowns, with aij , bi in a Euclidean domain R.
% We want to find all solutions x1, . . . , xn in R.
% (
%  )
%  ( )
%  ( )
% a11 \dots  a1n
%  b1
%  x1
% With A = | ( .
%  ..
%  .
%  ..
%  .
%  .. |
%  ), b = ( |
%  .. .
%  ), |
%  and x = |
%  ( ... ), |
%  this amounts to
% am1 \dots  amn
%  bm
%  xn
% `solving for x' in the matrix equation
% A \cdot  x = b.
% Of course, if m = n and the matrix A is square and invertible, then the
% solutions
% are obtained simply as
% x = A-1 \cdot  b.
% Here A-1 may be obtained through Gaussian elimination (cf. Proposition 2.9) or
% through determinants (\S{}3.2).
% 16 Of course one can deal with systems over arbitrary integral domains R by
% reducing to the
% field case, by embedding R in its field of fractions.
% Even without requiring anything special of the `matrix of coefficients' A, row
% and column operations will take it to the standard form presented in Proposi-
% tion 2.11; that is, it will yield invertible matrices M , N such that
% (
%  )
% d1 \dots  0 0
% | . .
%  |
% M \cdot  A \cdot  N = |
%  | |
%  ..
%  . . ... ... |
%  |
%  |
% ( 0 \dots  dr 0 )
%  \cdot \cdot \cdot 
% with notation as in Proposition 2.11. Gaussian elimination is a constructive
% pro-
% cedure: watching yourself as you switch/combine rows or columns will produce M
% and N explicitly. Now, letting y = (yj ) and c = M b, the system
% (M \cdot  A \cdot  N ) \cdot  y = c
% solves itself:
%  \{
% |
%  |
%  |
%  |
%  |
%  |
%  |
%  d1y1 \cdot \cdot \cdot 
%  = c1
% dr yr = cr
% |
%  |
%  |
%  |
%  |
%  |
%  \}
%  0 = cr+1
% \cdot \cdot \cdot 
% has solutions if and only if dj | cj for all j and cj = 0 for j > r; moreover,
% in this
% case yj = d-1
%  j cj for j = 1, . . . , r, and yj is arbitrary for j > r. This yields y, and
% the reader will check that x = N y gives all solutions to the original system.
% Such arguments can be packaged into convenient explicit procedures to solve
% systems of linear equations. Again, it seems futile to list (or try to remember)
% any
% such procedure; it seems more important to know what is behind such techniques,
% so as to be able to come up with one when needed.
% In any case, the reader should be able to justify such recipes rigorously. The
% most famous one is possibly Cramer's rule (Proposition 3.6) which relies on de-
% terminants and is, incidentally, essentially useless in practice for any
% decent-size
% problem.

\subsection{The determinant}
\label{subsec:linalg:the-determinant}
% Let \alpha  : F \to  G be a homomorphism of free R-modules

% of the same rank, and let A be the matrix representing \alpha  with respect to a
% choice
% of bases for F and G. As a consequence of Lemma 2.3 (cf. Exercise 2.3), \alpha
% is an
% isomorphism if and only if A is a unit in Mn (R), that is, if and only if it is
% invertible
% as a matrix with entries in R. This may be detected by computing the determinant
% of A.
% Definition 3.1. Let A = (aij ) \in  M n(R) be a square matrix. Then the
% determinant
% of A is the element
% ?
%  ?
%  n
% det(A) =
%  (-1)\sigma 
%  ai\sigma (i) \in  R.
%  ?
% \sigma \in Sn
%  i=1
% Here Sn denotes the symmetric group on {1, . . . , n}, and I write17 \sigma (i)
% for the
% action of \sigma  \in  Sn on i \in  {1, . . . , n}; (-1)\sigma  is the sign of a
% permutation (Defini-
% tion IV.4.10, Lemma IV.4.12).
% The reader is surely familiar with determinants, at least over fields. They
% satisfy a number of remarkable properties; here is a selection:
% ---The determinant of a matrix A equals the determinant of its transpose At =
% (atij ), defined by setting
% atij = aji
% for all i and j (that is, the rows of At are the columns of A). Indeed,
% ?
%  ?
%  n
%  ?
%  ?
%  n
%  ?
%  ?
%  n
% det(At ) =
%  (-1)\sigma 
%  ati\sigma (i) =
%  (-1)\sigma 
%  a\sigma (i)i =
%  (-1)\sigma 
%  ai\sigma -1 (i)
% \sigma \in Sn
%  i=1
%  \sigma \in Sn
%  i=1
%  \sigma \in Sn
%  i=1
% -1by the commutativity of the product in R; and \sigma  ranges over all
% permutations
% of {1, . . . , n} as \sigma  does the same, and with the same sign, so the
% right-most term
% equals det(A).
% ---If two rows or columns of a square matrix A agree, then det(A) = 0. Indeed,
% it is enough to check this for matching columns; the case for rows follows by
% applying
% \setminus the previous observation. If columns j and j of A are equal, the contribution to
% det(A) due to a \sigma  \in  Sn is equal and opposite in sign to the
% contribution due to the
% product of \sigma  and the transposition (jj\setminus  ), so det(A) = 0.
% ---Suppose A = (aij ) and B = (bij ) agree on all but at most one row: aij = bij
% if i = k, for all j and some fixed k. Let cij := aij = bij for i = k, ckj := akj
% + bkj ,
% and let C := (c ij ). Then
% det(C) = det(A) + det(B).
% This follows immediately from Definition 3.1 and distributivity. Applying this
% observation to the transpose matrices gives an analogous statement for matrices
% differing at most along a column.
% I will record more officially the effect of elementary operations on
% determinants:
% Lemma 3.2. Let A be a square matrix with entries in an integral domain R.
% \bullet  Let A\setminus  be obtained from A by switching two rows or two columns. Then
% det(A\setminus  ) = - det(A).
% \bullet  Let A\setminus  be obtained from A by adding to a row (column) a multiple of another
% row (column). Then det(A\setminus  ) = det(A).
% \bullet  Let A\setminus  be obtained from A by multiplying a row (column) by an element18
% c \in  R. Then det(A\setminus  ) = c det(A).
% In other words, the effect of an elementary operation on det(A) is the same as
% multiplying det(A) by the determinant of the corresponding elementary matrix.
% 17 Consistency with previous encounters with the symmetric group (e.g.,
% \S{}IV.4) would demand
% that I write i\sigma ; but then I would end up with things like `aii\sigma  ',
% which are very hard to parse.
% 18For an elementary operation, c should be a unit; this restriction is not
% necessary here.
% Proof. These are all essentially immediate from Definition 3.1. For example,
% switching two columns amounts to correcting each \sigma  in the
%  ? definition by a fixed
% transposition, changing the sign of all contributions to the
%  in the definition. The
% third point is immediate from distributivity. Combining the third operation and
% the two remarks preceding the statement yields the second point. Details are
% left
% to the reader (Exercise 3.2).
%  ?
% This observation simplifies the theory of determinants drastically. If R = k is
% a field and P \in  Mn(k), Gaussian elimination (Proposition 2.10) shows
% ?
%  ?
% Ir 0
%  \setminus  \setminus A = E 1 \dots  Ea \cdot 
%  \cdot  E1 \dots  Eb ,
% 0 0
% where r \leq  n and Ei , Ej \setminus  are elementary matrices. Then Lemma 3.2
% gives that
% ?
%  ?
% ?
%  ?
%  \setminus 
%  Ir 0
% det(A) =
%  det(Ei )
%  det(E j ) det
%  .
% 0 0
% i
%  j
% (In particular, det A = 0 only if r = n.) Useful facts about the determinant
% follow
% from this remark. For example,
% Proposition 3.3. Let R be a commutative ring.
% \bullet  A square matrix A \in  Mn (R) is invertible if and only if det(A) is a unit in R.
% \bullet  The determinant is a homomorphism19 GLn (R) \to  (R\ast  , \cdot ). More generally,
% for A, B \in  Mn (R),
% det(A \cdot  B) = det(A) det(B).
% Proof for R = a field. If R = k is a field, we can use the considerations imme-
% diately preceding the statement. The first point is reduced to the case of a
% block
% matrix
%  ?
%  ?
% Ir 0
% ,
% 0 0
% for which it is immediate. In fact, this shows that det(A) = 0 if and only if
% the
% linear map kn \to  k n corresponding to A is not an isomorphism. In particular,
% for
% all A, B we have det(AB) = 0 if and only if AB is not an isomorphism, if and
% only
% if A or B is not an isomorphism, if and only if det(A) = 0 or det(B) = 0. So we
% only need to check the homomorphism property for invertible matrices. These are
% products of elementary matrices `on the nose', and the homomorphism property
% then follows from Lemma 3.2.
%  ?
% Before giving the (easy) extension to the case of arbitrary commutative rings,
% it is helpful to note the following explicit formulas. A `submatrix' obtained
% from a
% given matrix A by removing a number of rows and columns is called a minor of A;
% more properly, this term refers to the determinants of square submatrices
% obtained
% 19 Recall that (R \ast  , \cdot ) denotes the groups of units of R.
% in these ways. If A \in  Mn (R), the cofactors of A are the (n - 1) \times  (n -
% 1) minors
% of A, corrected by a sign. More precisely, for A = (aij ) I will let
% (
%  )
% a11
%  \cdot \cdot \cdot 
%  a1j-1
%  a1j+1
%  \cdot \cdot \cdot 
%  a1n
% | | ..
%  .
%  ..
%  .
%  ..
%  .
%  ..
%  .
%  ..
%  .
%  .. . |
%  |
% |
%  |
% A(ij) := (-1)i+j det |ai-11
%  |
%  \cdot \cdot \cdot 
%  ai-1j-1
%  ai-1j+1
%  \cdot \cdot \cdot 
%  ai-1n |
%  | .
% |ai+11
%  |
%  \cdot \cdot \cdot 
%  ai+1j-1
%  ai+1j+1
%  \cdot \cdot \cdot 
%  ai+1n |
%  |
% | ( .
%  ..
%  ..
%  .
%  ..
%  .
%  ..
%  .
%  ..
%  .
%  .. . |
%  )
% an1
%  \cdot \cdot \cdot 
%  anj-1
%  anj+1
%  \cdot \cdot \cdot 
%  ann
% Lemma 3.4. With notation as above,
% ?n
% \bullet  for all i = 1, . . . , n, det(A) = j=1 aij A(ij) ,
% ?n
% \bullet  for all j = 1, . . . , n, det(A) = i=1 aij A(ij).
% Proof. This is a simple (if slightly messy) induction on n, which I leave to the
% diligent reader.
%  ?
% Of course Lemma 3.4 is simply stating the famous strategy computing a deter-
% minant by expanding it with respect to your favorite row or column. This works
% wonderfully for very small matrices and is totally useless for large ones, since
% the
% number of computations needed to apply it grows as the factorial of the size of
% the
% matrix. From a computational point of view, it makes much better sense to apply
% Gaussian elimination and use the considerations preceding Proposition 3.3.
% But Lemma 3.4 has the following important implication:
% Corollary 3.5. Let R be a commutative ring and A \in  Mn (R). Then
% (
%  ) ( (11)
%  )
% A(11)
%  \cdot \cdot \cdot 
%  A(n1)
%  A
%  \cdot \cdot \cdot 
%  A(n1)
% A \cdot  | ( ..
%  .
%  ..
%  .
%  .. . | ) = | ( ..
%  .
%  ..
%  .
%  .. . | )
%  \cdot  A = det(A) In
%  .
% A(1n)
%  \cdot \cdot \cdot 
%  A (nn)
%  A(1n)
%  \cdot \cdot \cdot 
%  A(nn)
% Note the switch in the role of i and j in the matrix of cofactors. This matrix
% is called the adjoint matrix of A.
% Proof. Along the diagonal of the right-hand side, this is a restatement of Lemma
% 3.4.
% Off the diagonal, one is evaluating (for example)
% ?
%  n
% j=1
% ai?j A(ij)
% for i\setminus  = i. By Lemma 3.4 this is the same as the determinant of the
% matrix obtained
% by replacing the i-th row with the i\setminus  -th row; the resulting matrix has
% two equal rows,
% so its determinant is 0, as needed.
%  ?
% In particular, Corollary 3.5 proves that we can invert a matrix if we can invert
% its determinant:
%  ( (11)
%  )
% A
%  \dots  A(n1)
% A-1 = det(A)-1 |
%  ( ...
%  ..
%  .
%  .. . | )
%  ;
% A(1n) \dots  A(nn)
% this holds over any commutative ring, as soon as det(A) is a unit. In practice
% the computation of cofactors is `expensive', so in any given concrete case
% Gaussian
% elimination is likely a better alternative (at least over fields); cf. Exercise
% 3.5.
% But this formula for the inverse has good theoretical significance. For example,
% we are now in a position to complete the proof of Proposition 3.3.
% Proof of Proposition 3.3 for commutative rings. The first point follows from
% the second and from what we have just seen. Indeed, we have checked that
% A \in  Mn (R) admits an inverse A-1 \in  Mn (R) if det(A) is a unit in R;
% conversely,
% if A admits an inverse A-1 \in  Mn(R), then
% det(A) det(A-1) = det(AA-1 ) = det(In ) = 1
% by the second statement, so that det(A-1 ) is the inverse of det(A) in R.
% Thus, we just need to verify the second point, that is, the homomorphism
% property of determinants. In order to verify this over every (commutative) ring,
% it suffices to verify the `universal' identity obtained by writing out the
% claimed
% equality, for matrices with indeterminate entries. For example, for n = 2 the
% statement is
% x1
%  x2
%  y1
%  y2
%  x1y1
%  + x2 y3 x1 y2 + x2 y4
% det det = det ,
% x3 x4
%  y3 y4
%  x3y1 + x4 y3 x3 y2 + x4 y4
% which translates into the identity
% (x1x4 - x2x3 )(y1 y4 - y2y3 ) = (x1 y1 + x2 y3 )(x3 y2 + x4 y4 ) - (x 1 y2 + x2
% y 4 )(x3y1 + x4 y3);
% since this identity holds in Z[x1 , . . . , y4 ], it must hold in any
% commutative ring, for
% any choice of x1 , . . . , y4 : indeed, Z is initial in Ring.
% Now, we have verified that the homomorphism property holds over fields; in
% particular it holds over the field of fractions of Z[x11 , . . . , xnn, y11, . .
% . , ynn]. It
% follows that it does hold in Z[x11 , . . . , xnn , y11, . . . , ynn ], and we
% are done.
%  ?
% The `universal identity' argument extending the result from fields to arbitrary
% commutative rings is a useful device, and the reader is invited to contemplate
% it
% carefully.
% As an application of determinants (and especially cofactors) we can now go
% back to the special case of a system of n equations in n unknowns
% A \cdot  x = b,
% cf. \S{}3.1, in the case in which det(A) is a unit.
% Proposition 3.6 (Cramer's rule). Assume det(A) is a unit, and let A(j) be the
% matrix obtained by replacing the j-th column of A by the column vector b. Then
% xj = det(A)-1 det(A(j)).
% Proof. Using Lemma 3.4, expand det(A(j) ) with respect to the j-th column:
% ?
%  n
% det(A (j)) =
%  A(ij)bi .
% i=1
% Therefore
% (
%  )
%  ( (11)
%  ) ( )
% x1
%  A
%  \dots  A(n1)
%  b1
% ( | . .. |
%  ) = A-1
%  b = det(A) -1 | ( .
%  ..
%  ..
%  .
%  . .. ) | \cdot  | ( . .. )
%  |
% xn
%  A
% (1n)
%  \cdot \cdot \cdot  A
% (nn)
%  bn
% ( ?n
%  i=1 A
% (i1)
%  bi
%  ) (
%  det(A)-1
%  det(A(1) )
%  )
% = det(A)-1 |
%  (
%  .
%  ..
%  | ) = (
%  |
%  ..
%  .
%  ) |
%  ,
% ?n
%  (in)
%  -1
%  (n)
% i=1 A
%  bi
%  det(A) det(A )
% which gives the statement.
%  ?

\subsection{Rank and nullity}
\label{subsec:linalg:rank-and-nullity}
% According to Proposition 2.10, each equivalence class of

% matrices over a field has a representative of the type
% ?
%  ?
% Ir 0
% .
% 0 0
% The reason why two different m \times  n matrices of this type are surely
% inequivalent
% is that if \alpha  : V \to  W (with V and W free modules, vector spaces in this
% case) is
% represented by a matrix of this form, then r is the dimension of the image of
% \alpha .
% Therefore, matrices with different r cannot represent the same \alpha .
% This integer r is called the rank of the matrix and deserves some attention. I
% will discuss it for matrices with entries in a field, leaving generalizations to
% more
% general rings to the reader (for now at least).
% The column (row) space of a matrix P over a field k is the span of the columns
% (rows) of P . The column (row) rank of P is the dimension of the column (row)
% space of P .
% Proposition 3.7. The row rank of a matrix over a field k equals its column rank.
% Proof. Equivalent matrices have the same ranks. Indeed, let P \in  Mm,n (k); the
% row space of P consists of all row vectors
% ?
%  ? ?
%  ?
% a 1 \dots  am = v 1 \dots  vm \cdot  P
% ?
%  obtained
%  as each
%  ? v ? i ranges in k. ? If Q = M P N , with M and N invertible, let
% w1 \dots  wm = v1 \dots  vm \cdot  M -1 ; then
% ?
%  ?
%  ?
%  ?
%  -1 ?
%  ?
% w1 \dots  wm \cdot  Q = v1 \dots  vm M (M P N ) = a1 \dots  am \cdot  N.
% This shows that multiplication on the right by N maps the row space of P
% (isomor-
% phically, since N is invertible) to the row space of Q; thus the two spaces have
% the
% same dimension, as claimed. Minimal variations on the same argument show that
% the column ranks of P and Q agree. (Cf. Exercise 2.10.)
% Applying this observation and Proposition 2.10 reduces the question to matrices
% ?
%  ?
% Ir 0
% ,
% 0 0
% for which row rank = r = column rank, proving the statement.
%  ?
% The result upgrades easily to matrices over an arbitrary integral domain R
% (applying the usual trick of embedding R in its field of fractions).
% In view of Proposition 3.7, we can simply talk about the rank of a matrix:
% Definition 3.8. Let M \in  Mm,n (k) be a matrix over a field k. The rank of M is
% the dimension of its column (or, equivalently, row) space.
%  ?
% One way to rephrase Proposition 2.10 is that matrices over a field are
% classified
% up to equivalence by their rank.
% The foregoing considerations translate nicely in more abstract terms for a
% linear
% map \alpha  : V \to  W between finite-dimensional vector spaces over a field k.
% Using the
% convenient language introduced in \S{}III.7.1, note that each \alpha  determines
% an exact
% sequence of vector spaces
%  / ker \alpha 
%  / V
%  / im \alpha 
%  / 0 .
% Definition 3.9. The rank of \alpha , denoted rk \alpha , is the dimension of im
% \alpha . The nullity
% of \alpha  is dim(ker \alpha ).
%  ?
% Claim 3.10. Let \alpha  : V \to  W be a linear map of finite-dimensional vector
% spaces.
% Then
% (rank of \alpha ) + (nullity of \alpha ) = dim V.
% Proof. Let n = dim V and m = dim W . By Proposition 2.10 we can represent \alpha
% by an m \times  n matrix of the form
% ?
%  ?
% Ir 0
% .
% 0 0
% From this representation it is immediate that rk \alpha  = r and the nullity of
% \alpha  is n - r,
% with the stated consequence.
%  ?
% Summarizing, rk \alpha  equals the (column) rank of any matrix P representing
% \alpha ;
% similarly, the nullity of \alpha  equals `dim V minus the (row) rank' of P .
% Claim 3.10 is
% the abstract version of the equality of row rank and column rank.

\subsection{Euler characteristic and the Grothendieck group}
\label{subsec:linalg:euler-characteristic-and-the-grothendieck-group}
% Against my best ef-

% forts, I cannot resist extending these simple observations to more general
% complexes.
% Claim 3.10 may be reformulated as follows:
% Proposition 3.11. Let
%  / U
%  / V
%  / W
%  / 0
% be a short exact sequence of finite-dimensional vector spaces. Then
% dim(V ) = dim(U ) + dim(W ).
% Equivalently, this amounts to the relation dim(V /U ) = dim(V ) - dim(U ).
% Consider then a complex of finite-dimensional vector spaces and linear maps:
% V\bullet  :
%  / VN
%  \alpha N
%  / VN -1 \alpha N -1 / \dots 
%  \alpha 2
%  / V1
%  \alpha 1
%  / V0
%  / 0
% (cf. \S{}III.7.1). Thus, \alpha i-1 \circ  \alpha i = 0 for all i. This
% condition is equivalent to the
% requirement that im(\alpha i+1 ) \subseteq  ker(\alpha i ); recall that the
% homology of this complex is
% defined as the collection of spaces
% ker(\alpha i)
% Hi (V\bullet  ) =
%  .
% im(\alpha i+1 )
% The complex is exact if im(\alpha i+1 ) = ker(\alpha i ) for all i, that is, if
% Hi (V\bullet  ) = 0 for all i.
% Definition 3.12. The Euler characteristic of V\bullet  is the integer
% ?
% \chi (V\bullet  ) :=
%  (-1)i dim(Vi ).
%  ?
% i
% The original motivation for the introduction of this number is topological: with
% suitable positions, this Euler characteristic equals the Euler characteristic
% obtained
% by triangulating a manifold and then computing the number of vertices of the
% triangulation, minus the number of edges, plus the number of faces, etc.
% The following simple result is then a straightforward (and very useful) gener-
% alization of Proposition 3.11:
% Proposition 3.13. With notation as above,
% ?
%  N
% \chi (V\bullet  ) =
%  (-1)i dim(Hi (V\bullet )).
% i=0
% In particular, if V\bullet  is exact, then \chi (V\bullet ) = 0.
% Proof. There is nothing to show for N = 0, and the result follows directly from
% Proposition 3.11 if N = 1 (Exercise 3.15). Arguing by induction, given a complex
% V\bullet  :
%  / VN
%  \alpha N
%  / VN -1 \alpha N -1 / \dots 
%  \alpha 2
%  / V1
%  \alpha 1
%  / V0
%  / 0,
% we may assume that the result is known for `shorter' complexes. Consider then
% the
% truncation
% V\bullet  \setminus  :
%  / VN -1 \alpha N -1 / \dots 
%  \alpha 2
%  / V1
%  \alpha 1
%  / V0
%  / 0 .
% Then
% \chi (V\bullet  ) = \chi (V\bullet  \setminus  ) + (-1)N dim(VN ),
% and
% Hi (V\bullet ) = Hi(V\bullet  \setminus )
%  for 0 \leq  i \leq  N - 2,
% while
% ? ker(\alpha N -1 )
% HN -1 (V\bullet  ) = ker(\alpha N -1 ),
%  HN -1 (V\bullet  ) =
%  ,
%  HN (V\bullet  ) = ker(\alpha N ).
% im(\alpha N )
% By Proposition 3.11 (cf. Claim 3.10),
% dim(VN ) = dim(im(\alpha N )) + dim(ker(\alpha N ))
% and
% dim(HN -1(V\bullet  )) = dim(ker(\alpha N -1 )) - dim(im(\alpha N ));
% therefore
% dim(HN -1 (V\bullet  \setminus  )) - dim(VN ) = dim(HN -1 (V\bullet  )) - dim(HN
% (V\bullet )).
% Putting all of this together with the induction hypothesis,
% N
%  ?
%  -1
% \chi (V\bullet  \setminus  ) =
%  (-1)i dim(Hi (V\bullet  \setminus  ))
% i=0
% gives
% \chi (V\bullet ) = \chi (V\bullet  \setminus  ) + (-1)N dim(VN )
% N
%  ?
%  -1
% =
%  (-1)i dim(Hi (V\bullet  \setminus  )) + (-1)N dim(VN )
% i=0
% N
%  ?
%  -2
% =
%  (-1)i dim(Hi (V\bullet  \setminus  )) + (-1)N -1 (dim(HN -1 (V\bullet
%  \setminus  )) - dim(VN ))
% i=0
% N
%  ?
%  -2
% =
%  (-1)i dim(Hi (V\bullet  )) + (-1)N -1 (dim(HN -1 (V\bullet )) - dim(HN
%  (V\bullet  )))
% i=0
% ?
%  N
% i=
%  (-1) dim(Hi (V\bullet ))
% i=0
% as needed.
% ?
% In terms of the topological motivation recalled above, Proposition 3.13 tells us
% that the Euler characteristic of a manifold may be computed as the alternating
% sum
% of the ranks of its homology, that is, of its Betti numbers.
% Having come this far, I cannot refrain from mentioning the next, equally simple-
% minded, generalization. The reader has surely noticed that the only tool used in
% the proof of Proposition 3.13 was the `additivity' property of dimension,
% established
% in Proposition 3.11: if
%  / U
%  / V
%  / W
%  / 0
% is exact, then
% dim(V ) = dim(U ) + dim(W ).
% Proposition 3.13 is a formal consequence of this one property of dim.
% With this in mind, we can reinterpret what we have just done in the following
% curious way. Consider the category k-Vectf of finite-dimensional k-vector
% spaces.
% Each object V of k-Vectf determines an isomorphism class [V ]. Let F (k-Vectf )
% be
% the free abelian group on the set of these isomorphism classes; further, let E
% be the
% subgroup generated by the elements
% [V ] - [U ] - [W ]
% for all short exact sequences
%  / U
%  / V
%  / W
%  / 0
% in k-Vectf . The quotient group
% F (k-Vectf )
% K(k-Vectf ) :=
% E
% is called the Grothendieck group of the category k-Vectf . The element
% determined
% by V in the Grothendieck group is still denoted [V ].
% More generally, a Grothendieck group may be defined for any category admitting
% a notion of exact sequence.
% Every complex V\bullet  determines an element in K(k-Vectf ), namely
% ?
% \chi K (V\bullet  ) :=
%  (-1)i[Vi ] \in  K(k-Vectf ).
% i
% Claim 3.14. With notation as above, we have the following:
% \bullet  \chi K `is an Euler characteristic', in the sense that it satisfies the formula given
% in Proposition 3.13:
% ?
% \chi K (V\bullet  ) =
%  (-1)i [Hi(V\bullet  )].
% i
% \bullet  \chi K is a `universal Euler characteristic', in the following sense. Let G be an
% abelian group, and let \delta  be a function associating an element of G to each
% finite-
% dimensional vector space, such that \delta (V ) = \delta (V \setminus  ) if V
% \sim
%  = V \setminus  and \delta (V /U ) =
% \delta (V ) - \delta (U ). For V\bullet  a complex, define
% ?
% \chi G (V\bullet  ) =
%  (-1)i\delta (Vi).
% i
% Then \delta  induces a (unique) group homomorphism
% K(k-Vectf ) \to  G
% mapping \chi K (V\bullet  ) to \chi G (V\bullet  ).
% \bullet  In particular, \delta  = dim induces a group homomorphism
% K(k-Vectf ) \to  Z
% such that \chi  K (V\bullet ) ?\to  \chi (V\bullet ).
% \bullet  This is in fact an isomorphism.
% The best way to convince the reader that this impressive claim is completely
% trivial is to leave its proof to the reader (Exercise 3.16). The first point is
% proved
% by adapting the proof of Proposition 3.13; the second point is a harmless
% mixture
% of universal properties; the third point follows from the second; and the last
% point
% follows from the fact that dim(k) = 1 and the IBN property.
% The last point is, in fact, rather anticlimactic: if the impressively abstract
% Grothendieck group turns out to just be a copy of the integers, why bother
% defin-
% ing it? The answer is, of course, that this only stresses how special the
% category
% k-Vectf is. The definition of the Grothendieck group can be given in any context
% in
% which complexes and a notion of exactness are available (for example, in the
% cate-
% gory of finitely generated20 modules over any ring). The formal arguments
% proving
% Claim 3.14 will go through in any such context and provide us with a useful
% notion
% of `universal Euler characteristic'.
% We will come back to complexes and homology in Chapter IX.

\subsection*{Exercises}

% 3.1. Use Gaussian elimination to find all integer solutions of the system of
% equations
% ?
% 7x - 36y + 12z = 1,
% -8x + 42y - 14z = 2.
% 3.2. ? Provide details for the proof of Lemma 3.2. [\S{}3.2]
% 3.3. Redo Exercise II.8.8.
% 3.4. Formalize the discussion of `universal identities': by what cocktail of
% universal
% properties is it true that if an identity holds in Z[x1 , . . . , xr ], then it
% holds over
% every commutative ring R, for every choice of xi \in  R? (Is the commutativity
% of R
% necessary?)
% 3.5. ? Let A be an n \times  n square invertible matrix with entries in a field,
% and
% consider the n \times  (2n) matrix B = (A|In ) obtained by placing the identity
% matrix
% to the side of A. Perform elementary row operations on B so as to reduce A to In
% (cf. Exercise 2.15). Prove that this transforms B into (In |A-1 ).
% (This is a much more efficient way to compute the inverse of a matrix than by
% using determinants as in \S{}3.2.) [\S{}2.3, \S{}3.2]
% 3.6. \lnot  Let R be a commutative ring and M = ?m1 , .(
%  . . , m )
%  r ? a finitely
%  (
%  ) generated
% m1
% R-module. Let A \in  Mr (R) be a matrix such that A \cdot  | ( ... ) | = ( | ...
% |
%  ). Prove that
% mr
% det(A)m = 0 for all m \in  M . (Hint: Multiply by the adjoint.) [3.7]
% 3.7. \lnot  Let R be a commutative ring, M a finitely generated R-module, and
% let J
% be an ideal of R. Assume JM = M . Prove that there exists an element b \in  J
% such
% that (1 + b)M = 0. (Let m1 , . . . , mr be generators
%  for M . Find an r \times  r matrix
% (
%  )
%  ( )m1
%  m1
% B with entries in J such that ( | . .. ) |
%  = B \cdot  ( | . .. ). |
%  Then use Exercise 3.6.) [3.8,
% mr
%  mr
% VIII.1.18]
% 20 Note that some finiteness condition is likely to be necessary.
%  We cannot define a
% `Grothendieck group of k-Vect' because the isomorphism classes of objects in
% k-Vect do not form
% a set: there is one for each cardinal number, and cardinal numbers do not form a
% set.
% 3.8. \lnot  Let R be a commutative ring, M be a finitely generated R-module, and
% let
% J be an ideal of R contained in the Jacobson radical of R (Exercise V.3.14).
% Prove
% that M = 0 \iff  JM = M . (Use Exercise 3.7. This is Nakayama's lemma, a
% result with important applications in commutative algebra and algebraic
% geometry.
% A particular case was given as Exercise III.5.16.) [III.5.16, 3.9, 5.5]
% 3.9. \lnot  Let R be a commutative local ring, that is, a ring with a single
% maximal
% ideal m, and let M , N be finitely generated R-modules. Prove that if M = mM +N
% ,
% then M = N . (Apply Nakayama's lemma, that is, Exercise 3.8, to M/N . Note
% that the Jacobson radical of R is m.) [3.10]
% 3.10. \lnot  Let R be a commutative local ring, and let M be a finitely
% generated R-
% module. Note that M/mM is a finite-dimensional vector space over the field R/m;
% let m1 , . . . , mr \in  M be elements whose cosets mod mM form a basis of M/mM
% .
% Prove that m1 , . . . , mr generate M .
% (Show that ?m1 , . . . , mr ? + mM = M ; then apply Nakayama's lemma in the
% form of Exercise 3.9.) [5.5, VIII.2.24]
% 3.11. Explain how to use Gaussian elimination to find bases for the row space
% and
% the column space of a matrix over a field.
% 3.12. \lnot  Let R be an integral domain, and let M \in  Mm,n(R), with m < n.
% Prove
% that the columns of M are linearly dependent over R. [5.6]
% 3.13. Let k be a field. Prove that a matrix M \in  Mm,n (k) has rank \leq  r if
% and only
% if there exist matrices P \in  Mm,r (k), Q \in  Mr,n (k) such that M = P Q.
% (Thus the
% rank of M is the smallest such integer.)
% 3.14. Generalize Proposition 3.11 to the case of finitely generated free modules
% over any integral domain. (Embed the integral domain in its field of fractions.)
% 3.15. ? Prove Proposition 3.13 for the case N = 1. [\S{}3.4]
% 3.16. ? Prove Claim 3.14. [\S{}3.4]
% 3.17. Extend the definition of Grothendieck group of vector spaces given in
% \S{}3.4 to
% the category of vector spaces of countable (possibly infinite) dimension, and
% prove
% that it is the trivial group.
% 3.18. Let Abf g be the category of finitely generated abelian groups. Define a
% Grothendieck group of this category in the style of the construction of
% K(k-Vectf ),
% and prove that K(Abf g ) \sim 
%  = Z.
% 3.19. \lnot  Let Abf be the category of finite abelian groups. Prove that
% assigning to ev-
% ery finite abelian group its order extends to a homomorphism from the
% Grothendieck
% group K(Abf ) to the multiplicative group (Q\ast  , \cdot ). [3.20]
% 3.20. Let R-Modf be the category of modules of finite length (cf. Exercise 1.16)
% over a ring R. Let G be an abelian group, and let \delta  be a function
% assigning an
% element of G to every simple R-module. Prove that \delta  extends to a
% homomorphism
% from the Grothendieck group of R-Modf to G.
% Explain why Exercise 3.19 is a particular case of this observation.
% (For another example, letting \delta (M ) = 1 \in  Z for every simple module M
% shows
% that length itself extends to a homomorphism from the Grothendieck group of
% R-Modf to Z.)

\section{Presentations and resolutions}
\label{sec:linalg:presentations-and-resolutions}

% After this excursion into the idyllic world of free modules, we can come back to
% earth
% and see if we have learned something that may be useful for more general
% situations.
% Modules over a field are necessarily free (Proposition 1.7), not so for modules
% over
% more general rings. In fact, this property will turn out to be a
% characterization of
% fields (Proposition 4.10).
% It is important that we develop some understanding of nonfree modules. In this
% section we will see that homomorphisms of free modules carry enough information
% to allow us to deal with many nonfree modules.

\subsection{Torsion}
\label{subsec:linalg:torsion}
% There are several ways in which a module M may fail to be free:

% the most spectacular one is that M may have torsion.
% Definition 4.1. Let M be an R-module. An element m \in  M is a torsion element
% if {m} is linearly dependent, that is, if \exists r \in  R, r = 0, such that rm
% = 0. The
% subset of torsion elements of M is denoted TorR (M ). A module M is torsion-free
% if TorR (M ) = {0}. A torsion module is a module M in which every element is a
% torsion element.
%  ?
% The subscript R is usually omitted if there is no uncertainty about the base
% ring.
% A commutative ring is torsion-free as a module over itself if and only if it
% is an integral domain; this is a good reason to limit the discussion to integral
% domains in this chapter. Also, the reader will check (Exercise 4.1) that if R is
% an
% integral domain, then Tor(M ) is a submodule of M . Equally easy is the
% following
% observation:
% Lemma 4.2. Submodules and direct sums of torsion-free modules are torsion-free.
% Free modules over an integral domain are torsion-free.
% Proof. The first statement is immediate; the second follows from the first,
% since
% an integral domain is torsion-free as a module over itself.
%  ?
% Lemma 4.2 gives a good source of torsion-free modules: for example, ideals
% in an integral domain R are torsion-free (because they are submodules of the
% free
% module R1 ). In fact, ideals provide us with examples of another mechanism in
% which a module may fail to be free.
% Example 4.3. Let R = Z[x], and let I = (2, x). Then I is not a free R-module.
% More generally, let I be any nonprincipal ideal of an integral domain R; then I
% is
% a torsion-free module which is not free.
% Indeed, if I were free, then its rank would have to be 1 at most, by Propo-
% sition 1.9 (a basis for I would be a linearly independent subset of R, and R has
% rank 1 over itself); thus one element would suffice to generate I, and I would
% be
% principal.
%  ?
% What is behind this example is a characterization of PIDs in terms of `torsion-
% free submodules of a rank-1 free module' (Exercise 4.3). This is a facet of the
% main
% result towards which we are heading, that is, the classification of finitely
% generated
% modules over PIDs (Theorem 5.6). The gist of this classification is that
% finitely
% generated modules over a PID can be decomposed into cyclic modules. We have
% essentially proved this fact already for modules over Euclidean domains (it
% follows
% from Proposition 2.11; see Exercise 2.19), and we have looked in great detail at
% the
% particular case of Z-modules, a.k.a. abelian groups (\S{}IV.6); we are almost
% ready to
% deal with the general case of PIDs.
% Definition 4.4. An R-module M is cyclic if it is generated by a singleton, that
% is,
% if M \sim 
%  = R/I for some ideal I of R.
%  ?
% The equivalence in the definition is hopefully clear to our reader, as an
% immedi-
% ate consequence of the first isomorphism theorem for modules (Corollary
% III.5.16).
% If not, go back and (re)do Exercise III.6.16.
% Cyclic modules are witness to the difference between fields and more general
% rings: over a field k, a cyclic module is just a 1-dimensional vector space,
% that is a
% `copy of k'; over more general rings, cyclic modules may be very interesting
% (think
% of the many hours spent contemplating cyclic groups). In fact, we can tell that
% a
% ring is a field by just looking at its cyclic modules:
% Lemma 4.5. Let R be an integral domain. Assume that every cyclic R-module is
% torsion-free. Then R is a field.
% Proof. Let c \in  R, c = 0; then M = R/(c) is a cyclic module. Note that Tor(M )
% =
% M : indeed, the class of 1 generates R/(c) and belongs to Tor(M ) since c \cdot
% 1 is 0
% mod (c) and c = 0. However, by hypothesis M is torsion-free; that is, Tor(M ) =
% {0}. Therefore M = Tor(M ) is the zero module.
% This shows R/(c) is the zero R-module; that is, (c) = (1). Therefore, c is a
% unit. Thus every nonzero element of R is a unit, proving that R is a field.
%  ?
% Lemma 4.5 is a simple-minded illustration of the fact that we can study a ring R
% by studying the module structure over R, that is, the category R-Mod, and that
% we may not even need to look at the whole of R-Mod to be able to draw strong
% conclusions about R.

\subsection{Finitely presented modules and free resolutions}
\label{subsec:linalg:finitely-presented-modules-and-free-resolutions}
% The `right' way to

% think of a cyclic R-module M is as a module which admits an epimorphism from R,
% viewing the latter as the free rank-1 R-module21 :
% R1
%  / M
%  / 0.
% The fact that M is surjected upon by a free R-module is nothing special. In
% fact,
% every module M admits such an epimorphism:
% R\oplus A
%  / M
%  / 0
% provided that we are willing to take A large enough; if we are desperate, A = M
% will surely do. This is immediate from the universal property of free modules;
% if
% the reader does not agree, it is time to go back and review \S{}III.6.3. What
% makes
% cyclic modules special is that A can be chosen to be a singleton.
% We are now going to focus on a case which is also special, but not quite as
% special as cyclic modules: finitely generated modules are modules for which we
% can
% choose A to be a finite set (cf. \S{}III.6.4). Thus, we will assume that M
% admits an
% epimorphism from a finite-rank free module:
% Rm
%  \pi 
%  / M
%  / 0
% for some integer m. The image by \pi  of the m vectors in a basis of Rm is a set
% of
% generators for M .
% Finitely generated modules are much easier to handle than arbitrary modules.
% For example, an ideal of R can tell us whether a finitely generated module is
% torsion.
% Definition 4.6. The annihilator of an R-module M is
% Ann R (M ) := {r \in  R | \forall m \in  M, rm = 0}.
%  ?
% The subscript is usually omitted. The reader will check (Exercise 4.4) that
% Ann(M ) is an ideal of R and that if M is a finitely generated module and R is
% an
% integral domain, then M is torsion if and only if Ann(M ) = 0.
% We would like to develop tools to deal with finitely generated modules. It
% turns out that matrices allow us to describe a comfortably large collection of
% such
% modules.
% Definition 4.7. An R-module M is finitely presented if for some positive
% integers
% m, n there is an exact sequence
% R n
%  \phi 
%  / Rm
%  / M
%  / 0.
% Such a sequence is called a presentation of M .
%  ?
% In other words, finitely presented modules are cokernels (cf. III.6.2) of homo-
% morphisms between finitely generated free modules. Everything about M must
% be encoded in the homomorphism \phi ; therefore, we should be able to describe
% the
% module M by studying the matrix corresponding to \phi .
% There is a gap between finitely presented modules and finitely generated mod-
% ules, but on reasonable rings the two notions coincide:
% 21 In context the exactness of a sequence of R-modules will be understood, so
% the displayed
% sequence is a way to denote the fact that there exists a surjective homomorphism
% of R-modules
% from R to M ; cf. Example III.7.2. Also note the convention of denoting R by R1
% when it is viewed
% as a module over itself.
% Lemma 4.8. If R is a Noetherian ring, then every finitely generated R-module is
% finitely presented.
% Proof. If M is a finitely generated module, there is an exact sequence
% Rm
%  \pi 
%  / M
%  / 0
% for some m. Since R is Noetherian, Rm is Noetherian as an R-module (Corol-
% lary III.6.8). Thus ker \pi  is finitely generated; that is, there is an exact
% sequence
% Rn
%  / ker \pi 
%  / 0
% for some n. Putting together the two sequences gives a presentation of M .
%  ?
% Once we have gone one step to obtain generators and two steps to get a pre-
% sentation, we should hit upon the idea to keep going:
% Definition 4.9. A resolution of an R-module M by finitely generated free modules
% is an exact complex
% ...
%  / Rm3
%  / R m2
%  / Rm1
%  / Rm0
%  / M
%  / 0.
%  ?
% Iterating the argument proving Lemma 4.8 shows that if R is Noetherian, then
% every finitely generated module has a resolution as in Definition 4.9.
% It is an important conceptual step to realize that M may be studied by studying
% an exact complex of free modules
% ...
%  / R m3
%  / Rm2
%  / Rm1
%  / Rm0
% resolving M , that is, such that M is the cokernel of the last map. The Rm0
% piece
% keeps track of the generators of M ; Rm1 accounts for the relations among these
% generators; Rm2 records relations among the relations; and so on.
% Developing this idea in full generality would take us too far for now: for exam-
% ple, we would have to deal with the fact that every module admits many different
% resolutions (for example, we can bump up every mi by one by direct-summing each
% term in the complex with a copy of R 1 , sent to itself by the maps in the
% complex).
% We will do this very carefully later on, in Chapter IX.
% However, we can already learn something by considering coarse questions, such
% as `how long' a resolution can be. A priori, there is no reason to expect a free
% resolutions to be `finite', that is, such that mi = 0 for i ? 0. Such finiteness
% conditions tell us something special about the base ring R.
% The first natural question of this type is, for which rings R is it the case
% that
% every finitely generated R-module M has a free resolution `of length 0', that
% is,
% stopping at m0 ? That would mean that there is an exact sequence
%  / Rm0
%  / M
%  / 0.
% Therefore, M itself must be free. What does this say about R?
% Proposition 4.10. Let R be an integral domain. Then R is a field if and only if
% every finitely generated R-module is free.
% Proof. If R is a field, then every R-module is free, by Proposition 1.7. For the
% converse, assume that every finitely generated R-module is free; in particular,
% every
% cyclic module is free; in particular, every cyclic module is torsion-free. But
% then R
% is a field, by Lemma 4.5.
%  ?
% The next natural question concerns rings for which finitely generated modules
% admit free resolutions of length 1. It is convenient to phrase the question in
% stronger
% terms, that is, to require that for every finitely generated R-module M and
% every
% beginning of a free resolution
% R m0
%  \pi 
%  / M
%  / 0,
% the resolution can be completed to a length 1 free resolution. This would amount
% to demanding that there exist an integer m1 and an R-module homomorphism
% Rm1 \to  Rm0 such that the sequence
%  / Rm1
%  / R m0
%  \pi 
%  / M
%  / 0
% is exact. Equivalently, this condition requires that the module ker \pi  of
% relations
% among the m0 generators necessarily be free.
% Claim 4.11. Let R be an integral domain satisfying this property. Then R is
% a PID.
% Proof. Let I be an ideal of R, and apply the condition to M = R/I. Since we
% have an epimorphism
% R1
%  \pi 
%  / R/I
%  / 0,
% the condition says that ker \pi  is free; that is, I is free. Since I is a free
% submodule
% of R, which is free of rank 1, I must be free of rank \leq  1 by Proposition
% 1.9.
% Therefore I is generated by one element, as needed.
%  ?
% The classification result for finitely generated modules over PIDs (Theorem
% 5.6),
% which I keep bringing up, will essentially be a converse to Claim 4.11: the
% mysterious
% condition requiring free resolutions of finitely generated modules to have
% length at
% most 1 turns out to be a characterization of PIDs, just as the length 0
% condition is
% a characterization of fields (as proved in Proposition 4.10). We will work this
% out
% in \S{}5.2.

\subsection{Reading a presentation}
\label{subsec:linalg:reading-a-presentation}
% Let us return to the brilliant idea of studying a

% finitely presented module M by studying a homomorphism of free modules
% (*)
%  \phi  : Rn -\to  Rm
% such that M = coker \phi . As we know, we can describe \phi  completely by
% considering
% a matrix A representing it, and therefore we can describe any finitely presented
% module by giving a matrix corresponding to (a homomorphism corresponding to) it.
% In many cases, judicious use of the material developed in \S{}2 allows us to
% deter-
% mine the module M explicitly. For example, take
% (
%  )
% 1 3
% (2 3 ) ;
% 5 9
% this matrix corresponds to a homomorphism Z2 \to  Z 3 , hence to a Z-module,
% that is,
% a finitely generated abelian group G. The reader should figure out what G is
% more
% explicitly (in terms of the classification of \S{}IV.6; cf. Exercise 2.19)
% before reading
% on. In the rest of this section I will simply tie up loose ends into a more
% concrete
% recipe to perform these operations.
% Incidentally, a number of software packages can perform sophisticated oper-
% ations on modules (say, over polynomial rings); a personal favorite is
% Macaulay2.
% These packages rely precisely on the correspondence between modules and
% matrices:
% with due care, every operation on modules (such as direct sums, tensors,
% quotients,
% etc.) can be executed on the corresponding matrices. For example,
% Lemma 4.12. Let A, B be matrices with entries in an integral domain R, and let
% M , N denote the corresponding R-modules. Then M \oplus  N corresponds to the
% block
% matrix
%  ?
%  ?
% A 0
% .
% 0 B
% Proof. This follows immediately from Exercise 4.16.
%  ?
% Coming back to (*), note that the module M cannot know which bases we
% have chosen for Rn or Rm ; that is, M = coker \phi  really depends on the
% homomor-
% phism \phi , not on the specific matrix representation we have chosen for \phi .
% This is an
% issue that we have already encountered, and treated rather thoroughly, in
% \S{}2.2 and
% following: `equivalent' matrices represent the same homomorphism and hence the
% same module. In the context we are exploring now, Proposition 2.5 tells us that
% two matrices A, B represent the same module M if there exist invertible matrices
% P , Q such that B = P AQ.
% But this is not the whole story. Two different homomorphisms \phi 1 , \phi 2 may
% have isomorphic cokernels, even if they act between different modules: the
% extreme
% case being any isomorphism
% R m -\to  Rm ,
% whose cokernel is 0 (regardless of the isomorphism and no matter what m is).
% Therefore, if a matrix A\setminus  corresponds to a module M , then (by Lemma
% 4.12) so
% does the block matrix
%  ?
%  ?
% Ir 0
% A =
%  ,
% 0 A\setminus 
% where Ir is the r \times  r identity matrix (and r is any nonnegative integer);
% in fact, Ir
% could be replaced here by any invertible matrix.
% The following proposition attempts to formalize these observations.
% Proposition 4.13. Let A be a matrix with entries in an integral domain R, and
% let B be obtained from A by any sequence of the following operations:
% \bullet  switch two rows or two columns;
% \bullet  add to one row (resp., column) a multiple of another row (resp., column);
% \bullet  multiply all entries in one row (or column) by a unit of R;
% \bullet  if a unit is the only nonzero entry in a row (or column), remove the row and
% column containing that entry.
% Then B represents the same R-module as A, up to isomorphism.
% Proof. The first three operations are the `elementary operations' of \S{}2.3,
% and they
% transform a matrix into an equivalent one (by Proposition 2.7); as observed
% above,
% this does not affect the corresponding module, up to isomorphism.
% As for the fourth operation, if u is a unit and the only nonzero entry in (say)
% a
% row, then by applications of the second elementary operation we may assume that
% u is also the only nonzero entry in its column; without loss of generality we
% may
% assume that u is in fact the (1,1) entry of the matrix; that is, the matrix is
% in block
% form:
%  ?
%  ?
% u 0
% A =
%  .
% 0 A\setminus 
% But then A and A\setminus  represent the same module, as needed.
%  ?
% Example 4.14. The matrix with integer entries
% (
%  )
% 1 3
% (2 3)
% 5 9
% determines an abelian group G. Subtract three times the first column from the
% second column, obtaining
%  (
%  )
% 1 0
% (2 -3) ;
% 5 -6
% the (1, 1) entry is a unit and the only nonzero entry in the first row, so we
% can
% remove the first row and column:
% -3
% ;
% -6
% now change the sign, and subtract twice the first row from the second, leaving
% .
% Therefore G is isomorphic to the cokernel of the homomorphism
% \phi  : Z -\to  Z \oplus  Z
% mapping 1 to (3, 0). This homomorphism is injective and identifies Z with the
% subgroup 3Z \oplus  0 of the target. Therefore
% G \sim 
%  = coker \phi  \sim 
%  =
%  Z \oplus  Z \sim  =
%  Z
%  \oplus  Z.
%  ?
% 3Z \oplus  0
%  3Z

\subsection*{Exercises}

% By virtue of Gaussian elimination, the `algorithm' implicitly described in
% Propo-
% sition 4.13 will work without fail over Euclidean domains (e.g., over the
% polynomial
% ring in one variable over a field), in the sense that it will identify the
% finitely
% generated module corresponding to a matrix with an explicit direct sum of cyclic
% modules, as in Example 4.14. This is too much to expect over more general rings,
% since in general elementary transformations do not generate GL; cf. Remark 2.12.
% 4.1. ? Prove that if R is an integral domain and M is an R-module, then Tor(M )
% is a submodule of M . Give an example showing that the hypothesis that R is an
% integral domain is necessary. [\S{}4.1]
% 4.2. ? Let M be a module over an integral domain R, and let N be a torsion-free
% module. Prove that HomR(M, N ) is torsion-free. In particular, HomR (M, R) is
% torsion-free. (We will run into this fact again; see Proposition VIII.5.16.)
% [\S{}VIII.5.5]
% 4.3. ? Prove that an integral domain R is a PID if and only if every submodule
% of R itself is free. [\S{}4.1, 5.13]
% 4.4. ? Let R be a commutative ring and M an R-module.
% \bullet  Prove that Ann(M ) is an ideal of R.
% \bullet  If R is an integral domain and M is finitely generated, prove that M is torsion
% if and only if Ann(M ) = 0.
% \bullet  Give an example of a torsion module M over an integral domain, such that
% Ann(M ) = 0. (Of course this example cannot be finitely generated!)
% [\S{}4.2, \S{}5.3]
% 4.5. \lnot  Let M be a module over a commutative ring R. Prove that an ideal I
% of R
% is the annihilator of an element of M if and only if M contains an isomorphic
% copy
% of R/I (viewed as an R-module).
% The associated primes of M are the prime ideals among the ideals Ann(m),
% for m \in  M . The set of the associated primes of a module M is denoted AssR (M
% ).
% Note that every prime in AssR (M ) contains AnnR (M ). [4.6, 4.7, 5.16]
% 4.6. \lnot  Let M be a module over a commutative ring R, and consider the family
% of ideals Ann(m), as m ranges over the nonzero elements of M . Prove that the
% maximal elements in this family are prime ideals of R. Conclude that if R is
% Noetherian, then AssR(M ) = \emptyset  (cf. Exercise 4.5). [4.7, 4.9]
% 4.7. \lnot  Let R be a commutative Noetherian ring, and let M be a finitely
% generated
% module over R. Prove that M admits a finite series
% M = M0 ? M1 ? \dots  ? Mm = ?0?
% in which all quotients Mi /Mi+1 are of the form R/p for some prime ideal p of R.
% \setminus 
% (Hint: Use Exercises 4.5 and 4.6 to show that M contains an isomorphic copy M348
% of R/p1 for some prime p1 . Then do the same with M/M \setminus  , producing an
% M \setminus \setminus  \supseteq  M \setminus
% such that M \setminus \setminus  /M \setminus  = \sim 
%  R/p2 for some prime p2 . Why must this process stop after
% finitely many steps?) [4.8]
% 4.8. Let R be a commutative Noetherian ring, and let M be a finitely generated
% module over R. Prove that every prime in AssR (M ) appears in the list of primes
% produced by the procedure presented in Exercise 4.7. (If p is an associated
% prime,
% then M contains an isomorphic copy N of R/p. With notation as in the hint in
% Exercise 4.7, prove that either p1 = p or N \cap  M \setminus  = 0. In the
% latter case, N maps
% \setminus isomorphically to a copy of R/p in M/M ; iterate the reasoning.)
% In particular, if M is a finitely generated module over a Noetherian ring, then
% Ass(M ) is finite.
% 4.9. Let M be a module over a commutative Noetherian ring R. Prove that the
% union of all annihilators of nonzero elements of M equals the union of all
% associated
% primes of M . (Use Exercise 4.6.)
% Deduce that the union of the associated primes of a Noetherian ring R (viewed
% as a module over itself) equals the set of zero-divisors of R.
% 4.10. Let R be a commutative Noetherian ring. One can prove that the minimal
% primes of Ann(M ) (cf. Exercise V.1.9) are in Ass(M ). Assuming this, prove that
% the intersection of the associated primes of a Noetherian ring R (viewed as a
% module
% over itself) equals the nilradical of R.
% 4.11. Review the notion of presentation of a group, (\S{}II.8.2), and relate it
% to the
% notion of presentation introduced in \S{}4.2.
% 4.12. Let p be a prime ideal of a polynomial ring k[x1 , . . . , xn ] over a
% field k, and
% let R = k[x1 , . . . , xn ]/p. Prove that every finitely generated module over R
% has a
% finite presentation.
% 4.13. \lnot  Let R be a commutative ring. A tuple (a1, a2, . . . , an ) of
% elements of R is a
% regular sequence if a1 is a non-zero-divisor in R, a2 is a non-zero-divisor
% modulo22
% (a1), a3 is a non-zero-divisor modulo (a1 , a2), and so on.
% For a, b in R, consider the following complex of R-modules:
% (*)
%  / R
%  d2
%  / R \oplus  R
%  d1
%  / R
%  \pi 
%  /
%  (a,b)
%  R
%  / 0
% where \pi  is the canonical projection, d1 (r, s) = ra + sb, and d2 (t) = (bt,
% -at). Put
% otherwise, d1 and d2 correspond, respectively, to the matrices
% ?
%  ?
%  b
% a b ,
%  .
% -a
% \bullet  Prove that this is indeed a complex, for every a and b.
% \bullet  Prove that if (a, b) is a regular sequence, this complex is exact.
% The complex (*) is called the Koszul complex of (a, b). Thus, when (a, b) is a
% regular sequence, the Koszul complex provides us with a free resolution of the
% module R/(a, b). [4.14, 5.4, VIII.4.22]
% 22That is, the class of a
% 2 in R/(a1 ) is a non-zero-divisor in R/(a1 ).

\section{Classification of finitely generated modules over PIDs}
\label{sec:linalg:classification-of-finitely-generated-modules-over-pids}


%\lnot  A Koszul complex may be defined for any sequence a1, . . . , an of elements}

% of a commutative ring R. The case n = 2 seen in Exercise 4.13 and the case n = 3
% reviewed here will hopefully suffice to get a gist of the general construction;
% the
% general case will be given in Exercise VIII.4.22.
% Let a, b, c \in  R. Consider the following complex:
%  / R
%  d3
%  / R \oplus  R \oplus  R
%  d2
%  / R \oplus  R \oplus  R
%  d1
%  / R
%  \pi 
%  /
%  (a,b,c)
%  R
%  / 0
% where \pi  is the canonical projection and the matrices for d1 , d2, d3 are,
% respectively,
% (
%  )
%  ( )
% 0 -c -b
%  a
% ?
%  ?
% a b c , (-c 0
%  a ) , (-b) .
% b
%  a
%  c
% \bullet  Prove that this is indeed a complex, for every a, b, c.
% \bullet  Prove that if (a, b, c) is a regular sequence, this complex is exact.
% Koszul complexes are very important in commutative algebra and algebraic geom-
% etry. [VIII.4.22]

%View Z as a module over the ring R = Z[x, y], where x and y act by 0}

% Find a free resolution of Z over R. [VIII.4.21]

%Let \phi  : Rn \to  Rm and \psi  : Rp \to  Rq be two R-module homomorphisms,}

% and let
% \phi  \oplus  \psi  : Rn \oplus  Rp \to  Rm \oplus  Rq
% be the morphism induced on direct sums. Prove that
% coker(\phi  \oplus  \psi ) = coker \phi  \oplus  coker \psi .

%Determine (as a better known entity) the module represented by the matrix}

% (
%  )
% 1 + 3x
%  2x
%  3x
% (1 + 2x 1 + 2x - x 2 2x)
% x
%  x2
%  x
% over the polynomial ring k[x] over a field.
% It is finally time to prove the classification theorem for finitely generated
% modules
% over arbitrary PIDs. We have already proved this statement in the special case
% of
% finite Z-modules (\S{}IV.6), and the diligent reader has worked out a proof in
% the less
% special case of finitely generated modules over Euclidean domains, in Exercise
% 2.19.
% Now we go for the real thing.

\subsection{Submodules of free modules}
\label{subsec:linalg:submodules-of-free-modules}
% Recall (Lemma 4.2, Example 4.3) that a

% submodule of a free module over an arbitrary integral domain R is necessarily
% torsion-free but need not be free. For example, the ideal I = (x, y) of R = k[x,
% y]
% (with k a field, for example) is torsion-free as an R-module, but not free: for
% this
% to become really, really evident, it is helpful to rename x = a, y = b when
% these
% are viewed as elements of I and observe that a and b are not linearly
% independent
% over k[x, y], since ya - xb = 0.
% On the other hand, submodules of a free module over a field are automatically
% free: simply because every module over a field is free (Proposition 1.7). It is
% reasonable to expect that `some property in between' being a field and being a
% friendly UFD such as k[x, y] will guarantee that a submodule of a free module is
% free. We will now prove that this property is precisely that of being a
% principal
% ideal domain.
% Proposition 5.1. Let R be a PID, let F be a finitely generated free module over
% R,
% and let M \subseteq  F be a submodule. Then M is free.
% We will actually prove a more precise result, in view of the full statement of
% the classification theorem: we will show that there is a basis (x1 , . . . , xn)
% of F and
% elements a1 , . . . , am of R (with m \leq  n) such that
% y1 = a1 x1 , . . . ,
%  ym = am xm
% form a basis of M . That is, not only do we prove that M is free, but we also
% show
% that there are `compatible' bases of F and M .
% In order to do this, we may of course assume that M = 0: otherwise there is
% nothing to prove. Most of the work will then go into showing that if M = 0, we
% can
% split one direct summand off M ; iterating this process will prove the
% proposition23 .
% This is where the PID condition is used, so I will single out the main technical
% point
% into the following statement.
% Lemma 5.2. Let R be a PID, let F be a finitely generated free module over R, and
% let M \subseteq  F be a nonzero submodule. Then there exist a \in  R, x \in  F ,
% y \in  M , and
% submodules F \setminus  \subseteq  F and M \setminus  \subseteq  M , such that y
% = ax = 0, M \setminus  = F \setminus  \cap  M , and
% F = ?x? \oplus  F \setminus  ,
%  M = ?y? \oplus  M \setminus .
% The reader would find it instructive to pause a moment and try to imagine how
% the PID hypothesis may enter into the proof of this lemma. The PID hypothesis is
% a hypothesis on R, so the question is, where is the special copy of R to which
% we
% can apply it? Don't read on until you have spent a moment thinking about this.
% It is tempting to look for this copy of R among the `factors' of F = Rn : for
% example, we could map F to R by projecting onto the first component:
% \pi (r1 , . . . , rn ) = r1.
% But there is nothing special about the first component of R n; in fact there is
% nothing special about the particular basis chosen for F , that is, the
% particular
% 23 In a loose sense, this is the same strategy we employed in the proof of the
% classification
% result for finite abelian groups in \S{}IV.6.
% representation of F as Rn . Therefore, this does not look too promising. The way
% out of this bind is to democratically map F to R in every possible way: we
% consider
% all homomorphisms \phi  : F \to  R. This is a set that does not depend on any
% extra
% choice, so it is more likely to carry the information we need.
% For each \phi , \phi (M ) is a submodule of R, that is, an ideal of R; so we
% have a
% chance to use the PID hypothesis with profit. In fact, the much weaker fact that
% R is Noetherian guarantees already that there must be some homomorphism \alpha
% for
% which \alpha (M ) is maximal among all ideals \phi (M ). This copy of R, that
% is, the target
% of such a homomorphism \alpha , is special for a good reason, independent of
% inessential
% choices. The fact that R is a PID tells us that \alpha (M ) is principal, and
% then we are
% clearly in business.
% Here is the formal argument:
% Proof. For all \phi  \in  HomR (F, R), \phi (M ) is a submodule of R, that is,
% an ideal.
% The family of all these ideals is nonempty, and PIDs are Noetherian; therefore
% (by
% Proposition V.1.1) there exists a maximal element in the family, say \alpha (M
% ), for a
% homomorphism \alpha  : M \to  R. The fact that M = 0 implies immediately that
% some
% \phi (M ) = 0 (for example, take for \phi  the projection to a suitable factor of Rn ); hence
% \alpha (M ) = 0.
% Since R is a PID, \alpha (M ) is principal: \alpha (M ) = (a) for some a \in  R,
% a = 0. Since
% a \in  \alpha (M ), there exists an element y \in  M , y = 0, such that \alpha
% (y) = a. These are
% the elements a, y mentioned in the statement.
% I claim that a divides \phi (y) for all \phi  \in  HomR (F, R). Indeed, let b be
% a generator
% of (a, \phi (y)) (which exists since R is a PID; of course b is simply a gcd of
% a and \phi (y)),
% and let r, s \in  R such that b = ra+s\phi (y); consider the homomorphism \psi
% := r\alpha +s\phi .
% Since a \in  (b), we have \alpha (M ) \subseteq  (b). On the other hand
% b = ra + s\phi (y) = (r\alpha  + s\phi )(y) = \psi (y) \in  \psi (M );
% therefore (b) \subseteq  \psi (M ). It follows that \alpha (M ) \subseteq  \psi
% (M ), and by maximality \alpha (M ) =
% \psi (M ); hence (a) = (b), and in particular a divides \phi (y), as claimed.
% Let y = (s1 , . . . , sn ) as an element of F = Rn . Each s i is the image of y
% by a
% homomorphism F \to  R (that is, the i-th projection), so a divides all of them
% by
% what we just proved. Therefore \exists r1 , . . . , rn \in  R such that si = ari
% ; let
% x = (r1 , . . . , rn ) \in  F.
% This is the element x mentioned in the statement. By construction, y = ax. Fur-
% ther, a = \alpha (y) = \alpha (ax) = a\alpha (x); since R is an integral domain
% and a = 0, this
% implies \alpha (x) = 1.
% Finally, we let F \setminus  = ker \alpha  and M \setminus  = F \setminus  \cap
% M , and we can proceed to verify the
% stated direct sums.
% First, every z \in  F may be written as
% z = \alpha (z)x + (z - \alpha (z)x);
% by linearity
% \alpha (z - \alpha (z)x) = \alpha (z) - \alpha (z)\alpha (x) = \alpha (z) - \alpha (z) = 0,
% that is, z - \alpha (z)x \in  ker \alpha . This implies that F = ?x? + F
% \setminus . On the other hand,
% rx \in  F \setminus  \implies  \alpha (rx) = 0 \implies  r\alpha (x) = 0
% \implies  r = 0: that is, ?x? \cap  F \setminus  = 0.
% Therefore
% F = ?x? \oplus  F \setminus ,
% as claimed (cf. Exercise 5.1).
% Second, if z \in  M , then a divides \alpha (z): indeed, \alpha (z) \in  \alpha
% (M ) = (a). Writing
% \alpha (z) = ca, we have \alpha (z)x = cax = cy; splitting z as above, we note
% z - \alpha (z)x = z - cy \in  M \cap  F \setminus  = M \setminus  ,
% and this leads as before to
% concluding the proof.
% M = ?y? \oplus  M \setminus  ,
% ?
% Once Lemma 5.2 is established, the proof of Proposition 5.1 is mere busywork:
% Proof of Proposition 5.1. If M = 0, we are done. If not, applying Lemma 5.2
% to M \subseteq  F produces an element y1 \in  M and a submodule M (1) \subseteq
% M such that
% M = ?y1? \oplus  M (1) .
% If M (1) = 0, we are done; otherwise, apply Lemma 5.2 again (to M (1) \subseteq
% F ) to
% obtain y2 \in  M (1) and M (2) \subseteq  M (1) so that
% M = ?y1 ? \oplus  (?y2 ? \oplus  M (2)).
% This process may be continued, producing elements y1 , . . . , ym \in  M such
% that
% M = ?y1 ? \oplus  \dots  \oplus  ?ym ? \oplus  M (m) ,
% so long as the module M (m) is nonzero. However, by Proposition 1.9 we know
% m \leq  n, since y1 , . . . , ym are linearly independent in F . It follows that
% the process
% must stop; that is, M (m) = 0 for some m \leq  n. That is,
% M = ?y1 ? \oplus  \dots  \oplus  ?ym ?
% is free, as needed.
%  ?
% Note that the proof has not used part of the result of Lemma 5.2, that is, the
% fact that the `factor' ?y? of M is a submodule of a corresponding factor ?x? of
% F .
% This is needed in order to upgrade Proposition 5.1 along the lines mentioned
% after
% its statement. Here is that stronger statement:
% Corollary 5.3. Let R be a PID, let F be a finitely generated free module over R,
% and let M \subseteq  F be a submodule. Then there exist a basis (x1 , . . . , xn
% ) of F and
% nonzero elements a1, . . . , am of R (m \leq  n) such that (a1 x1, . . . , am xm
% ) is a basis
% of M . Further, we may assume a1 | a 2 | \dots  | am .
% This statement should be compared with Proposition 2.11: it amounts to a
% `Smith normal form over PIDs'; cf. Remark 2.12.
% Proof. Now that we know that submodules of a free module are free, we see that
% the submodule F \setminus  \subseteq  F produced in Lemma 5.2 is free. The first
% part of the
% statement then follows from Lemma 5.2, by an inductive argument analogous to
% the proof of Proposition 5.1, and is left to the reader.
% The most delicate part of the statement is the divisibility condition. By induc-
% tion it suffices to prove a1 | a2 , and for this we refer back to the proof of
% Lemma 5.2:
% (a1) is maximal among the ideals \phi (M ) of R, as \phi  ranges over all
% homomorphisms
% F \to  R. Consider then any24 homomorphism \phi  such that \phi (x1 ) = \phi (x2
% ) = 1.
% Since \phi (y1) = \phi (a1x1 ) = a1 , we have (a1 ) \subseteq  \phi (M ); by
% maximality (a1 ) = \phi (M ).
% Therefore a2 = \phi (a2 x2 ) = \phi (y2 ) \in  \phi (M ) = (a1 ), proving a1 |
% a2 as needed.
%  ?
% The logic of the argument given in this section is somewhat tricky: it takes
% Lemma 5.2 to prove Proposition 5.1; but then it takes Proposition 5.1 (proving
% \setminus that F is free) to revisit Lemma 5.2 and squeeze from it the stronger Corollary 5.3.

\subsection{PIDs and resolutions}
\label{subsec:linalg:pids-and-resolutions}
% Proposition 5.1 allows us to complete the circle of

% ideas begun in \S{}4.2.
% Proposition 5.4. Let R be an integral domain. Then R is a PID if and only if
% for every finitely generated R-module M and every epimorphism
% Rm0
%  \pi 0
%  / M
%  / 0,
% there exist a free R-module R m1 and a homomorphism \pi 1 : Rm1 \to  Rm0 such
% that
% the sequence
%  / Rm1
%  \pi 1
%  / Rm0
%  \pi 0
%  / M
%  / 0
% is exact.
% Proof. The fact that the stated condition implies that R is a PID was proved
% in Claim 4.11. For the converse, let \pi 0 : Rm0 \to  M be an epimorphism; then
% ker \pi 0 is free by Proposition 5.1; the result follows by choosing any
% isomorphism
% \pi 1 : R m1 \to  ker(\pi 0).
%  ?
% The content of Proposition 5.4 is possibly simpler than its awkward formulation:
% Proposition 5.4 is simply a characterization for PIDs analogous to the
% characteri-
% zation of fields worked out in Proposition 4.10. This is another example of the
% fact
% that we can study a ring R by looking at how R-Mod is put together.
% The reader can well imagine the `length-n' version of the condition appearing
% in these results: we may ask for the class of integral domains R with the
% property
% that for all finitely generated modules M and all partial free resolutions of M
% ,
% Rmn-1
%  \pi n-1
%  / \cdot \cdot \cdot 
%  \pi 1
%  / Rm0
%  \pi 0
%  / M
%  / 0 ,
% 24 In order to define a homomorphism on a free module F , one may define it on a
% basis of F
% and then extend it by linearity; there is no restriction on the possible choices
% for the images of
% basis elements. The reader who is not sure about this should review the
% universal property of free
% modules!
% there exists a homomorphism \pi n : Rmn \to  Rmn-1 such that
%  / Rmn
%  \pi n
%  / Rm n-1 \pi n-1 / \dots 
%  \pi 1
%  / Rm0
%  \pi 0
%  / M
% / 0
% is exact. We have proved that this condition characterizes fields for n = 0 and
% PIDs
% for n = 1.
% The alert reader should now have a d\'ej\`{a} vu, as we have already encountered
% a notion that is 0 precisely for fields and 1 for PIDs: the Krull dimension of a
% ring; cf. Example III.4.14. Of course this is not a coincidence, but the full
% relation
% between the Krull dimension and the bound on the length of free resolutions is
% be-
% yond the scope of this book25 . `Most' rings (even rings with finite Krull
% dimension)
% do not satisfy any such bound for all modules. The local rings satisfying the
% finite-
% ness condition described above for some n correspond in the language of
% algebraic
% geometry to `smooth' points on varieties of dimension \leq  n.

\subsection{The classification theorem}
\label{subsec:linalg:the-classification-theorem}
% The notion of the rank of a free module (Defi-

% nition 1.14) extends naturally to every finitely generated module M over an
% integral
% domain R.
% Definition 5.5. Let R be an integral domain. The rank rk M of a finitely
% generated
% R-module M is the maximum number of linearly independent elements in M . ?
% It should be clear that this number is finite (Exercise 5.6).
% Theorem 5.6. Let R be a PID, and let M be a finitely generated R-module. Then
% the following hold:
% \bullet  There exist distinct prime ideals (q1 ), . . . , (qn ) \subseteq  R, positive integers rij , and
% an isomorphism
%  (
%  )
% !
%  RM \sim 
%  = Rrk M \oplus  (
%  rij
%  )
%  .
% i,j
%  (q
% i )
% \bullet  There exist nonzero, nonunit ideals (a1 ), . . . , (am ) of R, such that (a1 ) \supseteq 
% (a2) \supseteq  \dots  \supseteq  (am), and an isomorphism
% M \sim 
%  = Rrk M \oplus 
%  R
%  \oplus  \cdot \cdot \cdot  \oplus 
%  R
%  .
% (a1)
%  (am )
% These decompositions are unique (in the evident sense).
% After all our preparatory work, this statement practically proves itself. I will
% describe the arguments in general terms, leaving the details to the enterprising
% reader.
% The two forms taken by the theorem go under the name of invariant factors
% and elementary divisors, as in the case of abelian groups. As in the case of
% abelian
% groups, the equivalence between the two formulations amounts to careful
% bookkeep-
% ing, and we will not prove it formally; see \S{}IV.6.2 for a reminder of how to
% go from
% 25 The most persistent of readers will take a more sophisticated look at these
% bounds in the
% exercises to \S{}IX.8.
% one to the other. The role played by Lemma IV.6.1 in that discussion is taken
% over
% by the Chinese remainder theorem, Theorem V.6.1, in this more general setting.
% As the two formulations are equivalent, it suffices to prove the existence and
% uniqueness of the decompositions given in Theorem 5.6 for any one of the two.
% The existence is a direct consequence of the results in \S{}5.1. Indeed, let M
% be a
% finitely generated module; thus there is an epimorphism
% Rn
%  \pi 
%  / M
%  / 0
% where n is the number of generators of M . Apply Corollary 5.3 to the submodule
% ker \pi  \subseteq  Rn : there exist a basis (x1, . . . , xn ) of Rn and nonzero
% elements a1 , . . . , am
% of R such that (a1 x1 , . . . , amxm ) is a basis of ker \pi  and further a1 |
% \dots  | am .
% That is, M is presented by the n \times  m matrix
% (
%  )
% a1 \dots 
% | | .. .
%  . .
%  .
%  .. . |
%  |
% |
%  |
% | |
%  0 \dots  am |
%  | .
% | 0 ...
%  0 |
% |
%  |
% | ( . ..
%  .
%  ..
%  .. . |
%  )
%  ...
% By Proposition 4.13, we may assume that a1 , . . . , am are not units: if any
% one of
% them is, the corresponding row and column may be omitted from the matrix (and
% n corrected accordingly). It follows that
% M \sim 
%  =
%  R
%  \oplus  \cdot \cdot \cdot  \oplus 
%  R
%  \oplus  R (n-m)
% (a 1 )
%  (am )
% with a1 | \dots  | am nonzero nonunits, as prescribed by Theorem 5.6, and the
% existence
% is proved.
% As for the uniqueness of the representations, if
% M \sim 
%  = Rr \oplus  T,
% with T a torsion module26 , then r = rk M and T \sim 
%  = TorR (M ) (Exercise 5.10). Thus,
% the `free part' and the `torsion part' in the decompositions given in Theorem
% 5.6
% are determined uniquely.
% This reduces the uniqueness question to the structure of torsion modules: as-
% suming that
% ! R
%  ! R
% (*)
%  rij \sim 
%  =
%  sk? ,
% i,j
%  (pi )
%  k,?
%  (qk )
% with pi, qk irreducible elements of R, the task is to show that the range of the
% indices is the same and, up to reordering, (pi) = (qi) and rij = sij for all i,
% j.
% Since a finitely generated module is torsion if and only if its annihilator is
% nonzero (Exercise 4.4), it is reasonable that Ann(M ) will be of some use here.
% 26 Recall that this means that every element of T is a torsion element; cf.
% \S{}4.1.
% Lemma 5.7. Let M be a torsion module, expressed as in Theorem 5.6 (with rk M =
% 0). Then Ann(M ) = (am). Further, the prime ideals (qi ) are precisely the prime
% ideals of R containing Ann(M ).
% Proof. By hypothesis
% M \sim 
%  =
%  R
%  \oplus  \cdot \cdot \cdot  \oplus 
%  R
%  ,
% (a1 )
%  (am )
% with a1 | \dots  | am . If r \in  Ann(M ), then
% 0 = r(1, . . . , 1) = (r, . . . , r).
% In particular r \equiv  0 modulo (am ); that is, r \in  (am ). Thus Ann(M )
% \subseteq  (a m ).
% For the reverse inclusion, assume r \in  (am ) and y \in  M . Identifying M with
% its
% decomposition, write y = (y1 , . . . , ym ), with yi \in  R/(ai ). Since r \in
% (a m ) \subseteq  (ai ), we
% have ryi = 0 for all i; therefore ry = 0, and r \in  Ann(M ) as needed.
% For the second part of the statement, tracing the equivalence between the two
% decompositions in Theorem 5.6 shows that the qi's are precisely the irreducible
% factors of am , so the assertion follows from the first part.
%  ?
% By Lemma 5.7, the sets {pi }, {qk } of irreducibles appearing in (*) must co-
% incide (up to inessential units): the ideals they generate are precisely the
% prime
% ideals containing Ann(M ). The fact that the sets coincide also follows from the
% observation that the isomorphism in (*) must match like primes, in the sense
% that
% an element of a factor
% R
% r
% ij(pi )
% in the source must land in a combination of quotients in the target by powers of
% the
% same prime pi ; this follows by comparing annihilators (cf. Exercise 5.11).
% Therefore,
% uniqueness is reduced to the case of a single irreducible q \in  R: it suffices
% to show
% that if
% R
%  \oplus  \dots  \oplus  R
%  = \sim 
%  R
%  \oplus  \dots  \oplus  R
%  ,
% (q r
%  )
%  (qrm )
%  (q s1 )
%  (qsn )
% with r1 \geq  \dots  \geq  rm and s1 \geq  \dots  \geq  sn , then m = n and ri =
% si for all i. This
% situation reproduces precisely the key step the reader used in the proof of
% uniqueness
% for the particular case of abelian groups, in Exercise IV.6.1. Of course I will
% not
% spoil the reader's fun by giving any more details this time (Exercise 5.12).
% Remark 5.8. To summarize, the information carried by a torsion module over a
% PID is equivalent to the choice of a selection of nonzero, nonunit ideals (a1 ),
% . . . , (am )
% such that
% (a1) \supseteq  (a2 ) \supseteq  \dots  \supseteq  (am ).
% We have seen that (am) is in fact the annihilator ideal of the module; it would
% be
% nice if we had a similarly explicit description of all invariants (a1 ), . . . ,
% (am ). As a
% preview of coming attractions, we could call the ideal
% (a1 \dots  am)
% the characteristic ideal of the module27. In situations in which we can compute
% both the annihilator and the characteristic ideal, comparing them may lead to
% 27 Warning: This does not seem to be standard terminology.

\subsection*{Exercises}

% strong conclusions about the module. For example, a torsion module is cyclic if
% and only if its annihilator and characteristic ideals coincide.
% Further, the prime ideals appearing in Theorem 5.6 have been characterized
% as the prime ideals containing the annihilator ideal. They may just as well be
% characterized as those prime ideals containing the characteristic ideal, as the
% reader
% will check (Exercise 5.15).
% Finally, note that from this point of view it is immediate that the
% characteristic
% ideal is contained in the annihilator ideal. We will run into this fact again,
% in an
% important application (Theorem 6.11).
%  ?
% 5.1. ? Let N , P be submodules of a module M , such that N \cap  P = {0} and
% M = N + P . Prove that M \sim 
%  = N \oplus  P . (This is a word-for-word repetition of
% Proposition IV.5.3 for modules.) [\S{}5.1]
% 5.2. Let R be an integral domain, and let M be a finitely generated R-module.
% Prove that M is torsion if and only if rk M = 0.
% 5.3. Complete the proof of Corollary 5.3.
% 5.4. Let R be an integral domain, and assume that a, b \in  R are such that a =
% 0,
% b \in 
%  / (a), and R/(a), R/(a, b) are both integral domains.
% \bullet  Prove that the Krull dimension of R is at least 2.
% \bullet  Prove that if R satisfies the finiteness condition discussed in \S{}5.2 for some n,
% then n \geq  2.
% You can prove this second point by appealing to Proposition 5.4. For a more
% concrete argument, you should look for an R-module admitting a free resolution
% of
% length 2 which cannot be shortened.
% \bullet  Prove that (a, b) is a regular sequence in R (Exercise 4.13).
% \bullet  Prove that the R-module R/(a, b) has a free resolution of length exactly 2.
% Can you see how to construct analogous situations with n \geq  3 elements a1, .
% . . , an ?
% 5.5. \lnot  Recall (Exercise V.4.11) that a commutative ring is local if it has
% a single
% maximal ideal m. Let R be a local ring, and let M be a direct summand of a
% finitely
% generated free R-module: that is, there exists an R-module N such that M \oplus
% N is
% a free R-module.
% \bullet  Choose elements m1 , . . . , mr \in  M whose cosets mod mM are a basis of M/mM
% as a vector space over the field R/m. By Nakayama's lemma, M = ?m1 , . . . , mr
% ?
% (Exercise 3.10).
% \bullet  Obtain a surjective homomorphism \pi  : F = R\oplus r \to  M .
% \bullet  Show that \pi  splits, giving an isomorphism F = \sim 
%  M \oplus  ker \pi . (Apply Exer-
% cise III.6.9 to the surjective homomorphism \pi  and the free module M \oplus  N
% to
% obtain a splitting M \to  F ; then use Proposition III.7.5.)
% \bullet  Show ker \pi /m ker \pi  = 0. Use Nakayama's lemma (Exercise 3.8) to deduce that
% ker \pi  = 0.
% \bullet  Conclude that M \sim 
%  = F is in fact free. [VIII.2.24, VIII.6.8, VIII.6.11]
% Summarizing, over a local ring, every direct summand of a finitely generated28
% free R-module is free. Using the terminology we will introduce in Chapter VIII,
% we
% would say that `projective modules over local rings are free'. This result has
% strong
% implications in algebraic geometry, since it underlies the notion of vector
% bundle.
% Contrast this fact with Proposition 5.1, which shows that, over a PID, every
% submodule of a finitely generated free module is free.
% 5.6. ? Let R be an integral domain, and let M = ?m1 , . . . , mr ? be a finitely
% gener-
% ated module. Prove that rk M \leq  r. (Use Exercise 3.12.) [\S{}5.3]
% 5.7. Let R be an integral domain, and let M be a finitely generated module over
% R.
% Prove that rk M = rk(M/ Tor(M )).
% 5.8. Let R be an integral domain, and let M be a finitely generated module over
% R.
% Prove that rk M = r if and only if M has a free submodule N \sim 
%  = Rr , such that
% M/N is torsion.
% If R is a PID, then N may be chosen so that 0 \to  N \to  M \to  M/N \to  0
% splits.
% 5.9. Let R be an integral domain, and let
%  / M1
%  / M2
%  / M3
%  / 0
% be an exact sequence of finitely generated R-modules. Prove that rk M2 = rk M1 +
% rk M3 .
% Deduce that `rank' defines a homomorphism from the Grothendieck group of
% the category of finitely generated R-modules to Z (cf. \S{}3.4).
% 5.10. ? Let R be an integral domain, M an R-module, and assume M = \sim 
%  Rr \oplus  T ,
% with T a torsion module. Prove directly (that is, without using Theorem 5.6)
% that
% r = rk M and T \sim 
%  = TorR (M ). [\S{}5.3]
% 5.11. ? Let R be an integral domain, let M , N be R-modules, and let \phi  : M
% \to  N
% be a homomorphism. For m \in  M , show that Ann(?m?) \subseteq  Ann(?\phi (m)?).
% [\S{}5.3]
% 5.12. ? Complete the proof of uniqueness in Theorem 5.6. (The hint in Exer-
% cise IV.6.1 may be helpful.) [\S{}5.3]
% 5.13. Let M be a finitely generated module over an integral domain R.
% Prove that if R is a PID, then M is torsion-free if and only if it is free.
% Prove
% that this property characterizes PIDs. (Cf. Exercise 4.3.)
% 5.14. Give an example of a finitely generated module over an integral domain
% which
% is not isomorphic to a direct sum of cyclic modules.
% 28 The finite rank hypothesis is actually unnecessary, but the proof is harder
% without this
% condition.

\section{Linear transformations of a free module}
\label{sec:linalg:linear-transformations-of-a-free-module}


%Prove that the prime ideals appearing in the elementary divisor version of}

% the classification theorem for a torsion module M over a PID are the prime
% ideals
% containing the characteristic ideal of M , as defined in Remark 5.8. [\S{}5.3]

%Prove that the prime ideals appearing in the elementary divisor version of}

% the classification theorem for a module M over a PID are the associated primes
% of M , as defined in Exercise 4.5.

\subsection{Let R be a PID}
\label{subsec:linalg:let-r-be-a-pid}
% Prove that the Grothendieck group (cf. \S{}3.4) of the category

% of finitely generated R-modules is isomorphic to Z.
% One beautiful application of the classification theorem for finitely generated
% mod-
% ules over PIDs is the determination of `special' forms for matrices of linear
% maps of
% a vector space to itself.
% Several fundamental concepts are associated with the general notion of a linear
% map from a vector space to itself, and we will review these concepts in this
% section.
% Not surprisingly, much of the discussion may be carried out for free modules
% over
% any integral domain R, and we will stay at this level of generality in (most of)
% the
% section. Theorem 5.6 will be used with great profit when R is a field, in the
% next
% section.

\subsection{Endomorphisms and similarity. Let R be an integral domain}
\label{subsec:linalg:endomorphisms-and-similarity-let-r-be-an-integral-domain}
% We have

% considered in some detail the module HomR (F, G) of R-module homomorphisms
% F \to  G
% between two free modules. For example, I have argued that describing such homo-
% morphisms for finitely generated free modules F , G amounts to describing
% matrices
% with entries in R, up to `equivalence': the equivalence relation introduced in
% \S{}2.2
% accounts for the (arbitrary) choice of bases for F and G.
% We now shift the focus a little and consider the special case in which F = G,
% that is, the R-module EndR (F ) of endomorphisms \alpha  of a fixed free
% R-module F :
% F
%  \alpha 
%  / F .
% Note that EndR (F ) is in fact an R-algebra: the operation of composition makes
% it
% a ring, compatibly with the R-module structure (cf. Exercise 2.3).
% From the point of view championed in \S{}2, the two copies of F appearing
% in EndR (F ) = HomR (F, F ) are unrelated; they are just (any) isomorphic repre-
% sentatives of a free module of a given rank. There are circumstances, however,
% in
% which we need to really choose one representative F and stick with it: that is,
% view
% \alpha  as acting from a selected free module F to itself, not just to an isomorphic copy
% of itself. In this situation we also say that \alpha  is a linear transformation
% of F , or an
% operator on F .
% From this different point of view it makes sense to compare elements of F
% `before and after' we apply \alpha ; for example, we could ask whether for some
% v \in  F
% we may have \alpha (v) = \lambda v for some \lambda  \in  R, or more generally
% whether a submodule
% M of F may be sent to itself by \alpha . In other words, we can compare the
% action of \alpha
% with the identity29 I : F \to  F .
% In terms of matrix representations (in the finite rank case) the description of
% \alpha
% can be carried out as we did in \S{}2, but with one interesting twist. In
% \S{}2.2 we dealt
% with how the matrix representation of a homomorphism changes when we change
% bases in the source and in the target. As we are now identifying source and
% target,
% we must choose the same basis for both source and target. This leads to a
% different
% notion of equivalence of matrices:
% Definition 6.1. Two square matrices A, B \in  Mn (R) are similar if they
% represent
% the same homomorphism F \to  F of a free rank-n module F to itself, up to the
% choice of a basis for F .
%  ?
% This is clearly an equivalence relation. The analog of Proposition 2.5 for this
% new notion is readily obtained:
% Proposition 6.2. Two matrices A, B \in  Mn (R) are similar if and only if there
% exists an invertible matrix P such that
% -1B = P AP .
% The reader who has really understood Proposition 2.5 will not need any detailed
% proof of this statement: it should be apparent from staring at the butterfly
% diagram
% Rn I
%  I I I \phi 
%  I I I
%  $
%  \alpha 
%  z
% u
%  uuu
%  \phi 
%  u
%  u
%  u
%  R n
% \pi 
%  \sqcup  u u u \psi 
%  u
%  u u
%  : F
%  / F \psi  I \sqcup 
%  \pi 
% dI
%  II
% II
% Rn
%  R n
% which I am essentially copying from \S{}2.2. The difference here is that we are
% choosing
% the same basis for source and target; hence the two triangles keeping track of
% the
% change of basis are the same. Suppose A (resp., B) represents \alpha  with
% respect to
% the choice of basis dictated by \phi  (resp., \psi ); that is, A (resp., B) is
% the matrix of
% the `top' (resp., bottom) composition
% \phi -1 \circ  \alpha  \circ  \phi  : Rn \to  Rn
% (resp., \psi -1 \circ  \alpha  \circ  \psi ). If P is the matrix representing
% the change of basis \pi , then
% -1B = P AP simply because the diagram commutes:
% \psi  -1 \circ  \alpha  \circ  \psi  = (\pi  \circ  \phi -1 ) \circ  \alpha  \circ  (\phi  \circ  \pi  -1 ) = \pi  \circ  (\phi -1 \circ  \alpha  \circ  \phi ) \circ  \pi -1 .
% Proposition 6.2 suggests a useful equivalence relation among endomorphisms:
% Definition 6.3. Two R-module homomorphisms of a free module F to itself,
% \alpha , \beta  : F \to  F,
% 29 The existence of the identity is one of the axioms of a category; every now
% and then, it is
% handy to be able to invoke it.
% are similar if there exists an automorphism \pi  : F \to  F such that
% \beta  = \pi  \circ  \alpha  \circ  \pi -1 .
% ?
% In the finite rank case, similar endomorphisms are represented by similar matri-
% ces, and two endomorphisms \alpha , \beta  are similar if and only if they may
% be represented
% by the same matrix by choosing appropriate (possibly different) bases on F . The
% interesting invariants determined by similar endomorphisms will turn out (not
% sur-
% prisingly) to be the same.
% From a group-theoretic viewpoint, similarity is an eminently natural notion
% to study: the group GL(F ) = AutR (F ) of automorphisms of a free module acts
% on EndR (F ) by conjugation, and two endomorphisms \alpha , \beta  are similar
% if and only
% if they are in the same orbit under this action. If the reader prefers matrices,
% the
% group GLn (R) of invertible n \times  n matrices with entries in R acts by
% conjugation on
% the module Mn (R) of square n \times  n matrices, and two matrices are similar
% if and
% only if they are in the same orbit.
% Natural questions arise in this context. For example, we should look for ways
% to distinguish different orbits: given two endomorphisms \alpha , \beta , can we
% effectively
% decide whether \alpha  and \beta  are similar? One way to approach this question
% is to
% determine invariants which can distinguish different orbits. Better still, we
% can
% look for `special' representatives within each orbit, that is, of a given
% similarity
% class. In other words, given an endomorphism \alpha  : F \to  F , find a basis
% of F with
% respect to which the matrix description of \alpha  has a particular, predictable
% shape.
% Then \alpha , \beta  are similar if these distinguished representatives
% coincide.
% This is what we are eventually going to squeeze out of Theorem 5.6, in the
% friendly case of vector spaces over a fixed field.

\subsection{The characteristic and minimal polynomials of an endomorphism}
\label{subsec:linalg:the-characteristic-and-minimal-polynomials-of-an-endomorphism}

% Let \alpha  \in  EndR(F ) be an endomorphism of a free R-module; henceforth I am
% often
% tacitly going to assume that F is finitely generated (this is necessary anyway
% as
% we are aiming to translate what we do into the language of matrices). We want
% to identify invariants of the similarity class of \alpha , that is, quantities
% that will not
% change if we replace \alpha  with a similar linear map \beta  \in  EndR (F ).
% For example, the determinant is such a quantity:
% Definition 6.4. Let \alpha  \in  End R (F ). The determinant of \alpha  is
% det(\alpha ) := det(A),
% where A is the matrix representing \alpha  with respect to any choice of basis
% of F . ?
% Of course there is something to check here, that is, that det(A) does not depend
% on the choice of basis. But we know (Proposition 6.2) that A, B represent the
% same endomorphism \alpha  if and only if there exists an invertible matrix P
% such that
% -1B = P AP ; then
% -1 -1det(B) = det(P AP ) = det(P ) det(A) det(P ) = det(A)
% by Proposition 3.3. Thus the determinant is indeed independent of the choice of
% basis. Essentially the same argument shows that if \alpha  and \beta  are
% similar linear
% transformations, then det(\alpha ) = det(\beta ) (Exercise 6.3).
% By Proposition 3.3, a linear transformation \alpha  is invertible if and only if
% det(\alpha )
% is a unit in R. If R is a field, this of course means simply det(\alpha ) = 0.
% Even if R is
% not a field, det(\alpha ) = 0 says something interesting:
% Proposition 6.5. Let \alpha  be a linear transformation of a free R-module F =
% \sim
%  Rn .
% Then det(\alpha ) = 0 if and only if \alpha  is injective.
% Proof. Embed R in its field of fractions K, and view \alpha  as a linear
% transformation
% of K n ; note that the determinant of \alpha  is the same whether it is computed
% over R
% or over K. Then \alpha  is injective as a linear transformation Rn \to  Rn if
% and only if it
% is injective as a linear transformation K n \to  K n , if and only if it is
% invertible as a
% linear transformation K n \to  K n , if and only if det(\alpha ) = 0.
%  ?
% Of course over integral domains other than fields \alpha  may fail to be
% surjective even
% if det(\alpha ) = 0; and care is required even over fields, if we abandon the
% hypothesis
% that the free modules are finitely generated (cf. Exercises 6.4 and 6.5).
% Another quantity that is invariant under similarity is the trace. The trace of a
% square matrix A = (aij ) \in  Mn (R) is
% ?
%  n
% tr(A) :=
%  aii ,
% i=1
% that is, the sum of its diagonal entries.
% Definition 6.6. Let \alpha  \in  EndR (F ). The trace of \alpha  is defined to
% be tr(\alpha ) := tr(A),
% where A is the matrix representing \alpha  with respect to any choice of basis
% of F . ?
% Again, we have to check that this is independent of the choice of basis; the key
% is the following computation.
% Lemma 6.7. Let A, B \in  Mn (R). Then tr(AB) = tr(BA).
% ?n
% Proof. Let A = (aij ), B = (bij ). Then AB = (
%  k=1 aik bkj ); hence
% ?
%  n ?
%  n
% tr(AB) =
%  aik bki.
% i=1 k=1
% This expression is symmetric in A, B, so it must equal tr(BA).
%  ?
% With this understood, if B = P AP -1, then
% tr(B) = tr((P A)P -1) = tr(P -1 (P A)) = tr((P -1 P )A) = tr(A),
% showing (by Proposition 6.2) that similar matrices have the same trace, as
% needed.
% Again, the reader will check that the same token shows that similar linear
% transformations have the same trace.
% The trace and determinant of an endomorphism \alpha  of a rank-n free module are
% in fact just two of a sequence of n invariants, which can be nicely collected
% together
% into the characteristic polynomial of \alpha .
% Definition 6.8. Let F be a free R-module, and let \alpha  \in  EndR (F ). Denote
% by I
% the identity map F \to  F . The characteristic polynomial of \alpha  is the
% polynomial
% P\alpha  (t) := det(tI - \alpha ) \in  R[t].
%  ?
% Proposition 6.9. Let F be a free R-module of rank n, and let \alpha  \in  EndR
% (F ).
% \bullet  The characteristic polynomial P\alpha  (t) is a monic polynomial of degree n.
% \bullet  The coefficient of tn-1 in P\alpha  (t) equals - tr(\alpha ).
% \bullet  The constant term of P\alpha  (t) equals (-1)n det(\alpha ).
% \bullet  If \alpha  and \beta  are similar, then P\alpha  (t) = P\beta  (t).
% Proof. The first point is immediate, and the third is checked by setting t = 0.
% To
% verify the second assertion, let A = (aij ) be a matrix representing \alpha
% with respect
% to any basis for F , so that
% (
%  )
% t - a11 -a12 . . .
%  -a1n
% | |
%  -a21 t - a22 . . .
%  -a2n |
%  |
% P\alpha  (t) = det | ( ..
%  .
%  .
%  ..
%  .
%  ..
%  . .. )
%  | .
% -an1
%  -an2 . . . t - ann
% Expanding the determinant according to Definition 3.1 (or in any other way), we
% see that the only contributions to the coefficient of tn-1 come from the
% diagonal
% entries, in the form
% ?
%  n
% t \dots  t \cdot  (-aii) \cdot  t \dots  t,
% i=1
% and the statement follows.
% Finally, assume that \alpha  and \beta  are similar. Then there exists an
% invertible \pi
% such that \beta  = \pi  \circ  \alpha  \circ  \pi  -1 , and hence
% tI - \beta  = \pi  \circ  (tI - \alpha ) \circ  \pi  -1
% are similar (as endomorphisms of R[t]n ). By Exercise 6.3 these two
% transformations
% must have the same determinant, and this proves the fourth point.
%  ?
% By Proposition 6.9, all coefficients in the characteristic polynomial
% tn - tr(\alpha )tn-1 + \dots  + (-1)ndet(\alpha )
% are invariant under similarity; as far as I know, trace and determinant are the
% only
% ones that have special names.
% Determinants, traces, and more generally the characteristic polynomial can
% show very quickly that two linear transformations are not similar; but do they
% tell
% us unfailingly when two transformations are similar? In general, the answer is
% no,
% even over fields. For example, the two matrices
% 1 0
%  1 1
% I =
%  ,
%  A =
% 0 1
%  0 1
% both have characteristic polynomial (t - 1)2, but they are not similar. Indeed,
% I is
% -1the identity and clearly only the identity is similar to the identity: P IP =
% I for
% all invertible matrices P .
% Therefore, the equivalence relation defined by prescribing that two transforma-
% tions have the same characteristic polynomial is coarser than similarity. This
% hints
% that there must be other interesting quantities associated with a linear
% transforma-
% tion and invariant under similarity.
% We are going to understand this much more thoroughly in a short while (at least
% over fields), but we can already gain some insight by contemplating another kind
% of `polynomial' information, which also turns out to be invariant under
% similarity.
% As EndR (F ) is an R-algebra, we can evaluate every polynomial
% f (t) = rm tm + rm-1 tm-1 + \dots  + r0 \in  R[t]
% at any \alpha  \in  EndR (F ):
% f (\alpha ) = rm \alpha m + rm-1 \alpha m-1 + \dots  + r0 \in  EndR (F ).
% In other words, we can perform these operations in the ring EndR (F );
% multiplica-
% tion by r \in  R amounts to composition with rI \in  EndR (F ), and \alpha k
% stands for the
% k-fold composition \alpha  \circ  \dots  \circ  \alpha  of \alpha  with itself.
% The set of polynomials such that
% f (\alpha ) = 0 is an ideal of R[t] (Exercise 6.7), which I will denote I\alpha
% and call the
% annihilator ideal of \alpha .
% Lemma 6.10. If \alpha  and \beta  are similar, then I\alpha  = I\beta  .
% Proof. By hypothesis there exists an invertible \pi  such that \beta  = \pi
% \circ  \alpha  \circ  \pi -1 . As
% \beta  k = (\pi  \circ  \alpha  \circ  \pi  -1 )k = (\pi  \circ  \alpha  \circ  \pi  -1 ) \circ  (\pi  \circ  \alpha  \circ  \pi -1 ) \circ  \dots  \circ  (\pi  \circ  \alpha  \circ  \pi -1 ) = \pi  \circ  \alpha k \circ  \pi  -1 ,
% we see that for all f (t) \in  R[t] we have
% f (\beta ) = \pi  \circ  f (\alpha ) \circ  \pi  -1 .
% It follows immediately that f (\alpha ) = 0 \iff  f (\beta ) = 0, which is the
% statement.
%  ?
% Going back to the simple example shown above, the polynomial t - 1 is in the
% annihilator ideal of the identity, while it is not in the annihilator ideal of
% the matrix
% 1 1
% A =
%  .
% 0 1
% An optimistic reader might now guess that two linear transformations are sim-
% ilar if and only if both their characteristic polynomials and annihilator ideals
% coin-
% cide. This is unfortunatley not the case in general, but I promise that the
% situation
% will be considerably clarified in short order (cf. Exercise 7.3).
% In any case, even the simple example given above allows me to point out a
% remarkable fact. Note that the (common) characteristic polynomial (t - 1)2 of I
% and A annihilates both:
%  0 0
% (A - I)2 =
%  =
%  .
%  0 0
% This is not a coincidence: P\alpha  (t) \in  I\alpha  for all linear
% transformations \alpha . That is,
% Theorem 6.11 (Cayley-Hamilton). Let P\alpha  (t) be the characteristic
% polynomial of
% the linear transformation \alpha  \in  End R (F ). Then
% P\alpha  (\alpha ) = 0.
% This beautiful observation can be proved directly by judicious use of Cramer's
% rule30 , in the form of Corollary 3.5; cf. Exercise 6.9. In any case, the
% Cayley-
% Hamilton theorem will become essentially evident once we connect these linear
% algebra considerations with the classification theorem for finitely generated
% modules
% over a PID; the adventurous reader can already look back at Remark 5.8 and
% figure
% out why the Cayley-Hamilton theorem is obvious.
% If R is an arbitrary integral domain, we cannot expect too much of R[t], and it
% seems hard to say something a priori concerning I\alpha  . However, consider the
% field of
% fractions K of R (\S{}V.4.2); viewing \alpha  as an element of EndK (K n ) (that
% is, viewing
% the entries of a matrix representation of \alpha  as elements of K, rather than
% R), \alpha  will
% have an annihilator ideal I\alpha  (K)
%  `over K', and it is clear that
% I\alpha  = I \alpha  (K) \cap  R[t].
% The advantage of considering I\alpha  (K)
%  \subseteq  K[t] is that K[t] is a PID, and it follows that
% I\alpha  (K)
%  has a (unique) monic generator.
% Definition 6.12. Let F be a free R-module, and let \alpha  \in  EndR (F ). Let K
% be
% the field of fractions of R. The minimal polynomial of \alpha  is the monic
% generator
% m\alpha  (t) \in  K[t] of I\alpha  (K)
%  .
%  ?
% With this terminology, the Cayley-Hamilton theorem amounts to the assertion
% that the minimal polynomial divides the characteristic polynomial: m\alpha  (t)
% | P\alpha  (t).
% Of course the situation is simplified, at least from an expository point of
% view,
% if R is itself a field: then K = R, m\alpha  (t) \in  R[t], and I\alpha  =
% (m\alpha  (t)). This is one
% reason why we will eventually assume that R is a field.

\subsection{Eigenvalues, eigenvectors, eigenspaces}
\label{subsec:linalg:eigenvalues-eigenvectors-eigenspaces}

% Definition 6.13. Let F be a free R-module, and let \alpha  \in  EndR (F ) be a
% linear
% transformation of F . A scalar \lambda  \in  R is an eigenvalue for \alpha  if
% there exists v \in  F ,
% v = 0, such that
% \alpha (v) = \lambda v.
%  ?
% For example, 0 is an eigenvalue for \alpha  precisely when \alpha  has a
% nontrivial kernel.
% The notion of eigenvalue is one of the most important in linear algebra, if not
% in
% algebra, if not in mathematics, if not in the whole of science. The set of
% eigenvalues
% of a linear transformation is called its spectrum. Spectra of operators show up
% everywhere, from number theory to differential equations to quantum mechanics.
% The spectrum of a ring, encountered briefly in \S{}III.4.3, was so named because
% it
% may be interpreted (in important motivating contexts) as a spectrum in the sense
% 30 Here is an even more direct `proof': P\alpha 
%  (\alpha ) = det(\alpha I - \alpha ) = det(0) = 0. Unfortunately this
% does not work: in the definition det(tI - \alpha ) of the characteristic
% polynomial, t is assumed to act
% as a scalar and cannot be replaced by \alpha . Applying Cramer's rule
% circumvents this obstacle.
% of Definition 6.13. The `spectrum of the hydrogen atom' is also a spectrum in
% this
% sense.
% It is hopefully immediate that similar transformations have the same spectrum:
% for if \beta  = \pi  \circ  \alpha  \circ  \pi  -1 and \alpha (v) = \lambda v,
% then
% \beta (\pi (v)) = \pi  \circ  \alpha  \circ  \pi  -1(\pi (v)) = \pi  \circ  (\alpha (v)) = \pi (\lambda v) = \lambda \pi (v);
% as \pi (v) = 0 if v = 0, this shows that every eigenvalue of \alpha  is an
% eigenvalue of \beta .
% If F is finitely generated, then we have the following useful translation of the
% notion of eigenvalue:
% Lemma 6.14. Let F be a finitely generated R-module, and let \alpha  \in  EndR (F
% ). Then
% the set of eigenvalues of \alpha  is precisely the set of roots in R of the
% characteristic
% polynomial P\alpha  (t).
% Proof. This is a straightforward consequence of Proposition 6.5:
% \lambda  is an eigenvalue for \alpha  \iff  \exists v = 0 such that \alpha (v) = \lambda I(v)
% \iff  \exists v = 0 such that (\lambda I - \alpha )(v) = 0
% \iff  \lambda I - \alpha  is not injective
% \iff  det(\lambda I - \alpha ) = 0
% \iff  P\alpha  (\lambda ) = 0
% as claimed.
% ?
% It is an evident consequence of Lemma 6.14 that eigenvalue considerations
% depend very much on the base ring R.
% Example 6.15. The matrix
% 0 -1
% 1 0
% has no eigenvalues over R, while it has eigenvalues over C: indeed, the
% characteristic
% polynomial t2 + 1 has no real roots and two complex roots. The reader should
% observe that, as a linear transformation of the real plane R2 , this matrix
% corresponds
% to a 90\circ  counterclockwise rotation; the reason why this transformation has
% no (real)
% eigenvalues is that no direction in the plane is preserved through a 90\circ
% rotation. ?
% Example 6.16. An example with a different flavor is given by the matrix
% 0 2
% .
% 1 0
% This has characteristic polynomial t2 - 2; hence it has eigenvalues over R, but
% not over Q. Geometrically, the corresponding transformation flips the plane
% about
% a line; but that line has irrational slope, so it contains no nonzero vectors
% with
% rational components.
%  ?
% As another benefit of Lemma 6.14, we may now introduce the following notion:
% Definition 6.17. The algebraic multiplicity of an eigenvalue of a linear
% transfor-
% mation \alpha  of a finitely generated free module is its multiplicity as a root
% of the
% characteristic polynomial of \alpha .
%  ?
% For example, the identity on F has the single eigenvalue 1, with (algebraic)
% multiplicity equal to the rank of F .
% The sum of the algebraic multiplicities is bounded by the degree of the charac-
% teristic polynomial, that is, the dimension of the space. In particular,
% Corollary 6.18. The number of eigenvalues of a linear transformation of Rn is
% at most n. If the base ring R is an algebraically closed field, then every
% linear
% transformation has exactly n eigenvalues (counted with algebraic multiplicity).
% Proof. Immediate from Lemmas 6.14, V.5.1, and V.5.10.
%  ?
% There is a different notion of multiplicity of an eigenvalue, related to how big
% the corresponding31 eigenspace may be.
% Definition 6.19. Let \lambda  be an eigenvalue of a linear transformation \alpha
% of a free
% R-module F . Then a nonzero v \in  F is an eigenvector for \alpha ,
% corresponding to the
% eigenvalue \lambda , if \alpha (v) = \lambda v, that is, if v \in  ker(\lambda I
% - \alpha ). The submodule ker(\lambda I - \alpha )
% is the eigenspace corresponding to \lambda .
%  ?
% Definition 6.20. The geometric multiplicity of an eigenvalue \lambda  is the
% rank of its
% eigenspace.
%  ?
% This is clearly invariant under similarity (Exercise 6.14). Geometric and alge-
% braic multiplicities do not necessarily coincide: for example, the matrix
% 1 1
% 0 1
% has characteristic polynomial (t - 1)2, and hence the
%  single eigenvalue 1, with
% ?v1?algebraic multiplicity 2. However, for a vector v = ,
% v21 1
%  v1
%  v1 + v2
% =
% 0 1
%  v2
%  v2
% equals 1v if and only if v2 = 0: that is, the eigenspace
%  ? ?
%  corresponding to the single
% eigenvalue 1 consists of the span of the vector 10 and has dimension 1. Thus,
% the
% geometric multiplicity of the eigenvalue is 1 in this example.
% Contemplating this example will convince the reader that the geometric multi-
% plicity of an eigenvalue is always less than its algebraic multiplicity. In the
% neatest
% possible situation, an operator \alpha  on a free module F of rank n may have
% all its n
% eigenvalues in the base ring R, and each algebraic multiplicity may agree with
% the
% corresponding geometric multiplicity. If this is the case, F may then be
% expressed as
% a direct sum of the eigenspaces, producing the so-called spectral decomposition
% of F
% determined by \alpha  (cf. Exercise 6.15 for a concrete instance of this
% situation). The
% action of \alpha  on F is then completely transparent, as it amounts to simply
% applying
% a (possibly) different scaling factor on each piece of the spectral
% decomposition.
% 31 If we were serious about pursuing the theory over arbitrary integral domains,
% I would feel
% compelled to call this object the eigenmodule of \alpha . However, we are
% eventually going to restrict
% our attention to the case in which the base ring is a field, so I will not
% bother.
% If A is a matrix representing \alpha  \in  EndR (V ) with respect to any basis,
% then \alpha
% admits a spectral decomposition if and only if A is similar to a diagonal
% matrix:
% (
%  )
% \lambda 1 0 . . . 0
% | 0 \lambda 2 . . . 0 |
% |
%  | -1
% A = P | ( .
%  ..
%  .
%  ..
%  ..
%  .
%  . .. )
%  | P ,
%  0 . . . \lambda n
% where \lambda 1 , . . . , \lambda n are the eigenvalues of \alpha . In this case
% we say that A (or \alpha ) is
% diagonalizable. A moment's thought reveals that the columns of P are then a
% basis
% of V consisting of eigenvectors of \alpha .
% Once more, I promise that the situation will become (even) clearer once we
% bring Theorem 5.6 into the picture.

\subsection*{Exercises}

% In these exercises, R denotes an integral domain.
% 6.1. Let k be an infinite field, and let n be any positive integer.
% \bullet  Prove that there are finitely many equivalence classes of matrices in Mn (k).
% \bullet  Prove that there are infinitely many similarity classes of matrices in Mn (k).
% 6.2. Let F be a free R-module of rank n, and let \alpha , \beta  \in  EndR (F ).
% Prove that
% det(\alpha  \circ  \beta ) = det(\alpha ) det(\beta ). Prove that \alpha  is
% invertible if and only if det(\alpha ) is
% invertible.
% 6.3. ? Prove that if \alpha  and \beta  are similar in the sense of Definition
% 6.3, then det(\alpha ) =
% det(\beta ) and tr(\alpha ) = tr(\beta ). (Do this without using Proposition
% 6.9!) [\S{}6.2]
% 6.4. ? Let F be a finitely generated free R-module, and let \alpha  be a linear
% trans-
% formation of F . Give an example of an injective \alpha  which is not
% surjective; in fact,
% prove that \alpha  is not surjective precisely when det(\alpha ) is not a unit.
% [\S{}6.2]
% 6.5. ? Let k be a field, and view k[t] as a vector space over k in the evident
% way.
% Give an example of a k-linear transformation k[t] \to  k[t] which is injective
% but not
% surjective; give an example of a linear transformation which is surjective but
% not
% injective. [\S{}6.2, \S{}VII.4.1]
% 6.6. Prove that two 2 \times  2 matrices have the same characteristic polynomial
% if and
% only if they have the same trace and determinant. Find two 3 \times  3 matrices
% with
% the same trace and determinant but different characteristic polynomials.
% 6.7. ? Let \alpha  \in  EndR (F ) be a linear transformation on a free R-module
% F . Prove
% that the set of polynomials f (t) \in  R[t] such that f (\alpha ) = 0 is an
% ideal of R[t]. [\S{}6.2]
% 6.8. ? Let A \in  Mn (R) be a square matrix, and let At be its transpose. Prove
% that
% A and At have the same characteristic polynomial and the same annihilator
% ideals.
% [\S{}7.2]
% 6.9. ? Prove the Cayley-Hamilton theorem, as follows. Recall that every square
% matrix M has an adjoint matrix, which I will denote adj(M ), and that we proved
% (Corollary 3.5) that adj(M ) \cdot  M = det(M ) \cdot  I. Applying this to M =
% tI - A (with
% A a matrix realization of \alpha  \in  EndR(F )) gives
% (*)
%  adj(tI - A) \cdot  (tI - A) = P\alpha  (t) \cdot  I.
% ?
% n-1
%  k
% Prove that there exist matrices Bk \in  Mn (R) such that adj(tI - A) = k=0 Bk t
% ;
% then use (*) to obtain P\alpha  (A) = 0, proving the Cayley-Hamilton theorem.
% [\S{}6.2]
% 6.10. ? Let F1 , F2 be free R-modules of finite rank, and let \alpha 1 , resp.,
% \alpha 2 , be linear
% transformations of F1 , resp., F2 . Let F = F1 \oplus  F2 , and let \alpha  =
% \alpha 1 \oplus  \alpha 2 be the
% linear transformation of F restricting to \alpha 1 on F1 and \alpha 2 on F2 .
% \bullet  Prove that P\alpha  (t) = P\alpha 1 (t)P\alpha 2 (t). That is, the characteristic polynomial is
% multiplicative under direct sums.
% \bullet  Find an example showing that the minimal polynomial is not multiplicative
% under direct sums.
% [6.11, \S{}7.2]
% 6.11. \lnot  Let \alpha  be a linear transformation of a finite-dimensional
% vector space V ,
% and let V1 be an invariant subspace, that is, such that \alpha (V1 ) \subseteq
% V1 . Let \alpha 1 be
% the restriction of \alpha  to V1, and let V2 = V /V1 . Prove that \alpha
% induces a linear
% transformation \alpha 2 on V2 , and (in the same vein as Exercise 6.10) show
% that P\alpha  (t) =
% P \alpha 1 (t)P\alpha 2 (t). Also, prove that tr(\alpha r ) = tr(\alpha 1 r ) +
% tr(\alpha 2 r ), for all r \geq  0. [6.12]
% 6.12. Let \alpha  be a linear transformation of a finite-dimensional C-vector
% space V .
% Prove the identity of formal power series with coefficients in C:
% ? ?
%  ?
%  \infty 
%  r t
% r
% = exp
%  tr(\alpha  )
%  .
% det(1 - \alpha t)
%  r
% r=1
% (Hint: The left-hand side is essentially the inverse of the characteristic
% polynomial
% of \alpha . Use Exercise 6.11 to show that both sides are multiplicative with
% respect
% to exact sequences 0 \to  V1 \to  V \to  V2 \to  0, where V1 is an invariant
% subspace.
% Use this and the fact that \alpha  admits nontrivial invariant subspaces since C
% is alge-
% braically closed (why?) to reduce to the case dim V = 1, for which the identity
% is
% an elementary calculus exercise.)
% With due care, the identity can be stated and proved (in the same way) over
% arbitrary fields of characteristic 0. It is an ingredient in the cohomological
% inter-
% pretation of the `Weil conjectures'.
% 6.13. Let A be a square matrix with integer entries. Prove that if \lambda  is a
% rational
% eigenvalue of A, then in fact \lambda  \in  Z. (Hint: Proposition V.5.5.)
% 6.14. ? Let \lambda  be an eigenvalue of two similar transformations \alpha ,
% \beta . Prove that the
% geometric multiplicities of \lambda  with respect to \alpha  and \beta
% coincide. [\S{}6.3]
% 6.15. ? Let \alpha  be a linear transformation on a free R-module F , and let v
% 1 , . . . , vn
% be eigenvectors corresponding to pairwise distinct eigenvalues \lambda 1 , . . .
% , \lambda n . Prove
% that v1, . . . , vn are linearly independent. (Hint: If not, there is a shortest
% linear
% combination r1vi1 + \dots  + rm vim = 0 with all rj \in  R, rj = 0. Compare the
% action
% of \alpha  on this linear combination with the product by \lambda i1 .)
% It follows that if F has rank n and \alpha  has n distinct eigenvalues, then
% \alpha  induces
% a spectral decomposition of F . [\S{}6.3, VII.6.14]
% 6.16. \lnot  The standard inner product on V = Rn is the map V \times  V \to  R
% defined by
% (v, w) := vt \cdot  w
% (viewing elements v \in  V as column vectors). The standard hermitian product
% on W = Cn is the map W \times  W \to  C defined by
% (v, w) := v \dagger  \cdot  w,
% \dagger where for any matrix M , M stands for the matrix obtained by taking the complex
% conjugates of the entries of the transpose M t .
% These products satisfy evident linearity properties: for example, for \lambda
% \in  C and
% v, w \in  W
% (\lambda v, w) = \lambda (v, w).
% Prove32 that a matrix M \in  Mn (R) belongs to On (R) if and only if it
% preserves
% the standard inner product on Rn :
% (\forall v, w \in  Rn)
%  (M v, M w) = (v, w).
% Likewise, prove that a matrix M \in  Mn (C) belongs to U(n) if and only if it
% preserves
% the standard hermitian product on Cn . [6.18]
% 6.17. \lnot  We say that two vectors v, w of Rn or Cn are orthogonal if (v, w) =
% 0.
% The orthogonal complement v\bot  of v is the set of vectors w that are
% orthogonal to v.
% Prove that if v = 0 in V = Rn or Cn, then v \bot  is a subspace of V of
% dimension n-1.
% [7.16, VIII.5.15]
% 6.18. \lnot  Let V = Rn, endowed with the standard inner product defined in
% Exer-
% cise 6.16. A set of distinct vectors v1 , . . . , vr is orthonormal if (vi , vj
% ) = 0 for i = j
% and 1 for i = j. Geometrically, this means that each vector has length 1, and
% dif-
% ferent vectors are orthogonal. The same terminology may be used in Cn, w.r.t.
% the
% standard hermitian product.
% Prove that M \in  On (R) if and only if the columns of M are orthonormal, if and
% only if the rows of M are orthonormal33 .
% Formulate and prove an analogous statement for U(n). (The group U(n) is
% called the unitary group. Note that, for n = 1, it consists of the complex
% numbers
% of norm 1.) [6.19, 7.16, VIII.5.9]
% 6.19. \lnot  Let v1 , . . . , vr form an orthonormal set of vectors in Rn .
% 32 See Exercise II.6.1 for the definitions of On
%  (R) and U(n). I trust that basic facts on inner
% products are not new to the reader. Among these, recall that for v, w \in  Rn,
% (v, v) equals the
% square of the length of v, and (v, w) equals the product of the lengths of v and
% w and the cosine
% of the angle formed by v and w. Thus, the reader is essentially asked to verify
% that M \in  On(R) if
% and only if M preserves lengths and angles, and to check an analogous statement
% in the complex
% environment.
% 33 The group On
%  (R) is called the orthogonal group; orthonormal group would probably be a
% more appropriate terminology.

\section{Canonical forms}
\label{sec:linalg:canonical-forms}

% \bullet  Prove that v1 , . . . , vr are linearly independent; so they form an orthonormal
% basis of the space V they span.
% \bullet  Let w = a1v1 + \dots  + ar vr be a vector of V . Prove that ai = (vi , w).
% \bullet  More generally, prove that if w \in  Rn , then (vi , w) is the component of vi in
% the orthogonal projection wV of w onto V . (That is, prove that w - wV is
% orthogonal to all vectors of V .)
% w
% wV
%  V
% For reasons such as these, it is convenient to work with orthonormal bases. Note
% that, by Exercise 6.18, a matrix is in On (R) if and only if its columns form an
% orthonormal basis of R n. Again, the reader should formulate parallel statements
% for hermitian products in Cn . [6.20]

\subsection{The Gram-Schmidt process takes as input a choice of linearly independent}
\label{subsec:linalg:the-gram-schmidt-process-takes-as-input-a-choice-of-linearly-independent}

% vectors v1 , . . . , vr of Rn and returns vectors w1 , . . . , wr spanning the
% same space V
% spanned by v1 , . . . , vr and such that (wi, wj ) = 0 for i = j. Scaling each
% wi by its
% length, this yields an orthonormal basis for V .
% Here is how the Gram-Schmidt process works:
% --- w1 := v1 .
% --- For k \geq  1, wk := vk - orthogonal projection of v k onto ?w1 , . . . ,
% wk-1 ?.
% (Cf. Exercise 6.19.) Prove that this process accomplishes the stated goal.

%\lnot  A matrix M \in  Mn(Rn) is symmetric if M t = M}
% Prove that M is

% symmetric if and only if (\forall v, w \in  Cn), (M v, w) = (v, M w).
% A matrix M \in  Mn (Cn ) is hermitian if M \dagger  = M . Prove that M is
% hermitian if
% and only if (\forall v, w \in  Cn ), (M v, w) = (v, M w).
% In both cases, one may say that M is self-adjoint; this means that shuttling it
% from one side of the product to the other does not change the result of the
% operation.
% A hermitian matrix with real entries is symmetric. It is in fact useful to think
% of real symmetric matrices as particular cases of hermitian matrices. [6.22]

%Prove that the eigenvalues of a hermitian matrix (Exercise 6.21) are real}

% Also, prove that if v, w are eigenvectors of a hermitian matrix, corresponding
% to
% different eigenvalues, then (v, w) = 0. (Thus, eigenvectors with distinct
% eigenvalues
% for a real symmetric matrix are orthogonal.) [7.20]

\subsection{Linear transformations of free modules; actions of polynomial rings}
\label{subsec:linalg:linear-transformations-of-free-modules-actions-of-polynomial-rings}

% I have promised to use Theorem 5.6 to better understand the similarity relation
% and to obtain special forms for square matrices with entries in a field. I am
% now
% ready to make good on my promise.
% The questions the reader should be anxiously asking are, where is the PID?
% Where is the finitely generated module? Why would a classification theorem for
% these things have anything to tell us about matrices? Once these questions are
% answered, everything else will follow easily. In a sense, this is the one issue
% that I
% keep in my mind concerning all these considerations: once I remember how to get
% an interesting module out of a linear transformation of vector spaces, the rest
% is
% essentially an immediate consequence of standard notions.
% Once more, while the main application will be to vector spaces and matrices
% over a field, we do not need to preoccupy ourselves with specializing the
% situation
% before we get to the key point. Therefore, let's keep working for a while longer
% on
% our given integral domain 34 R.
% Claim 7.1. Giving a linear transformation on a free R-module F is the same as
% giving an R[t]-module structure on F , compatible with its R-module structure.
% Here R[t] is the polynomial ring in one indeterminate t over the ring R. The
% claim is much less impressive than it sounds; in fact, it is completely
% tautological.
% Giving an R[t]-module structure on F compatible with the R-module structure is
% the same as giving a homomorphism of R-algebras
% \phi  : R[t] \to  EndR (F )
% (this is an insignificant upgrade of the very definition of module from
% \S{}III.5.1).
% By the universal property satisfied by polynomial rings (\S{}III.2.2, especially
% Exam-
% ple III.2.3), giving \phi  as an extension of the basic R-algebra structure of
% EndR (F ) is
% just the same as specifying the image \phi (t), that is, choosing an element of
% EndR (F ).
% This is precisely what the claim says.
% My propensity for the use of a certain language may obfuscate the matter,
% which is very simple. Given a linear transformation \alpha  of a free module F ,
% we can
% define the action of a polynomial
% f (t) = rm tm + rm-1 tm-1 + \dots  + r0 \in  R[t]
% on F as follows: for every v \in  F , set
% f (t)(v) := rm \alpha m (v) + rm-1 \alpha  m-1 (v) + \dots  + r0 v,
% where \alpha k denotes the k-fold composition of \alpha  with itself; we have
% already run into
% this in \S{}6.2. Conversely, an action of R[t] on F determines in particular a
% map
% \alpha  : F \to  F , that is, `multiplication by t':
% v ?\to  tv;
% the compatibility with the R-module structure on F tells us precisely that
% \alpha  is a
% linear transformation of F . Again, this is the content of Claim 7.1.
% Tautological as it is, Claim 7.1 packs a good amount of information. Indeed,
% the R[t] module knows everything about similarity:
% 34 In fact, the requirements that R be an integral domain and that F be free are
% immaterial
% here; cf. Exercise III.5.11.
% Lemma 7.2. Let \alpha , \beta  be linear transformations of a free R-module F .
% Then the
% corresponding R[t]-module structures on F are isomorphic if and only if \alpha
% and \beta
% are similar.
% Proof. Denote by F\alpha , F\beta  the two R[t]-modules defined on F by \alpha ,
% \beta  as per
% Claim 7.1.
% Assume first that \alpha  and \beta  are similar. Then there exists an
% invertible R-linear
% transformation \pi  : F \to  F such that
% \beta  = \pi  \circ  \alpha  \circ  \pi -1 ;
% that is, \pi  \circ  \alpha  = \beta  \circ  \pi . We can view \pi  as an
% R-linear map
% F\alpha  \to  F\beta  .
% I claim that it is R[t]-linear: indeed, multiplication by t is \alpha  in
% F\alpha  and \beta  in F\beta  , so
% \pi (tv) = \pi  \circ  \alpha (v) = \beta  \circ  \pi (v) = t\pi (v).
% Thus \pi  is an invertible R[t]-linear map F\alpha  \to  F\beta  , proving that
% F\alpha  and F\beta  are
% isomorphic as R[t]-modules.
% The converse implication is obtained essentially by running this argument in
% reverse and is left to the reader (Exercise 7.1).
%  ?
% Corollary 7.3. There is a one-to-one correspondence between the similarity
% classes
% of R-linear transformations of a free R-module F and the isomorphism classes of
% R[t]-module structures on F .
% Of course the same statement holds for square matrices with entries in R, in
% the finite rank case.
% Corollary 7.3 is the tool we were looking for, bridging between classifying sim-
% ilarity classes of linear transformations or matrices and classifying modules.
% The
% task now becomes that of translating the notions we have encountered in \S{}6
% into
% the module language and seeing if this teaches us anything new.
% In any case, since we have classified finitely generated modules over PIDs,
% it is clear that we will be able to classify similarity classes of
% transformations of
% finite-rank free modules---provided that R[t] is a PID. This is where the
% additional
% hypothesis that R be a field enters the discussion (cf. Exercise V.2.12).

\subsection{k[t]-modules and the rational canonical form}
\label{subsec:linalg:k-t-modules-and-the-rational-canonical-form}
% It is finally time to spe-

% cialize to the case in which R = k is a field and F has finite rank; so F = V is
% simply a finite-dimensional vector space. Let n = dim V .
% By the preceding discussion, choosing a linear transformation of V is the same
% as giving V a k[t]-module structure (compatible with its vector space
% structure);
% similar linear transformations correspond to isomorphic k[t]-modules. Then V is
% a finitely generated module over k[t], and k[t] is a PID since k is a field
% (Exer-
% cise III.4.4 and \S{}V.2); therefore, we are precisely in the situation covered
% by the
% classification theorem.
% Even before spelling out the result of applying Theorem 5.6, we can make use
% of its slogan version: every finitely generated module over a PID is a direct
% sum
% of cyclic modules. This tells us that we can understand all linear
% transformations
% of finite-dimensional vector spaces, up to similarity, if we understand the
% linear
% transformation given by multiplication by t on a cyclic k[t]-module
% k[t]
% V =
%  ,
% (f (t))
% where f (t) is a nonconstant monic polynomial:
% f (t) = tn + rn-1 tn-1 + \dots  + r0 .
% It is worthwhile obtaining a good matrix representation of this poster case. We
% choose the basis
% 1, t, \dots  , tn-1
% of V (cf. Proposition III.4.6). Recall that the columns of the matrix
% corresponding
% to a transformation consist of the images of the chosen basis (cf. the comments
% preceding Corollary 2.2). Since multiplication by t on V acts as
% \{
% |
%  |
%  |
%  |
%  |
%  1 t ?\to t,
%  ?\to t2 ,
% |
%  |
%  |
%  |
%  \} \cdot \cdot \cdot 
% tn-1
%  ?\to tn = -rn-1 tn-1 - \dots  - r0 ,
% the matrix corresponding to this linear transformation is
% (
%  )
% 0 0 0 ... 0
%  -r0
% |1 |
%  0 0 . . . 0
%  -r1 |
%  |
% |0 |
%  1 0 . . . 0
%  -r2 |
%  |
% | | . .. . .. .. .
%  . .
%  .
%  ..
%  .
%  .. . |
%  | .
% |
%  |
% (0 0 0 . . . 0 -rn-2 )
% 0 0 0 . . . 1 -rn-1
% Definition 7.4. This is called the companion matrix of the polynomial f (t), de-
% noted Cf (t).
%  ?
% Theorem 5.6 tells us (among other things) that every linear transformation
% admits a matrix representation into blocks, each of which is the companion
% matrix
% of a polynomial. Here is the statement of Theorem 5.6 in this context:
% Theorem 7.5. Let k be a field, and let V be a finite-dimensional vector space.
% Let
% \alpha  be a linear transformation on V , and endow V with the corresponding k[t]-module
% structure, as in Claim 7.1. Then the following hold:
% \bullet  There exist distinct monic irreducible polynomials p1 (t), . . . , p s(t) \in  k[t] and
% positive integers rij such that
% !
%  k[t]
% V \sim 
%  =
% (p
% i
%  (t)rij )
% i,j
% as k[t]-modules.
% \bullet  There exist monic nonconstant polynomials f1 (t), . . . , fm (t) \in  k[t] such that
% f1 (t) | \dots  | fm (t) and
% V \sim 
%  =
%  k[t]
%  \oplus  \cdot \cdot \cdot  \oplus 
%  k[t]
% (f1 (t))
%  (fm (t))
% as k[t]-modules.
% Via these isomorphisms, the action of \alpha  on V corresponds to multiplication
% by t.
% Further, two linear transformations \alpha , \beta  are similar if and only if
% they have
% the same collections of invariants pi (t)rij (`elementary divisors'), fi (t)
% (`invariant
% factors').
% Proof. Since dim V is finite, V is finitely generated as a k-module and a
% fortiori as
% a k[t]-module. The two isomorphisms are then obtained by applying Theorem 5.6.
% All the relevant polynomials may be chosen to be monic since every polynomial
% over a field is the associate of a (unique) monic polynomial (Exercise III.4.7).
% The
% fact that the action of \alpha  corresponds to multiplication by t is precisely
% what de-
% fines the corresponding k[t]-module structure on V . The statement about similar
% transformations follows from Corollary 7.3.
%  ?
% Theorem 7.5 answers (over fields) the question raised in \S{}6.1: we now have a
% list of invariants which describe completely the similarity class of a linear
% transfor-
% mation. Further, these invariants provide us with a special matrix
% representation
% of a given linear transformation.
% Definition 7.6. The rational canonical form of a linear transformation \alpha
% of a
% vector space V is the block matrix
% (
%  )
% Cf1(t)
% |
%  |
% |
%  ..
%  .
%  | ,
% (
%  )
% Cfm(t)
% where f1 (t), . . . , fm (t) are the invariant factors of \alpha .
%  ?
% The rational canonical form of a square matrix is (of course) the rational
% canon-
% ical form of the corresponding linear transformation. The following statement is
% an
% immediate consequence of Theorem 7.5.
% Corollary 7.7. Every linear transformation admits a rational canonical form. Two
% linear transformations have the same rational canonical form if and only if they
% are
% similar.
% Remark 7.8. The `rational' in rational canonical form has nothing to do with Q;
% it is meant to remind the reader that this form can be found without leaving
% the base field. The other canonical form we will encounter will have entries in
% a possibly larger field, where the characteristic polynomial of the
% transformation
% factors completely.
%  ?
% Corollary 7.7 fulfills explicitly another wish expressed at the end of \S{}6.1,
% that
% is, to find one distinguished representative in each similarity class of
% matrices/linear
% transformations. The rational canonical form is such a representative.
% The next obvious question is how the invariants we just found relate to our
% more naive attempts to produce invariants of similar transformations in \S{}6.
% Proposition 7.9. Let f1 (t) | \dots  | fm (t) be the invariant factors of a
% linear trans-
% formation \alpha  on a vector space V . Then the minimal polynomial m\alpha  (t)
% equals fm (t),
% and the characteristic polynomial P\alpha  (t) equals the product f1 (t) \dots
% fm (t).
% Proof. Tracing definitions, (m\alpha (t)) is the annihilator ideal of V when
% this is viewed
% as a k[t]-module via \alpha  (as in Claim 7.1). Therefore the equality of
% m\alpha  (t) and fm (t)
% is a restatement of Lemma 5.7.
% Concerning the characteristic polynomial, by Exercise 6.10 it suffices to prove
% the statement in the cyclic case: that is, it suffices to prove that if f (t) is
% a monic
% polynomial, then f (t) equals the characteristic polynomial of the companion
% matrix
% Cf (t) . Explicitly, let
% f (t) = tn + rn-1 tn-1 + \dots  + r0
% be a monic polynomial; then the task amounts to showing that
% (
%  )
% t
%  0 0 ...
%  r0
% |-1 t 0 . . .
%  r1 |
% |
%  |
% | |
%  0 -1 t . . .
%  r2 |
%  |
% det | | .
%  ..
%  .
%  ..
%  .
%  .. . .
%  .
%  .
%  ..
%  ..
%  .
%  |
%  | = f (t).
% |
%  |
% ( 0
%  0 0 ...
%  t
%  rn-2 )
%  0 0 . . . -1 t + rn-1
% This can be done by an easy induction and is left to the reader (Exercise 7.2).
% ?
% Corollary 7.10 (Cayley-Hamilton). The minimal polynomial of a linear transfor-
% mation divides its characteristic polynomial.
% Proof. This has now become evident, as promised in \S{}6.2.
%  ?
% Finding the rational canonical form of a given matrix A amounts to working
% out the classification theorem for the specific k[t]-module corresponding to A
% as in
% Claim 7.1. As k[t] is in fact a Euclidean domain, we know that this can be done
% by
% Gaussian elimination (cf. \S{}2.4). In practice, the process consists of
% compiling the
% information obtained while diagonalizing tI - A over k[t] by elementary
% operations.
% The reader is invited to either produce an original algorithm or at least look
% it up
% in a standard reference to see what is involved.
% Major shortcuts may come to our aid. For example, if the minimal and char-
% acteristic polynomials coincide, then one knows a priori that the corresponding
% module is cyclic (cf. Remark 5.8), and it follows that the rational canonical
% form is
% simply the companion matrix of the characteristic polynomial.
% In any case, the strength of a canonical form rests in the fact that it allows
% us
% to reduce general facts to specific, standard cases. For example, the following
% state-
% ment is somewhat mysterious if all one knows are the definitions, but it becomes
% essentially trivial with a pinch of rational canonical forms:
% Proposition 7.11. Let A \in  M n(k) be a square matrix. Then A is similar to its
% transpose.
% Proof. If B is similar to A and we can prove that B is similar to its transpose
% Bt ,
% then A is similar to its transpose At : because B = P AP -1 , B t = QBQ-1 give
% At = (P t QP )A(P tQP )-1 .
% Therefore, it suffices to prove the statement for matrices in rational canonical
% form.
% Further, to prove the statement for a block matrix, it clearly suffices to prove
% it
% for each block; so we may assume that A is the companion matrix C of a
% polynomial
% f (t). Since the characteristic and minimal polynomials of the transpose C t
% coincide
% with those of C (Exercise 6.8), they are both equal to f (t). It follows that
% the
% rational canonical form of C t is again the companion matrix to f (t); therefore
% C t
% and C are similar, as needed.
%  ?

\subsection{Jordan canonical form}
\label{subsec:linalg:jordan-canonical-form}
% We have derived the rational canonical form from

% the invariant factors of the module corresponding to a linear transformation. A
% useful alternative can be obtained from the elementary divisors, at least in the
% case
% in which the characteristic polynomial factors completely over the field k. If
% this is
% not the case, one can enlarge k so as to include all the roots of the
% characteristic
% polynomial: the reader proved this in Exercise V.5.13. The price to pay will be
% that
% the Jordan canonical form of a linear transformation \alpha  \in  Endk (V ) may
% be a matrix
% with entries in a field larger than k. In any case, whether two transformations
% are
% similar or not is independent of the base field (Exercise 7.4), so this does not
% affect
% the issue at hand.
% Given \alpha  \in  End k (V ), obtain the elementary divisor decomposition of
% the corre-
% sponding k[t]-module, as in Theorem 7.5:
% ! k[t]
% V \sim 
%  =
%  .
% (pi (t)rij )
% i,j
% It is an immediate consequence of Proposition 7.9 and the equivalence between
% the elementary divisor and invariant factor formulations that the characteristic
% polynomial P\alpha  (t) equals the product
% ?
% pi (t)rij .
% i,j
% Lemma 7.12. Assume that the characteristic polynomial P\alpha  (t) factors
% completely;
% that is,
% ?
%  s
% P\alpha  (t) =
%  (t - \lambda i )mi
% i=1
% where \lambda i , i = 1, . . . , s, are the distinct eigenvalues of \alpha  (and
%  ? mi are their algebraic
% multiplicities; cf. \S{}6.3). Then pi(t) = (t - \lambda i ), and mi = j rij .
% In this situation, the minimal polynomial of \alpha  equals
% ?
%  s
% m\alpha  (t) =
%  (t - \lambda i )maxj {rij } .
% i=1
% Proof. The first statement follows from uniqueness of factorizations. The state-
% ment about the minimal polynomial is immediate from Proposition 7.9 and the
% bookkeeping giving the equivalence of the two formulations in Theorem 7.5.
%  ?
% The elementary divisor decomposition splits V into a different collection of
% cyclic modules than the decomposition in invariant factors: the basic cyclic
% bricks
% are now of the form k[t]/(p(t)r ) for a monic prime p(t); by Lemma 7.12,
% assuming
% that the characteristic polynomial factors completely over k, they are in fact
% of the
% form
% k[t]
% ((t - \lambda )r )
% for some \lambda  \in  k (which equals an eigenvalue of \alpha ) and r > 0. Just
% as we did for
% `companion matrices', we now look for a basis with respect to which \alpha  (=
% multipli-
% cation by t; keep in mind the fine print in Theorem 7.5) has a particularly
% simple
% matrix representation. This time we choose the basis
% (t - \lambda )r-1 ,
%  (t - \lambda )r-2 ,
%  \cdot \cdot \cdot  ,
%  (t - \lambda )0 = 1.
% Multiplication by t on V acts (in the first line, use the fact that (t-\lambda
% )r = 0 in V ) as
% \{
% | |
%  |
%  |
%  |
%  (t (t - \lambda )r-1 \lambda )r-2 ?\to t(t ?\to t(t - - \lambda )r-1 \lambda
%  )r-2 = = \lambda (t (t - - \lambda )r-1 \lambda )r-1 + + \lambda (t (t - -
%  \lambda )r-2,
%  \lambda )r = \lambda (t - \lambda )r-1 ,
% -|
%  |
%  |
%  |
%  \}
%  \cdot \cdot \cdot 
% 1 ?\to t = (t - \lambda ) + \lambda .
% Therefore, with respect to this basis the linear transformation has matrix
% (
%  )
% \lambda  1 0 ... 0 0
% | 0 \lambda  1 . . . 0 0|
% |
%  |
% | 0 0 \lambda  . . . 0 0|
% |
%  |
% | | .. . . .. .
%  .. . .
%  .
%  .. .
%  .. .
%  |
%  | .
% |
%  |
% ( 0 0 0 . . . \lambda  1)
% 0 0 0 ... 0 \lambda 
% Definition 7.13. This matrix is the Jordan block of size r corresponding to
% \lambda ,
% denoted J\lambda ,r .
%  ?
% We can put several blocks together for a given \lambda : for (rj ) = (r1 , . . .
% , r?), let
% (
%  )
% J\lambda ,r1
% |
%  |
% J\lambda ,(rj ) := |
%  (
%  ..
%  .
%  | )
%  .
% J\lambda ,r?
% With this notation in hand, we get new distinguished representatives of the
% similarity class of a given linear transformation:
% Definition 7.14. The Jordan canonical form of a linear transformation \alpha  of
% a
% vector space V is the block matrix
% (
%  )
% J\lambda 1 ,(r1j )
% |
%  |
% |
%  (
%  ..
%  .
%  | )
%  ,
% J\lambda s ,(rsj )
% where (t - \lambda i )rij are the elementary divisors of \alpha  (cf. Lemma
% 7.12).
%  ?
% This canonical form is `a little less canonical' than the rational canonical
% form,
% in the sense that while the rational canonical form is really unique, some
% ambiguity
% is left in the Jordan canonical form: there is no special way to choose the
% order
% of the different blocks any more than there is a way to order35 the eigenvalues
% \lambda 1 , . . . , \lambda s.
% Therefore, the analog of Corollary 7.7 should state that two linear transforma-
% tions are similar if and only if they have the same Jordan canonical form up to
% a
% reordering of the blocks. As we have seen, every linear transformation \alpha
% \in  Endk (V )
% admits a Jordan canonical form if P\alpha  (t) factors completely over k; for
% example, we
% can always find a Jordan canonical form with entries in the algebraic closure of
% k.
% Example 7.15. One use of the Jordan canonical form is the enumeration of all
% possible similarity classes of transformations with given eigenvalues. For
% example,
% there are 5 similarity classes of linear transformations with a single
% eigenvalue \lambda
% with algebraic multiplicity 4, over a 4-dimensional vector space: indeed, there
% are 5
% different ways to stack together Jordan blocks corresponding to the same
% eigenvalue,
% within a 4 \times  4 square matrix:
% (
%  )
%  (
%  )
%  (
%  )
% \lambda 
%  \lambda 
%  \lambda 
% |
%  (
%  | | |
%  \lambda 
%  \lambda 
%  0 0 |
%  |
%  | )
%  |
%  ,
%  |
%  ( |
%  | 0
%  \lambda 
%  \lambda 
%  0 0 )
%  |
%  | |
%  ,
%  | ( |
%  |
%  \lambda 
%  \lambda 
%  1 0 )
%  | |
%  |
%  ,
%  \lambda 
%  0 \lambda 
%  \lambda 
% (
%  )
%  (
%  )
% \lambda 
%  \lambda 
% | |
%  |
%  |
% |
%  \lambda 
%  0 |
%  | 0
%  \lambda 
%  0 |
% |
%  | , |
%  | .
%  ?
% ( 0
%  \lambda 
%  0 )
%  ( 0
%  \lambda 
%  1 )
%  \lambda 
%  \lambda 
% The Jordan canonical form clarifies the difference between algebraic and geo-
% metric multiplicities of an eigenvalue. The algebraic multiplicity of \lambda
% as an eigen-
% value of a linear transformation \alpha  is (of course) the number of times
% \lambda  appears in
% the Jordan canonical form of \alpha . Recall that the geometric multiplicity of
% \lambda  equals
% the dimension of the eigenspace of \lambda .
% Proposition 7.16. The geometric multiplicity of \lambda  as an eigenvalue of
% \alpha  equals
% the number of Jordan blocks corresponding to \lambda  in the Jordan canonical
% form of \alpha .
% 35 Of course if the ground field is Q or R, then we can order the \lambda i
%  's in, for example, increasing
% order. However, this is not an option on fields like C.
% Proof. As the geometric multiplicity is clearly additive in direct sums, it
% suffices
% to show that the geometric multiplicity of \lambda  for the transformation
% corresponding
% to a single Jordan block
% (
%  )
% \lambda  1 ... 0 0
% |0 \lambda  . . . 0 0|
% |
%  |
% J = |
%  | ... ... . . . ... ... |
%  |
% |
%  |
% (0 0 . . . \lambda  1)
% 0 0 ... 0 \lambda 
% is 1.
%  ( )
% v1
% Let v = | ( .. . |
%  ) be an eigenvector corresponding to \lambda . Then
% vr
% ( )
%  ( ) (
%  )
% v1
%  v1
%  \lambda v1 + v2
% |v2 |
%  |v2 | |\lambda v2 + v3 |
% | |
%  | | |
%  |
% \lambda  ( | .. . | )
%  = J | ( . .. ) | = (
%  |
%  ..
%  .
%  )
%  | ,
% vr
%  vr
%  \lambda vr
% yielding
% v2 = \dots  = vr = 0.
% That is, the eigenspace of \lambda  is generated by
% ( )
% |0|
% | |
% v = | . |
% ( .. )
% and has dimension 1, as needed.
%  ?

\subsection{Diagonalizability}
\label{subsec:linalg:diagonalizability}
% A linear transformation of a vector space V is diagonal-

% izable if it can be represented by a diagonal matrix.
% The question of whether a given linear transformation is or is not
% diagonalizable
% is interesting and important: as I pointed out already in \S{}6.3, when \alpha
% \in  Endk (V ) is
% diagonalizable, then the vector space V admits a corresponding spectral
% decomposi-
% tion: \alpha  is diagonalizable if and only if V admits a basis of eigenvectors
% of \alpha . Much
% more could be said on this important theme, but we are already in the position
% of
% giving useful criteria for diagonalizability.
% For example, the diagonal matrix representing \alpha  will necessarily be a
% Jordan
% canonical form for \alpha , consisting of blocks of size 1. Proposition 7.16
% tells us that
% this can be detected by the difference between algebraic and geometric
% multiplicity,
% so we get the following characterization of diagonalizable transformations:
% Corollary 7.17. Assume the characteristic polynomial of \alpha  \in  Endk (V )
% factors
% completely over k. Then \alpha  is diagonalizable if and only if the geometric
% and alge-
% braic multiplicities of all its eigenvalues coincide.
% Another view of the same result is the following:

\subsection*{Exercises}

% Proposition 7.18. Assume the characteristic polynomial of \alpha  \in  Endk (V )
% factors
% completely over k. Then \alpha  is diagonalizable if and only if the minimal
% polynomial
% of \alpha  has no multiple roots.
% Proof. Again, diagonalizability is equivalent to having all Jordan blocks of
% size 1
% in the Jordan canonical form of \alpha . Therefore, if the characteristic
% polynomial of \alpha
% factors completely, then \alpha  is diagonalizable if and only if all exponents
% rij appearing
% in Theorem 7.5 equal 1. By the second part of Lemma 7.12, this is equivalent to
% the requirement that the minimal polynomial of \alpha  have no multiple roots.
%  ?
% The condition of factorizability in these statements is necessary in order to
% guarantee that a Jordan canonical form exists. Not surprisingly, whether a
% matrix
% is diagonalizable or not depends on the ground field. For example,
% 0 -1
% 1 0
% is not diagonalizable over R (cf. Example 6.15), while it is diagonalizable over
% C.
% Endowing the vector space V with additional structure (such as an inner prod-
% uct) singles out important classes of operators, for which one can obtain
% precise and
% delicate results on the existence of spectral decompositions. For example, this
% can
% be used to prove that real symmetric matrices are diagonalizable (Exercise
% 7.20.)
% The reader will get a taste of these `spectral theorems' in the exercises.
% As above, k denotes a field.
% 7.1. ? Complete the proof of Lemma 7.2. [\S{}7.1]
% 7.2. ? Prove that the characteristic polynomial of the companion matrix of a
% monic
% polynomial f (t) equals f (t). [\S{}7.2]
% 7.3. ? Prove that two linear transformations of a vector space of dimension \leq
% 3 are
% similar if and only if they have the same characteristic and minimal
% polynomials.
% Is this true in dimension 4? [\S{}6.2]
% 7.4. ? Let k be a field, and let K be a field containing k. Two square matrices
% A, B \in  Mn (k) may be viewed as matrices with entries in the larger field K.
% Prove
% that A and B are similar over k if and only if they are similar over K.
% [\S{}7.3]
% 7.5. Find the rational canonical form of a diagonal matrix
% (
%  )
% \lambda 1 0 . . . 0
% | 0 \lambda 2 . . . 0 |
% |
%  |
% | ( ..
%  .
%  ..
%  .
%  ..
%  .
%  .. . | )
%  ,
%  0 . . . \lambda r
% assuming that the \lambda i are all distinct.
% 7.6. Let
%  (
%  )
% 6 -10
%  -10
% A = ( 3
%  -5
%  -6 ) .
% -1
% \bullet  Compute the characteristic polynomial of A.
% \bullet  Find the minimal polynomial of A (use the Cayley-Hamilton theorem!).
% \bullet  Find the invariant factors of A.
% \bullet  Find the rational canonical form of A.
% 7.7. Let V be a k-vector space of dimension n, and let \alpha  \in  Endk (V ).
% Prove that
% the minimal and characteristic polynomials of \alpha  coincide if and only if
% there is a
% vector v \in  V such that
% v,
%  \alpha (v),
%  \cdot \cdot \cdot  ,
%  \alpha n-1 (v)
% is a basis of V .
% 7.8. Let V be a k-vector space of dimension n, and let \alpha  \in  Endk (V ).
% Prove that the
% characteristic polynomial P\alpha  (t) divides a power of the minimal polynomial
% m\alpha  (t).
% 7.9. What is the number of distinct similarity classes of linear transformations
% on
% an n-dimensional vector space with one fixed eigenvalue \lambda  with algebraic
% multiplic-
% ity n?
% 7.10. Classify all square matrices A \in  Mn (k) such that A2 = A, up to
% similarity.
% Describe the action of such matrices `geometrically'.
% 7.11. A square matrix A \in  Mn (k) is nilpotent (cf. Exercise V.4.19) if Ak = 0
% for
% some integer k.
% \bullet  Characterize nilpotent matrices in terms of their Jordan canonical form.
% \bullet  Prove that if A k = 0 for some integer k, then Ak = 0 for some integer k no
% larger than n (= the size of the matrix).
% \bullet  Prove that the trace of a nilpotent matrix is 0.
% 7.12. \lnot  Let V be a finite-dimensional k-vector space, and let \alpha  \in
% Endk (V ) be a
% diagonalizable linear transformation. Assume that W \subseteq  V is an invariant
% subspace,
% so that \alpha  induces a linear transformation \alpha |W \in  Endk (W ). Prove
% that \alpha |W is also
% diagonalizable. (Use Proposition 7.18.) [7.14]
% 7.13. \lnot  Let R be an integral domain. Assume that A \in  Mn (R) is
% diagonalizable,
% with distinct eigenvalues. Let B \in  Mn (R) be such that AB = BA. Prove that
% B is also diagonalizable, and in fact it is diagonal w.r.t. a basis of
% eigenvectors
% -1 -1 -1of A. (If P is such that P AP is diagonal, note that P AP and P BP also
% commute.) [7.14]
% 7.14. Prove that `commuting transformations may be simultaneously diagonalized',
% in the following sense. Let V be a finite-dimensional vector space, and let
% \alpha , \beta  \in
% Endk (V ) be diagonalizable transformations. Assume that \alpha \beta  = \beta
% \alpha . Prove that V
% has a basis consisting of eigenvectors of both \alpha  and \beta . (Argue as in
% Exercise 7.13
% to reduce to the case in which V is an eigenspace for \alpha ; then use Exercise
% 7.12.)
% 7.15. A complete flag of subspaces of a vector space V of dimension n is a
% sequence
% of nested subspaces
% 0 = V0 ? V1 ? \dots  ? Vn-1 ? Vn = V
% with dim Vi = i. In other words, a complete flag is a composition series in the
% sense
% of Exercise 1.16.
% Let V be a finite-dimensional vector space over an algebraically closed field.
% Prove that every linear transformation \alpha  of V preserves a complete flag:
% that is,
% there is a complete flag as above and such that \alpha (Vi) \subseteq  Vi .
% Find a linear transformation of R2 that does not preserve a complete flag.
% 7.16. (Schur form) Let A \in  Mn (C). Prove that there exists a unitary matrix
% P \in  U(n) such that
% (
%  )
% \lambda 1 \ast 
%  \ast 
%  ...
%  \ast 
% | |
%  0 \lambda 2 \ast 
%  ...
%  \ast  |
%  |
% A = P | | ..
%  .
%  .
%  ..
%  .
%  ..
%  .
%  ..
%  . .. |
%  | P -1 .
% |
%  |
% ( 0
%  0 . . . \lambda  n-1 \ast  )
%  0 ...
%  \lambda n
% So not only is A similar to an upper triangular matrix (this is already
% guaranteed
% by the Jordan canonical form), but a matrix of change of basis P
% `triangularizing'
% the matrix A can in fact be chosen to be in U(n). (Argue inductively on n. Since
% C is algebraically closed, A has at least one eigenvector v1 , and we may assume
% \bot (v1, v1 ) = 1. For the induction step, consider the complement v1 ; cf. Exercise 6.17,
% and keep in mind Exercise 6.18.)
% The upper triangular matrix is a Schur form for A. It is (of course) not unique.
% 7.17. \lnot  A matrix M \in  Mn (C) is normal if M M \dagger  = M \dagger M .
% Note that unitary
% matrices and hermitian matrices are both normal.
% Prove that triangular normal matrices are diagonal. [7.18]
% 7.18. \lnot  Prove the spectral theorem for normal operators: if M is a normal
% matrix,
% then there exists an orthonormal basis of eigenvectors of M . Therefore, normal
% operators are diagonalizable; they determine a spectral decomposition of the
% vector
% space on which they act. (Consider the Schur form of M ; cf. Exercise 7.16. Use
% Exercise 7.17.) [7.19, 7.20]
% 7.19. Prove that a matrix M \in  Mn (C) is normal if and only if it admits an
% orthonormal basis of eigenvectors. (Exercise 7.18 gives one direction; prove the
% converse.)
% 7.20. ? Prove that real symmetric matrices are diagonalizable and admit an or-
% thonormal basis of eigenvectors. (This is an immediate consequence of Exercise
% 7.18
% and Exercise 6.22.) [\S{}7.4]

%%% Local Variables:
%%% mode: LaTeX
%%% TeX-engine: xetex
%%% TeX-master: "../master"
%%% End:
